% Generated from public Doxygen XML by tools/generate_api.py.
\begingroup\setlength{\parskip}{2pt}
\Needspace{9\baselineskip}
\section{nurbspath::application\_info}
Optional output identification of the application producing a stream.
\par\noindent\begin{minipage}{\linewidth}
\begin{lstlisting}[style=apireference]
std::string name;
\end{lstlisting}
{\footnotesize Nonempty application-name token.\par}
\end{minipage}\par\vspace{1pt}
\par\noindent\begin{minipage}{\linewidth}
\begin{lstlisting}[style=apireference]
std::string version;
\end{lstlisting}
{\footnotesize Nonempty application-version token.\par}
\end{minipage}\par\vspace{1pt}
\par\noindent\begin{minipage}{\linewidth}
\begin{lstlisting}[style=apireference]
std::ostream & write(std::ostream &output) const;
\end{lstlisting}
{\footnotesize Write one optional application v1 <name> <version> output row.\par}
{\footnotesize\textit{Returns:} The destination stream, with its formatting preserved.\par}
\end{minipage}\par\vspace{1pt}
\Needspace{9\baselineskip}
\section{nurbspath::circle2}
Circle in the independent two-dimensional Cartesian world.
\par\noindent\begin{minipage}{\linewidth}
\begin{lstlisting}[style=apireference]
circle2(const point2< REAL > &center_value, REAL radius_value);
\end{lstlisting}
{\footnotesize Construct a circle.\par}
\end{minipage}\par\vspace{1pt}
\par\noindent\begin{minipage}{\linewidth}
\begin{lstlisting}[style=apireference]
const point2< REAL > & center() const noexcept;
\end{lstlisting}
{\footnotesize Get the center.\par}
{\footnotesize\textit{Returns:} Constant center reference.\par}
\end{minipage}\par\vspace{1pt}
\par\noindent\begin{minipage}{\linewidth}
\begin{lstlisting}[style=apireference]
REAL radius() const noexcept;
\end{lstlisting}
{\footnotesize Get the radius.\par}
{\footnotesize\textit{Returns:} Radius in 2D world units.\par}
\end{minipage}\par\vspace{1pt}
\par\noindent\begin{minipage}{\linewidth}
\begin{lstlisting}[style=apireference]
const point2< REAL > & get_center() const noexcept;
\end{lstlisting}
{\footnotesize Get the center.\par}
{\footnotesize\textit{Returns:} Constant reference to the center in the 2D world.\par}
\end{minipage}\par\vspace{1pt}
\par\noindent\begin{minipage}{\linewidth}
\begin{lstlisting}[style=apireference]
REAL get_radius() const noexcept;
\end{lstlisting}
{\footnotesize Get the radius.\par}
{\footnotesize\textit{Returns:} Radius in 2D world units.\par}
\end{minipage}\par\vspace{1pt}
\par\noindent\begin{minipage}{\linewidth}
\begin{lstlisting}[style=apireference]
void set_center(const point2< REAL > &center_value) noexcept;
\end{lstlisting}
{\footnotesize Replace the center while preserving the radius.\par}
\end{minipage}\par\vspace{1pt}
\par\noindent\begin{minipage}{\linewidth}
\begin{lstlisting}[style=apireference]
template <typename VECTOR>
void set_center(const VECTOR &center_value) noexcept;
\end{lstlisting}
{\footnotesize Replace the center using vector components while preserving radius.\par}
\end{minipage}\par\vspace{1pt}
\par\noindent\begin{minipage}{\linewidth}
\begin{lstlisting}[style=apireference]
void set_radius(REAL radius_value);
\end{lstlisting}
{\footnotesize Replace the radius while preserving the center.\par}
\end{minipage}\par\vspace{1pt}
\par\noindent\begin{minipage}{\linewidth}
\begin{lstlisting}[style=apireference]
void set_center_and_radius(const point2< REAL > &center_value, REAL radius_value);
\end{lstlisting}
{\footnotesize Replace the center and radius together.\par}
\end{minipage}\par\vspace{1pt}
\par\noindent\begin{minipage}{\linewidth}
\begin{lstlisting}[style=apireference]
template <typename VECTOR>
void set_center_and_radius(const VECTOR &center_value, REAL radius_value);
\end{lstlisting}
{\footnotesize Replace the center from vector components and replace the radius.\par}
\end{minipage}\par\vspace{1pt}
\par\noindent\begin{minipage}{\linewidth}
\begin{lstlisting}[style=apireference]
circle2 & operator=(const vector2< REAL > &center_value) noexcept;
\end{lstlisting}
{\footnotesize Replace the center using vector components while preserving the radius.\par}
{\footnotesize\textit{Returns:} Reference to this circle.\par}
\end{minipage}\par\vspace{1pt}
\par\noindent\begin{minipage}{\linewidth}
\begin{lstlisting}[style=apireference]
point2< REAL > point_at(REAL u) const noexcept;
\end{lstlisting}
{\footnotesize Evaluate the circle parameterization.\par}
{\footnotesize\textit{Returns:} Point on the circle in the 2D world.\par}
\end{minipage}\par\vspace{1pt}
\par\noindent\begin{minipage}{\linewidth}
\begin{lstlisting}[style=apireference]
vector2< REAL > normal_at(const point2< REAL > &point) const;
\end{lstlisting}
{\footnotesize Compute the outward radial unit normal at a point.\par}
{\footnotesize\textit{Returns:} Outward unit direction after radial projection.\par}
\end{minipage}\par\vspace{1pt}
\par\noindent\begin{minipage}{\linewidth}
\begin{lstlisting}[style=apireference]
REAL parameter_of(const point2< REAL > &point, REAL tolerance=REAL(0)) const;
\end{lstlisting}
{\footnotesize Recover the angular parameter by radial projection.\par}
{\footnotesize\textit{Returns:} Counterclockwise angle in [0,2*pi).\par}
\end{minipage}\par\vspace{1pt}
\Needspace{9\baselineskip}
\section{nurbspath::diagnostic\_message3}
Optional output diagnostic anchored at a position in the 3D world.
\par\noindent\begin{minipage}{\linewidth}
\begin{lstlisting}[style=apireference]
diagnostic_level level = diagnostic_level::information;
\end{lstlisting}
{\footnotesize Output severity.\par}
\end{minipage}\par\vspace{1pt}
\par\noindent\begin{minipage}{\linewidth}
\begin{lstlisting}[style=apireference]
point3< REAL > position {};
\end{lstlisting}
{\footnotesize Finite world-space position of the annotation.\par}
\end{minipage}\par\vspace{1pt}
\par\noindent\begin{minipage}{\linewidth}
\begin{lstlisting}[style=apireference]
REAL size = REAL(12);
\end{lstlisting}
{\footnotesize Positive text size in consumer display units.\par}
\end{minipage}\par\vspace{1pt}
\par\noindent\begin{minipage}{\linewidth}
\begin{lstlisting}[style=apireference]
std::string value;
\end{lstlisting}
{\footnotesize Text, optionally containing escaped line breaks.\par}
\end{minipage}\par\vspace{1pt}
\par\noindent\begin{minipage}{\linewidth}
\begin{lstlisting}[style=apireference]
std::ostream & tag_write(std::size_t tag, std::ostream &output) const;
\end{lstlisting}
{\footnotesize Write an optional positioned warning, information, or error row.\par}
{\footnotesize\textit{Returns:} The destination stream, with its formatting preserved.\par}
\end{minipage}\par\vspace{1pt}
\Needspace{9\baselineskip}
\section{nurbspath::intersection2}
2D contact with parameters of both participating entities.
\par\noindent\begin{minipage}{\linewidth}
\begin{lstlisting}[style=apireference]
point2< REAL > point;
\end{lstlisting}
{\footnotesize Contact point in the independent 2D world.\par}
\end{minipage}\par\vspace{1pt}
\par\noindent\begin{minipage}{\linewidth}
\begin{lstlisting}[style=apireference]
REAL s = REAL(0);
\end{lstlisting}
{\footnotesize Ray or spline parameter.\par}
\end{minipage}\par\vspace{1pt}
\par\noindent\begin{minipage}{\linewidth}
\begin{lstlisting}[style=apireference]
REAL u = REAL(0);
\end{lstlisting}
{\footnotesize Circle angle in [0,2*pi).\par}
\end{minipage}\par\vspace{1pt}
\par\noindent\begin{minipage}{\linewidth}
\begin{lstlisting}[style=apireference]
REAL residual = REAL(0);
\end{lstlisting}
{\footnotesize Absolute 2D circle residual after refinement.\par}
\end{minipage}\par\vspace{1pt}
\Needspace{9\baselineskip}
\section{nurbspath::intersection3}
World-space contact with parameters of both participating entities.
\par\noindent\begin{minipage}{\linewidth}
\begin{lstlisting}[style=apireference]
point3< REAL > point;
\end{lstlisting}
{\footnotesize Contact point in world coordinates.\par}
\end{minipage}\par\vspace{1pt}
\par\noindent\begin{minipage}{\linewidth}
\begin{lstlisting}[style=apireference]
REAL s = REAL(0);
\end{lstlisting}
{\footnotesize Ray or spline parameter.\par}
\end{minipage}\par\vspace{1pt}
\par\noindent\begin{minipage}{\linewidth}
\begin{lstlisting}[style=apireference]
REAL u = REAL(0);
\end{lstlisting}
{\footnotesize Sphere or plane u parameter.\par}
\end{minipage}\par\vspace{1pt}
\par\noindent\begin{minipage}{\linewidth}
\begin{lstlisting}[style=apireference]
REAL v = REAL(0);
\end{lstlisting}
{\footnotesize Sphere or plane v parameter.\par}
\end{minipage}\par\vspace{1pt}
\par\noindent\begin{minipage}{\linewidth}
\begin{lstlisting}[style=apireference]
REAL residual = REAL(0);
\end{lstlisting}
{\footnotesize Absolute surface residual after refinement.\par}
\end{minipage}\par\vspace{1pt}
\Needspace{9\baselineskip}
\section{nurbspath::intersection\_result2}
Complete result of one two-dimensional intersection query.
\par\noindent\begin{minipage}{\linewidth}
\begin{lstlisting}[style=apireference]
intersection_kind kind = intersection_kind::none;
\end{lstlisting}
{\footnotesize Set classification.\par}
\end{minipage}\par\vspace{1pt}
\par\noindent\begin{minipage}{\linewidth}
\begin{lstlisting}[style=apireference]
std::vector< intersection2< REAL > > points;
\end{lstlisting}
{\footnotesize Isolated contacts sorted by s.\par}
\end{minipage}\par\vspace{1pt}
\par\noindent\begin{minipage}{\linewidth}
\begin{lstlisting}[style=apireference]
bool has_intersection() const noexcept;
\end{lstlisting}
{\footnotesize Test for isolated or coincident contact.\par}
{\footnotesize\textit{Returns:} True unless kind is none.\par}
\end{minipage}\par\vspace{1pt}
\par\noindent\begin{minipage}{\linewidth}
\begin{lstlisting}[style=apireference]
bool is_coincident() const noexcept;
\end{lstlisting}
{\footnotesize Test for continuous overlap.\par}
{\footnotesize\textit{Returns:} True when kind is coincident.\par}
\end{minipage}\par\vspace{1pt}
\Needspace{9\baselineskip}
\section{nurbspath::intersection\_result3}
Complete result of one intersection query.
\par\noindent\begin{minipage}{\linewidth}
\begin{lstlisting}[style=apireference]
intersection_kind kind = intersection_kind::none;
\end{lstlisting}
{\footnotesize Set classification.\par}
\end{minipage}\par\vspace{1pt}
\par\noindent\begin{minipage}{\linewidth}
\begin{lstlisting}[style=apireference]
std::vector< intersection3< REAL > > points;
\end{lstlisting}
{\footnotesize Isolated contacts sorted by s.\par}
\end{minipage}\par\vspace{1pt}
\par\noindent\begin{minipage}{\linewidth}
\begin{lstlisting}[style=apireference]
bool has_intersection() const noexcept;
\end{lstlisting}
{\footnotesize Test for any isolated or coincident contact.\par}
{\footnotesize\textit{Returns:} True unless kind is none.\par}
\end{minipage}\par\vspace{1pt}
\par\noindent\begin{minipage}{\linewidth}
\begin{lstlisting}[style=apireference]
bool is_coincident() const noexcept;
\end{lstlisting}
{\footnotesize Test for continuous overlap.\par}
{\footnotesize\textit{Returns:} True when kind is coincident.\par}
\end{minipage}\par\vspace{1pt}
\Needspace{9\baselineskip}
\section{nurbspath::line\_color}
Optional normalized RGBA line color associated with an application tag.
\par\noindent\begin{minipage}{\linewidth}
\begin{lstlisting}[style=apireference]
REAL red = REAL(0);
\end{lstlisting}
{\footnotesize Normalized red component in [0, 1].\par}
\end{minipage}\par\vspace{1pt}
\par\noindent\begin{minipage}{\linewidth}
\begin{lstlisting}[style=apireference]
REAL green = REAL(0);
\end{lstlisting}
{\footnotesize Normalized green component in [0, 1].\par}
\end{minipage}\par\vspace{1pt}
\par\noindent\begin{minipage}{\linewidth}
\begin{lstlisting}[style=apireference]
REAL blue = REAL(0);
\end{lstlisting}
{\footnotesize Normalized blue component in [0, 1].\par}
\end{minipage}\par\vspace{1pt}
\par\noindent\begin{minipage}{\linewidth}
\begin{lstlisting}[style=apireference]
REAL alpha = REAL(1);
\end{lstlisting}
{\footnotesize Normalized opacity in [0, 1].\par}
\end{minipage}\par\vspace{1pt}
\par\noindent\begin{minipage}{\linewidth}
\begin{lstlisting}[style=apireference]
std::ostream & tag_write(std::size_t tag, std::ostream &output) const;
\end{lstlisting}
{\footnotesize Write one version-1 linecolor row without changing stream format.\par}
{\footnotesize\textit{Returns:} The destination stream.\par}
\end{minipage}\par\vspace{1pt}
\par\noindent\begin{minipage}{\linewidth}
\begin{lstlisting}[style=apireference]
static std::optional< tagged_read_result< line_color > > tag_read(std::istream &input);
\end{lstlisting}
{\footnotesize Read one linecolor v1 row and allocate a validated color. Advisory application and diagnostic rows are skipped without decoding.\par}
{\footnotesize\textit{Returns:} Preserved tag and non-null color owner, or nullopt on EOF/failure. Malformed rows set failbit and consume only their own physical row.\par}
\end{minipage}\par\vspace{1pt}
\Needspace{9\baselineskip}
\section{nurbspath::line\_font}
Optional line width and repeating dash lengths for tagged line elements.
\par\noindent\begin{minipage}{\linewidth}
\begin{lstlisting}[style=apireference]
REAL width = REAL(1);
\end{lstlisting}
{\footnotesize Positive stroke width in consumer display units.\par}
\end{minipage}\par\vspace{1pt}
\par\noindent\begin{minipage}{\linewidth}
\begin{lstlisting}[style=apireference]
REAL on = REAL(0);
\end{lstlisting}
{\footnotesize Visible dash length; zero with off for a solid line.\par}
\end{minipage}\par\vspace{1pt}
\par\noindent\begin{minipage}{\linewidth}
\begin{lstlisting}[style=apireference]
REAL off = REAL(0);
\end{lstlisting}
{\footnotesize Gap length; zero with on for a solid line.\par}
\end{minipage}\par\vspace{1pt}
\par\noindent\begin{minipage}{\linewidth}
\begin{lstlisting}[style=apireference]
std::ostream & tag_write(std::size_t tag, std::ostream &output) const;
\end{lstlisting}
{\footnotesize Write one version-1 linefont row without changing stream format.\par}
{\footnotesize\textit{Returns:} The destination stream.\par}
\end{minipage}\par\vspace{1pt}
\par\noindent\begin{minipage}{\linewidth}
\begin{lstlisting}[style=apireference]
static std::optional< tagged_read_result< line_font > > tag_read(std::istream &input);
\end{lstlisting}
{\footnotesize Read one linefont v1 row and allocate a validated line font. Advisory application and diagnostic rows are skipped without decoding.\par}
{\footnotesize\textit{Returns:} Preserved tag and non-null font owner, or nullopt on EOF/failure. Malformed rows set failbit and consume only their own physical row.\par}
\end{minipage}\par\vspace{1pt}
\Needspace{9\baselineskip}
\section{nurbspath::numerical\_settings}
Controls iterative and sampled numerical algorithms.
\par\noindent\begin{minipage}{\linewidth}
\begin{lstlisting}[style=apireference]
REAL residual_tolerance = REAL(1e-8);
\end{lstlisting}
{\footnotesize Accepted absolute function residual.\par}
\end{minipage}\par\vspace{1pt}
\par\noindent\begin{minipage}{\linewidth}
\begin{lstlisting}[style=apireference]
REAL parameter_tolerance = REAL(1e-10);
\end{lstlisting}
{\footnotesize Accepted parameter interval width.\par}
\end{minipage}\par\vspace{1pt}
\par\noindent\begin{minipage}{\linewidth}
\begin{lstlisting}[style=apireference]
std::size_t max_iterations = 100;
\end{lstlisting}
{\footnotesize Maximum refinement iterations.\par}
\end{minipage}\par\vspace{1pt}
\par\noindent\begin{minipage}{\linewidth}
\begin{lstlisting}[style=apireference]
std::size_t sample_count = 256;
\end{lstlisting}
{\footnotesize Number of initial uniform intervals.\par}
\end{minipage}\par\vspace{1pt}
\par\noindent\begin{minipage}{\linewidth}
\begin{lstlisting}[style=apireference]
void validate() const;
\end{lstlisting}
{\footnotesize Validate all numerical controls.\par}
\end{minipage}\par\vspace{1pt}
\Needspace{9\baselineskip}
\section{nurbspath::nurbs\_arr\_spline2}
Self-contained 2D NURBS curve backed by one public scalar valarray.
\par\noindent\begin{minipage}{\linewidth}
\begin{lstlisting}[style=apireference]
using real_t = REAL;
\end{lstlisting}
{\footnotesize Floating-point scalar stored by this spline.\par}
\end{minipage}\par\vspace{1pt}
\par\noindent\begin{minipage}{\linewidth}
\begin{lstlisting}[style=apireference]
using parameter_index_range = std::ranges::iota_view< std::size_t, std::size_t >;
\end{lstlisting}
{\footnotesize Allocation-free range of flattened indices into parameters.\par}
\end{minipage}\par\vspace{1pt}
\par\noindent\begin{minipage}{\linewidth}
\begin{lstlisting}[style=apireference]
std::valarray< REAL > parameters;
\end{lstlisting}
{\footnotesize Flattened control coordinates, weights, and knots.\par}
\end{minipage}\par\vspace{1pt}
\par\noindent\begin{minipage}{\linewidth}
\begin{lstlisting}[style=apireference]
nurbs_arr_spline2(const point2< REAL > &start, const point2< REAL > &end, std::size_t control_point_count, std::size_t degree, REAL tolerance);
\end{lstlisting}
{\footnotesize Construct a standard open 2D spline with explicit tolerance.\par}
\end{minipage}\par\vspace{1pt}
\par\noindent\begin{minipage}{\linewidth}
\begin{lstlisting}[style=apireference]
template <std::same_as< bool > CLOSED = bool>
nurbs_arr_spline2(const point2< REAL > &start, const point2< REAL > &end, std::size_t control_point_count, std::size_t degree, CLOSED closed=false, REAL tolerance=REAL(64) *std::numeric_limits< REAL >::epsilon());
\end{lstlisting}
{\footnotesize Construct a standard open or closed spline between two points.\par}
\end{minipage}\par\vspace{1pt}
\par\noindent\begin{minipage}{\linewidth}
\begin{lstlisting}[style=apireference]
nurbs_arr_spline2(std::valarray< REAL > values, std::size_t degree, REAL tolerance=REAL(64) *std::numeric_limits< REAL >::epsilon());
\end{lstlisting}
{\footnotesize Construct an open array-backed spline from flattened parameters.\par}
\end{minipage}\par\vspace{1pt}
\par\noindent\begin{minipage}{\linewidth}
\begin{lstlisting}[style=apireference]
template <std::same_as< bool > CLOSED>
nurbs_arr_spline2(std::valarray< REAL > values, std::size_t degree, CLOSED closed, REAL tolerance=REAL(64) *std::numeric_limits< REAL >::epsilon());
\end{lstlisting}
{\footnotesize Construct an open or closed spline from flattened parameters.\par}
\end{minipage}\par\vspace{1pt}
\par\noindent\begin{minipage}{\linewidth}
\begin{lstlisting}[style=apireference]
nurbs_arr_spline2(std::vector< point2< REAL > > control_points, std::vector< REAL > weights, std::vector< REAL > knots, std::size_t degree, REAL tolerance=REAL(64) *std::numeric_limits< REAL >::epsilon());
\end{lstlisting}
{\footnotesize Construct an open array-backed spline from component vectors.\par}
\end{minipage}\par\vspace{1pt}
\par\noindent\begin{minipage}{\linewidth}
\begin{lstlisting}[style=apireference]
template <std::same_as< bool > CLOSED>
nurbs_arr_spline2(std::vector< point2< REAL > > control_points, std::vector< REAL > weights, std::vector< REAL > knots, std::size_t degree, CLOSED closed, REAL tolerance=REAL(64) *std::numeric_limits< REAL >::epsilon());
\end{lstlisting}
{\footnotesize Construct an open or closed spline from component vectors.\par}
\end{minipage}\par\vspace{1pt}
\par\noindent\begin{minipage}{\linewidth}
\begin{lstlisting}[style=apireference]
nurbs_arr_spline2(const nurbs_spline2< REAL > &spline);
\end{lstlisting}
{\footnotesize Copy an ordinary 2D NURBS spline into flattened storage.\par}
\end{minipage}\par\vspace{1pt}
\par\noindent\begin{minipage}{\linewidth}
\begin{lstlisting}[style=apireference]
std::size_t degree() const noexcept;
\end{lstlisting}
{\footnotesize Get the polynomial degree.\par}
{\footnotesize\textit{Returns:} Positive spline degree.\par}
\end{minipage}\par\vspace{1pt}
\par\noindent\begin{minipage}{\linewidth}
\begin{lstlisting}[style=apireference]
REAL tolerance() const noexcept;
\end{lstlisting}
{\footnotesize Get the definition and parameter-boundary tolerance.\par}
{\footnotesize\textit{Returns:} Current tolerance.\par}
\end{minipage}\par\vspace{1pt}
\par\noindent\begin{minipage}{\linewidth}
\begin{lstlisting}[style=apireference]
bool is_closed() const noexcept;
\end{lstlisting}
{\footnotesize Report whether the spline requires a closed seam.\par}
{\footnotesize\textit{Returns:} True when active-domain endpoints must coincide.\par}
\end{minipage}\par\vspace{1pt}
\par\noindent\begin{minipage}{\linewidth}
\begin{lstlisting}[style=apireference]
std::size_t number_of_parameters() const noexcept;
\end{lstlisting}
{\footnotesize Get the number of flattened scalar parameters.\par}
{\footnotesize\textit{Returns:} Current public array size, including malformed layouts.\par}
\end{minipage}\par\vspace{1pt}
\par\noindent\begin{minipage}{\linewidth}
\begin{lstlisting}[style=apireference]
std::size_t control_point_count() const;
\end{lstlisting}
{\footnotesize Derive the control-point count from the current array layout.\par}
{\footnotesize\textit{Returns:} Number of represented control points.\par}
\end{minipage}\par\vspace{1pt}
\par\noindent\begin{minipage}{\linewidth}
\begin{lstlisting}[style=apireference]
std::size_t knot_count() const;
\end{lstlisting}
{\footnotesize Derive the knot count from the current array layout.\par}
{\footnotesize\textit{Returns:} Number of represented knots.\par}
\end{minipage}\par\vspace{1pt}
\par\noindent\begin{minipage}{\linewidth}
\begin{lstlisting}[style=apireference]
std::size_t weight_count() const;
\end{lstlisting}
{\footnotesize Derive the weight count from the current array layout.\par}
{\footnotesize\textit{Returns:} Number of represented weights, equal to the control-point count.\par}
\end{minipage}\par\vspace{1pt}
\par\noindent\begin{minipage}{\linewidth}
\begin{lstlisting}[style=apireference]
parameter_index_range get_control_point_count_range() const;
\end{lstlisting}
{\footnotesize Get the flattened indices occupied by control-point coordinates.\par}
{\footnotesize\textit{Returns:} Iterable scalar-index range [0, 2 * control\_point\_count()).\par}
\end{minipage}\par\vspace{1pt}
\par\noindent\begin{minipage}{\linewidth}
\begin{lstlisting}[style=apireference]
parameter_index_range get_weight_count_range() const;
\end{lstlisting}
{\footnotesize Get the flattened indices occupied by rational weights.\par}
{\footnotesize\textit{Returns:} Iterable scalar-index range [2 * C, 3 * C), where C is the control-point count.\par}
\end{minipage}\par\vspace{1pt}
\par\noindent\begin{minipage}{\linewidth}
\begin{lstlisting}[style=apireference]
parameter_index_range get_knot_count_range() const;
\end{lstlisting}
{\footnotesize Get the flattened indices occupied by knots.\par}
{\footnotesize\textit{Returns:} Iterable scalar-index range [3 * C, parameters.size()), where C is the control-point count.\par}
\end{minipage}\par\vspace{1pt}
\par\noindent\begin{minipage}{\linewidth}
\begin{lstlisting}[style=apireference]
std::vector< point2< REAL > > get_control_points() const;
\end{lstlisting}
{\footnotesize Get all control points as detached point values.\par}
{\footnotesize\textit{Returns:} Control points reconstructed in index order.\par}
\end{minipage}\par\vspace{1pt}
\par\noindent\begin{minipage}{\linewidth}
\begin{lstlisting}[style=apireference]
std::vector< REAL > get_weights() const;
\end{lstlisting}
{\footnotesize Get all weights as a detached vector.\par}
{\footnotesize\textit{Returns:} Weights copied in index order.\par}
\end{minipage}\par\vspace{1pt}
\par\noindent\begin{minipage}{\linewidth}
\begin{lstlisting}[style=apireference]
std::vector< REAL > get_knots() const;
\end{lstlisting}
{\footnotesize Get all knots as a detached vector.\par}
{\footnotesize\textit{Returns:} Knots copied in index order.\par}
\end{minipage}\par\vspace{1pt}
\par\noindent\begin{minipage}{\linewidth}
\begin{lstlisting}[style=apireference]
point2< REAL > get_control_point(std::size_t index) const;
\end{lstlisting}
{\footnotesize Get one control point by index.\par}
{\footnotesize\textit{Returns:} Point reconstructed by value from two adjacent scalars.\par}
\end{minipage}\par\vspace{1pt}
\par\noindent\begin{minipage}{\linewidth}
\begin{lstlisting}[style=apireference]
REAL get_weight(std::size_t index) const;
\end{lstlisting}
{\footnotesize Get one rational weight by index.\par}
{\footnotesize\textit{Returns:} Selected weight value.\par}
\end{minipage}\par\vspace{1pt}
\par\noindent\begin{minipage}{\linewidth}
\begin{lstlisting}[style=apireference]
REAL get_knot(std::size_t index) const;
\end{lstlisting}
{\footnotesize Get one knot by index.\par}
{\footnotesize\textit{Returns:} Selected knot value.\par}
\end{minipage}\par\vspace{1pt}
\par\noindent\begin{minipage}{\linewidth}
\begin{lstlisting}[style=apireference]
REAL get_parameter(std::size_t index) const;
\end{lstlisting}
{\footnotesize Get one scalar by its flattened index.\par}
{\footnotesize\textit{Returns:} Selected scalar value.\par}
\end{minipage}\par\vspace{1pt}
\par\noindent\begin{minipage}{\linewidth}
\begin{lstlisting}[style=apireference]
std::string parameter_name(std::size_t index) const;
\end{lstlisting}
{\footnotesize Get the descriptive name of one flattened scalar.\par}
{\footnotesize\textit{Returns:} A P[i].x, P[i].y, W[i], or K[i] name.\par}
\end{minipage}\par\vspace{1pt}
\par\noindent\begin{minipage}{\linewidth}
\begin{lstlisting}[style=apireference]
void validate() const;
\end{lstlisting}
{\footnotesize Validate the current public array and all typed metadata.\par}
\end{minipage}\par\vspace{1pt}
\par\noindent\begin{minipage}{\linewidth}
\begin{lstlisting}[style=apireference]
void set_parameters(std::valarray< REAL > values);
\end{lstlisting}
{\footnotesize Atomically validate and replace the complete public scalar array.\par}
\end{minipage}\par\vspace{1pt}
\par\noindent\begin{minipage}{\linewidth}
\begin{lstlisting}[style=apireference]
void set_parameters(const REAL *values);
\end{lstlisting}
{\footnotesize Atomically copy and replace the complete public scalar array.\par}
\end{minipage}\par\vspace{1pt}
\par\noindent\begin{minipage}{\linewidth}
\begin{lstlisting}[style=apireference]
bool set_parameter(std::size_t index, REAL value);
\end{lstlisting}
{\footnotesize Atomically validate and replace one flattened scalar.\par}
{\footnotesize\textit{Returns:} True when the candidate was valid and committed; false for an out-of-range index, malformed current shape, or a rejected value.\par}
\end{minipage}\par\vspace{1pt}
\par\noindent\begin{minipage}{\linewidth}
\begin{lstlisting}[style=apireference]
void set_control_point(std::size_t index, const point2< REAL > &point);
\end{lstlisting}
{\footnotesize Atomically replace one control point.\par}
\end{minipage}\par\vspace{1pt}
\par\noindent\begin{minipage}{\linewidth}
\begin{lstlisting}[style=apireference]
void set_weight(std::size_t index, REAL value);
\end{lstlisting}
{\footnotesize Atomically replace one rational weight.\par}
\end{minipage}\par\vspace{1pt}
\par\noindent\begin{minipage}{\linewidth}
\begin{lstlisting}[style=apireference]
void set_knot(std::size_t index, REAL value);
\end{lstlisting}
{\footnotesize Atomically replace one knot.\par}
\end{minipage}\par\vspace{1pt}
\par\noindent\begin{minipage}{\linewidth}
\begin{lstlisting}[style=apireference]
void set_control_points(std::vector< point2< REAL > > control_points);
\end{lstlisting}
{\footnotesize Atomically replace all control points while retaining other data.\par}
\end{minipage}\par\vspace{1pt}
\par\noindent\begin{minipage}{\linewidth}
\begin{lstlisting}[style=apireference]
void set_weights(std::vector< REAL > values);
\end{lstlisting}
{\footnotesize Atomically replace all weights while retaining other data.\par}
\end{minipage}\par\vspace{1pt}
\par\noindent\begin{minipage}{\linewidth}
\begin{lstlisting}[style=apireference]
void set_knots(std::vector< REAL > values);
\end{lstlisting}
{\footnotesize Atomically replace all knots while retaining other data.\par}
\end{minipage}\par\vspace{1pt}
\par\noindent\begin{minipage}{\linewidth}
\begin{lstlisting}[style=apireference]
void set_standard_knots(REAL s_end);
\end{lstlisting}
{\footnotesize Set standard open-clamped knots on [0, s\_end].\par}
\end{minipage}\par\vspace{1pt}
\par\noindent\begin{minipage}{\linewidth}
\begin{lstlisting}[style=apireference]
void set_standard_knots(REAL s_start, REAL s_end);
\end{lstlisting}
{\footnotesize Set standard open-clamped knots on [s\_start, s\_end].\par}
\end{minipage}\par\vspace{1pt}
\par\noindent\begin{minipage}{\linewidth}
\begin{lstlisting}[style=apireference]
void set_tolerance(REAL value);
\end{lstlisting}
{\footnotesize Atomically replace the validation tolerance.\par}
\end{minipage}\par\vspace{1pt}
\par\noindent\begin{minipage}{\linewidth}
\begin{lstlisting}[style=apireference]
void set_closed(bool value);
\end{lstlisting}
{\footnotesize Atomically change the closed-seam requirement.\par}
\end{minipage}\par\vspace{1pt}
\par\noindent\begin{minipage}{\linewidth}
\begin{lstlisting}[style=apireference]
void set_definition(std::vector< point2< REAL > > control_points, std::vector< REAL > weights, std::vector< REAL > knots, bool closed, REAL tolerance);
\end{lstlisting}
{\footnotesize Atomically replace all numeric components and mutable metadata.\par}
\end{minipage}\par\vspace{1pt}
\par\noindent\begin{minipage}{\linewidth}
\begin{lstlisting}[style=apireference]
REAL s_min() const;
\end{lstlisting}
{\footnotesize Get the lower active parameter bound from the current knot data.\par}
{\footnotesize\textit{Returns:} Minimum valid native s.\par}
\end{minipage}\par\vspace{1pt}
\par\noindent\begin{minipage}{\linewidth}
\begin{lstlisting}[style=apireference]
REAL s_max() const;
\end{lstlisting}
{\footnotesize Get the upper active parameter bound from the current knot data.\par}
{\footnotesize\textit{Returns:} Maximum valid native s.\par}
\end{minipage}\par\vspace{1pt}
\par\noindent\begin{minipage}{\linewidth}
\begin{lstlisting}[style=apireference]
point2< REAL > get_start() const;
\end{lstlisting}
{\footnotesize Evaluate the current point at the active-domain start.\par}
{\footnotesize\textit{Returns:} Start point derived on demand from public storage.\par}
\end{minipage}\par\vspace{1pt}
\par\noindent\begin{minipage}{\linewidth}
\begin{lstlisting}[style=apireference]
point2< REAL > get_end() const;
\end{lstlisting}
{\footnotesize Evaluate the current point at the active-domain end.\par}
{\footnotesize\textit{Returns:} End point derived on demand from public storage.\par}
\end{minipage}\par\vspace{1pt}
\par\noindent\begin{minipage}{\linewidth}
\begin{lstlisting}[style=apireference]
spline_derivatives2< REAL > derivatives_at(REAL s) const;
\end{lstlisting}
{\footnotesize Evaluate position and the first two analytic derivatives.\par}
{\footnotesize\textit{Returns:} Position, first derivative, and second derivative in 2D.\par}
\end{minipage}\par\vspace{1pt}
\par\noindent\begin{minipage}{\linewidth}
\begin{lstlisting}[style=apireference]
point2< REAL > evaluate(REAL s) const;
\end{lstlisting}
{\footnotesize Evaluate the curve position.\par}
{\footnotesize\textit{Returns:} Curve point in the independent 2D world.\par}
\end{minipage}\par\vspace{1pt}
\par\noindent\begin{minipage}{\linewidth}
\begin{lstlisting}[style=apireference]
point2< REAL > point_at(REAL s) const;
\end{lstlisting}
{\footnotesize Evaluate the curve position.\par}
{\footnotesize\textit{Returns:} Curve point in the independent 2D world.\par}
\end{minipage}\par\vspace{1pt}
\par\noindent\begin{minipage}{\linewidth}
\begin{lstlisting}[style=apireference]
vector2< REAL > first_derivative(REAL s) const;
\end{lstlisting}
{\footnotesize Evaluate the first analytic derivative.\par}
{\footnotesize\textit{Returns:} First derivative with respect to native s.\par}
\end{minipage}\par\vspace{1pt}
\par\noindent\begin{minipage}{\linewidth}
\begin{lstlisting}[style=apireference]
vector2< REAL > second_derivative(REAL s) const;
\end{lstlisting}
{\footnotesize Evaluate the second analytic derivative.\par}
{\footnotesize\textit{Returns:} Second derivative with respect to native s.\par}
\end{minipage}\par\vspace{1pt}
\par\noindent\begin{minipage}{\linewidth}
\begin{lstlisting}[style=apireference]
vector2< REAL > third_derivative(REAL s) const;
\end{lstlisting}
{\footnotesize Evaluate the third analytic rational derivative.\par}
{\footnotesize\textit{Returns:} Third derivative with respect to native s.\par}
\end{minipage}\par\vspace{1pt}
\par\noindent\begin{minipage}{\linewidth}
\begin{lstlisting}[style=apireference]
vector2< REAL > tangent(REAL s, REAL tolerance=vector2< REAL >::default_tolerance()) const;
\end{lstlisting}
{\footnotesize Evaluate a unit tangent.\par}
{\footnotesize\textit{Returns:} Normalized first derivative.\par}
\end{minipage}\par\vspace{1pt}
\par\noindent\begin{minipage}{\linewidth}
\begin{lstlisting}[style=apireference]
REAL curvature(REAL s, REAL tangent_tolerance=vector2< REAL >::default_tolerance()) const;
\end{lstlisting}
{\footnotesize Evaluate unsigned geometric curvature in the 2D world.\par}
{\footnotesize\textit{Returns:} Nonnegative curvature magnitude in inverse 2D world units.\par}
\end{minipage}\par\vspace{1pt}
\par\noindent\begin{minipage}{\linewidth}
\begin{lstlisting}[style=apireference]
REAL get_polygon_length() const;
\end{lstlisting}
{\footnotesize Get the total length of the current 2D control polygon.\par}
{\footnotesize\textit{Returns:} Control-polygon length in 2D world units.\par}
\end{minipage}\par\vspace{1pt}
\par\noindent\begin{minipage}{\linewidth}
\begin{lstlisting}[style=apireference]
REAL approximate_arc_length(std::size_t segment_count=512) const;
\end{lstlisting}
{\footnotesize Estimate arc length using a uniform-native-s polyline.\par}
{\footnotesize\textit{Returns:} Approximate length in 2D world units.\par}
\end{minipage}\par\vspace{1pt}
\par\noindent\begin{minipage}{\linewidth}
\begin{lstlisting}[style=apireference]
void adopt_to_points(const std::vector< point2< REAL > > &samples, const std::vector< REAL > &arc_length_parameters, std::size_t requested_degree=3, REAL tolerance=REAL(1e-10));
\end{lstlisting}
{\footnotesize Replace this curve with an open interpolant through samples.\par}
\end{minipage}\par\vspace{1pt}
\par\noindent\begin{minipage}{\linewidth}
\begin{lstlisting}[style=apireference]
void adopt_to_points(const std::vector< point2< REAL > > &samples, const std::vector< REAL > &arc_length_parameters, std::size_t requested_degree, bool closed, REAL tolerance=REAL(1e-10));
\end{lstlisting}
{\footnotesize Replace this curve with an open or closed interpolant.\par}
\end{minipage}\par\vspace{1pt}
\par\noindent\begin{minipage}{\linewidth}
\begin{lstlisting}[style=apireference]
nurbs_spline2< REAL > to_nurbs_spline() const;
\end{lstlisting}
{\footnotesize Convert to the ordinary vector-backed 2D NURBS type.\par}
{\footnotesize\textit{Returns:} Valid ordinary spline with identical geometry and metadata.\par}
\end{minipage}\par\vspace{1pt}
\par\noindent\begin{minipage}{\linewidth}
\begin{lstlisting}[style=apireference]
static nurbs_arr_spline2 interpolate(const std::vector< point2< REAL > > &samples, const std::vector< REAL > &arc_length_parameters, std::size_t requested_degree=3, REAL tolerance=REAL(1e-10));
\end{lstlisting}
{\footnotesize Build an open unit-weight spline interpolating measured samples.\par}
{\footnotesize\textit{Returns:} Array-backed interpolating curve.\par}
\end{minipage}\par\vspace{1pt}
\par\noindent\begin{minipage}{\linewidth}
\begin{lstlisting}[style=apireference]
static nurbs_arr_spline2 interpolate(const std::vector< point2< REAL > > &samples, const std::vector< REAL > &arc_length_parameters, std::size_t requested_degree, bool closed, REAL tolerance=REAL(1e-10));
\end{lstlisting}
{\footnotesize Build an open or closed unit-weight interpolating spline.\par}
{\footnotesize\textit{Returns:} Array-backed interpolating curve.\par}
\end{minipage}\par\vspace{1pt}
\Needspace{9\baselineskip}
\section{nurbspath::nurbs\_arr\_spline3}
Self-contained 3D NURBS curve backed by one public scalar valarray.
\par\noindent\begin{minipage}{\linewidth}
\begin{lstlisting}[style=apireference]
using real_t = REAL;
\end{lstlisting}
{\footnotesize Scalar type stored by this spline.\par}
\end{minipage}\par\vspace{1pt}
\par\noindent\begin{minipage}{\linewidth}
\begin{lstlisting}[style=apireference]
using parameter_index_range = std::ranges::iota_view< std::size_t, std::size_t >;
\end{lstlisting}
{\footnotesize Allocation-free range of flattened indices into parameters.\par}
\end{minipage}\par\vspace{1pt}
\par\noindent\begin{minipage}{\linewidth}
\begin{lstlisting}[style=apireference]
std::valarray< REAL > parameters;
\end{lstlisting}
{\footnotesize Flattened control-coordinate, weight, and knot storage.\par}
\end{minipage}\par\vspace{1pt}
\par\noindent\begin{minipage}{\linewidth}
\begin{lstlisting}[style=apireference]
nurbs_arr_spline3(const point3< REAL > &start, const point3< REAL > &end, std::size_t control_point_count, std::size_t degree, REAL tolerance);
\end{lstlisting}
{\footnotesize Construct a standard open 3D spline with explicit tolerance.\par}
\end{minipage}\par\vspace{1pt}
\par\noindent\begin{minipage}{\linewidth}
\begin{lstlisting}[style=apireference]
template <std::same_as< bool > CLOSED = bool>
nurbs_arr_spline3(const point3< REAL > &start, const point3< REAL > &end, std::size_t control_point_count, std::size_t degree, CLOSED closed=false, REAL tolerance=REAL(64) *std::numeric_limits< REAL >::epsilon());
\end{lstlisting}
{\footnotesize Construct a standard open or closed spline between two 3D points.\par}
\end{minipage}\par\vspace{1pt}
\par\noindent\begin{minipage}{\linewidth}
\begin{lstlisting}[style=apireference]
nurbs_arr_spline3(std::valarray< REAL > parameter_values, std::size_t degree, REAL tolerance=REAL(64) *std::numeric_limits< REAL >::epsilon());
\end{lstlisting}
{\footnotesize Construct an open array spline from flattened scalar parameters.\par}
\end{minipage}\par\vspace{1pt}
\par\noindent\begin{minipage}{\linewidth}
\begin{lstlisting}[style=apireference]
template <std::same_as< bool > CLOSED>
nurbs_arr_spline3(std::valarray< REAL > parameter_values, std::size_t degree, CLOSED closed, REAL tolerance=REAL(64) *std::numeric_limits< REAL >::epsilon());
\end{lstlisting}
{\footnotesize Construct an open or closed array spline from flattened scalars.\par}
\end{minipage}\par\vspace{1pt}
\par\noindent\begin{minipage}{\linewidth}
\begin{lstlisting}[style=apireference]
nurbs_arr_spline3(std::vector< point3< REAL > > control_points, std::vector< REAL > weights, std::vector< REAL > knots, std::size_t degree, REAL tolerance=REAL(64) *std::numeric_limits< REAL >::epsilon());
\end{lstlisting}
{\footnotesize Construct an open array spline from ordinary vector collections.\par}
\end{minipage}\par\vspace{1pt}
\par\noindent\begin{minipage}{\linewidth}
\begin{lstlisting}[style=apireference]
template <std::same_as< bool > CLOSED>
nurbs_arr_spline3(std::vector< point3< REAL > > control_points, std::vector< REAL > weights, std::vector< REAL > knots, std::size_t degree, CLOSED closed, REAL tolerance=REAL(64) *std::numeric_limits< REAL >::epsilon());
\end{lstlisting}
{\footnotesize Construct an open or closed array spline from vector collections.\par}
\end{minipage}\par\vspace{1pt}
\par\noindent\begin{minipage}{\linewidth}
\begin{lstlisting}[style=apireference]
nurbs_arr_spline3(const nurbs_spline3< REAL > &spline);
\end{lstlisting}
{\footnotesize Copy an ordinary 3D NURBS into independent array storage.\par}
\end{minipage}\par\vspace{1pt}
\par\noindent\begin{minipage}{\linewidth}
\begin{lstlisting}[style=apireference]
std::size_t number_of_parameters() const noexcept;
\end{lstlisting}
{\footnotesize Get the current flattened scalar count.\par}
{\footnotesize\textit{Returns:} Number of values in parameters, including an invalid draft.\par}
\end{minipage}\par\vspace{1pt}
\par\noindent\begin{minipage}{\linewidth}
\begin{lstlisting}[style=apireference]
std::size_t control_point_count() const;
\end{lstlisting}
{\footnotesize Derive the current control-point count from shape and degree.\par}
{\footnotesize\textit{Returns:} Number of control points and weights.\par}
\end{minipage}\par\vspace{1pt}
\par\noindent\begin{minipage}{\linewidth}
\begin{lstlisting}[style=apireference]
std::size_t weight_count() const;
\end{lstlisting}
{\footnotesize Derive the current weight count.\par}
{\footnotesize\textit{Returns:} Number of rational weights, equal to control\_point\_count().\par}
\end{minipage}\par\vspace{1pt}
\par\noindent\begin{minipage}{\linewidth}
\begin{lstlisting}[style=apireference]
std::size_t knot_count() const;
\end{lstlisting}
{\footnotesize Derive the current knot count.\par}
{\footnotesize\textit{Returns:} Complete knot-vector length.\par}
\end{minipage}\par\vspace{1pt}
\par\noindent\begin{minipage}{\linewidth}
\begin{lstlisting}[style=apireference]
parameter_index_range get_control_point_count_range() const;
\end{lstlisting}
{\footnotesize Get the flattened control-coordinate index range.\par}
{\footnotesize\textit{Returns:} Lazy ascending range of valid control-coordinate scalar indices.\par}
\end{minipage}\par\vspace{1pt}
\par\noindent\begin{minipage}{\linewidth}
\begin{lstlisting}[style=apireference]
parameter_index_range get_weight_count_range() const;
\end{lstlisting}
{\footnotesize Get the flattened weight index range.\par}
{\footnotesize\textit{Returns:} Lazy ascending range of valid flattened weight indices.\par}
\end{minipage}\par\vspace{1pt}
\par\noindent\begin{minipage}{\linewidth}
\begin{lstlisting}[style=apireference]
parameter_index_range get_knot_count_range() const;
\end{lstlisting}
{\footnotesize Get the flattened knot index range.\par}
{\footnotesize\textit{Returns:} Lazy ascending range of valid flattened knot indices.\par}
\end{minipage}\par\vspace{1pt}
\par\noindent\begin{minipage}{\linewidth}
\begin{lstlisting}[style=apireference]
std::size_t degree() const noexcept;
\end{lstlisting}
{\footnotesize Get the polynomial degree.\par}
{\footnotesize\textit{Returns:} Immutable spline degree.\par}
\end{minipage}\par\vspace{1pt}
\par\noindent\begin{minipage}{\linewidth}
\begin{lstlisting}[style=apireference]
REAL tolerance() const noexcept;
\end{lstlisting}
{\footnotesize Get the definition tolerance.\par}
{\footnotesize\textit{Returns:} Positive tolerance.\par}
\end{minipage}\par\vspace{1pt}
\par\noindent\begin{minipage}{\linewidth}
\begin{lstlisting}[style=apireference]
bool is_closed() const noexcept;
\end{lstlisting}
{\footnotesize Report closure state.\par}
{\footnotesize\textit{Returns:} True when closure is required.\par}
\end{minipage}\par\vspace{1pt}
\par\noindent\begin{minipage}{\linewidth}
\begin{lstlisting}[style=apireference]
std::vector< point3< REAL > > get_control_points() const;
\end{lstlisting}
{\footnotesize Copy all control points out of scalar storage.\par}
{\footnotesize\textit{Returns:} Point vector in control-point order.\par}
\end{minipage}\par\vspace{1pt}
\par\noindent\begin{minipage}{\linewidth}
\begin{lstlisting}[style=apireference]
std::vector< REAL > get_weights() const;
\end{lstlisting}
{\footnotesize Copy all rational weights out of scalar storage.\par}
{\footnotesize\textit{Returns:} Weight vector in control-point order.\par}
\end{minipage}\par\vspace{1pt}
\par\noindent\begin{minipage}{\linewidth}
\begin{lstlisting}[style=apireference]
std::vector< REAL > get_knots() const;
\end{lstlisting}
{\footnotesize Copy the complete knot vector out of scalar storage.\par}
{\footnotesize\textit{Returns:} Knot vector in stored order.\par}
\end{minipage}\par\vspace{1pt}
\par\noindent\begin{minipage}{\linewidth}
\begin{lstlisting}[style=apireference]
REAL get_parameter(std::size_t index) const;
\end{lstlisting}
{\footnotesize Get one flattened scalar parameter.\par}
{\footnotesize\textit{Returns:} Stored parameter value.\par}
\end{minipage}\par\vspace{1pt}
\par\noindent\begin{minipage}{\linewidth}
\begin{lstlisting}[style=apireference]
std::string parameter_name(std::size_t index) const;
\end{lstlisting}
{\footnotesize Get the descriptive name of one flattened scalar.\par}
{\footnotesize\textit{Returns:} P[i].x, P[i].y, P[i].z, W[i], or K[i].\par}
\end{minipage}\par\vspace{1pt}
\par\noindent\begin{minipage}{\linewidth}
\begin{lstlisting}[style=apireference]
bool set_parameter(std::size_t index, REAL value);
\end{lstlisting}
{\footnotesize Atomically validate and replace one scalar parameter.\par}
{\footnotesize\textit{Returns:} True when the candidate is valid and committed; false for an invalid shape, out-of-range index, or rejected candidate definition.\par}
\end{minipage}\par\vspace{1pt}
\par\noindent\begin{minipage}{\linewidth}
\begin{lstlisting}[style=apireference]
point3< REAL > get_control_point(std::size_t index) const;
\end{lstlisting}
{\footnotesize Materialize one control point from three stored scalars.\par}
{\footnotesize\textit{Returns:} Point value copied from X/Y/Z storage.\par}
\end{minipage}\par\vspace{1pt}
\par\noindent\begin{minipage}{\linewidth}
\begin{lstlisting}[style=apireference]
REAL get_weight(std::size_t index) const;
\end{lstlisting}
{\footnotesize Get one rational weight.\par}
{\footnotesize\textit{Returns:} Stored weight value.\par}
\end{minipage}\par\vspace{1pt}
\par\noindent\begin{minipage}{\linewidth}
\begin{lstlisting}[style=apireference]
REAL get_knot(std::size_t index) const;
\end{lstlisting}
{\footnotesize Get one knot.\par}
{\footnotesize\textit{Returns:} Stored knot value.\par}
\end{minipage}\par\vspace{1pt}
\par\noindent\begin{minipage}{\linewidth}
\begin{lstlisting}[style=apireference]
void validate() const;
\end{lstlisting}
{\footnotesize Validate the current public array and all spline metadata.\par}
\end{minipage}\par\vspace{1pt}
\par\noindent\begin{minipage}{\linewidth}
\begin{lstlisting}[style=apireference]
void set_parameters(std::valarray< REAL > parameter_values);
\end{lstlisting}
{\footnotesize Atomically replace the complete flattened parameter array.\par}
\end{minipage}\par\vspace{1pt}
\par\noindent\begin{minipage}{\linewidth}
\begin{lstlisting}[style=apireference]
void set_parameters(const REAL *parameter_values);
\end{lstlisting}
{\footnotesize Atomically copy and replace the complete flattened parameter array.\par}
\end{minipage}\par\vspace{1pt}
\par\noindent\begin{minipage}{\linewidth}
\begin{lstlisting}[style=apireference]
void set_control_point(std::size_t index, const point3< REAL > &point);
\end{lstlisting}
{\footnotesize Atomically replace one control point.\par}
\end{minipage}\par\vspace{1pt}
\par\noindent\begin{minipage}{\linewidth}
\begin{lstlisting}[style=apireference]
void set_weight(std::size_t index, REAL value);
\end{lstlisting}
{\footnotesize Atomically replace one rational weight.\par}
\end{minipage}\par\vspace{1pt}
\par\noindent\begin{minipage}{\linewidth}
\begin{lstlisting}[style=apireference]
void set_knot(std::size_t index, REAL value);
\end{lstlisting}
{\footnotesize Atomically replace one knot.\par}
\end{minipage}\par\vspace{1pt}
\par\noindent\begin{minipage}{\linewidth}
\begin{lstlisting}[style=apireference]
void set_control_points(std::vector< point3< REAL > > control_points);
\end{lstlisting}
{\footnotesize Atomically replace all control points while retaining other data.\par}
\end{minipage}\par\vspace{1pt}
\par\noindent\begin{minipage}{\linewidth}
\begin{lstlisting}[style=apireference]
void set_weights(std::vector< REAL > weight_values);
\end{lstlisting}
{\footnotesize Atomically replace all weights while retaining other data.\par}
\end{minipage}\par\vspace{1pt}
\par\noindent\begin{minipage}{\linewidth}
\begin{lstlisting}[style=apireference]
void set_knots(std::vector< REAL > knot_values);
\end{lstlisting}
{\footnotesize Atomically replace the complete knot vector.\par}
\end{minipage}\par\vspace{1pt}
\par\noindent\begin{minipage}{\linewidth}
\begin{lstlisting}[style=apireference]
void set_definition(std::vector< point3< REAL > > control_points, std::vector< REAL > weight_values, std::vector< REAL > knot_values, bool closed, REAL tolerance);
\end{lstlisting}
{\footnotesize Atomically replace all vector definition fields and metadata.\par}
\end{minipage}\par\vspace{1pt}
\par\noindent\begin{minipage}{\linewidth}
\begin{lstlisting}[style=apireference]
void set_tolerance(REAL tolerance);
\end{lstlisting}
{\footnotesize Atomically change the validation tolerance.\par}
\end{minipage}\par\vspace{1pt}
\par\noindent\begin{minipage}{\linewidth}
\begin{lstlisting}[style=apireference]
void set_closed(bool closed);
\end{lstlisting}
{\footnotesize Atomically change the closed-seam requirement.\par}
\end{minipage}\par\vspace{1pt}
\par\noindent\begin{minipage}{\linewidth}
\begin{lstlisting}[style=apireference]
void set_standard_knots(REAL s_end);
\end{lstlisting}
{\footnotesize Set standard open-clamped knots on [0, s\_end].\par}
\end{minipage}\par\vspace{1pt}
\par\noindent\begin{minipage}{\linewidth}
\begin{lstlisting}[style=apireference]
void set_standard_knots(REAL s_start, REAL s_end);
\end{lstlisting}
{\footnotesize Set standard open-clamped knots on [s\_start, s\_end].\par}
\end{minipage}\par\vspace{1pt}
\par\noindent\begin{minipage}{\linewidth}
\begin{lstlisting}[style=apireference]
point3< REAL > get_start() const;
\end{lstlisting}
{\footnotesize Compute the current point at the start of the active domain.\par}
{\footnotesize\textit{Returns:} Evaluated start point by value; no endpoint is cached.\par}
\end{minipage}\par\vspace{1pt}
\par\noindent\begin{minipage}{\linewidth}
\begin{lstlisting}[style=apireference]
point3< REAL > get_end() const;
\end{lstlisting}
{\footnotesize Compute the current point at the end of the active domain.\par}
{\footnotesize\textit{Returns:} Evaluated end point by value; no endpoint is cached.\par}
\end{minipage}\par\vspace{1pt}
\par\noindent\begin{minipage}{\linewidth}
\begin{lstlisting}[style=apireference]
REAL s_min() const;
\end{lstlisting}
{\footnotesize Get the current lower active parameter bound.\par}
{\footnotesize\textit{Returns:} Knot at index degree().\par}
\end{minipage}\par\vspace{1pt}
\par\noindent\begin{minipage}{\linewidth}
\begin{lstlisting}[style=apireference]
REAL s_max() const;
\end{lstlisting}
{\footnotesize Get the current upper active parameter bound.\par}
{\footnotesize\textit{Returns:} Knot at index control\_point\_count().\par}
\end{minipage}\par\vspace{1pt}
\par\noindent\begin{minipage}{\linewidth}
\begin{lstlisting}[style=apireference]
spline_derivatives3< REAL > derivatives_at(REAL s) const;
\end{lstlisting}
{\footnotesize Evaluate position and the first two analytic derivatives.\par}
{\footnotesize\textit{Returns:} Position and derivatives with respect to native s.\par}
\end{minipage}\par\vspace{1pt}
\par\noindent\begin{minipage}{\linewidth}
\begin{lstlisting}[style=apireference]
point3< REAL > evaluate(REAL s) const;
\end{lstlisting}
{\footnotesize Evaluate a world-space curve point.\par}
{\footnotesize\textit{Returns:} Evaluated point.\par}
\end{minipage}\par\vspace{1pt}
\par\noindent\begin{minipage}{\linewidth}
\begin{lstlisting}[style=apireference]
point3< REAL > point_at(REAL s) const;
\end{lstlisting}
{\footnotesize Evaluate a world-space curve point.\par}
{\footnotesize\textit{Returns:} Evaluated point.\par}
\end{minipage}\par\vspace{1pt}
\par\noindent\begin{minipage}{\linewidth}
\begin{lstlisting}[style=apireference]
vector3< REAL > first_derivative(REAL s) const;
\end{lstlisting}
{\footnotesize Evaluate the first analytic derivative.\par}
{\footnotesize\textit{Returns:} First derivative with respect to native s.\par}
\end{minipage}\par\vspace{1pt}
\par\noindent\begin{minipage}{\linewidth}
\begin{lstlisting}[style=apireference]
vector3< REAL > second_derivative(REAL s) const;
\end{lstlisting}
{\footnotesize Evaluate the second analytic derivative.\par}
{\footnotesize\textit{Returns:} Second derivative with respect to native s.\par}
\end{minipage}\par\vspace{1pt}
\par\noindent\begin{minipage}{\linewidth}
\begin{lstlisting}[style=apireference]
vector3< REAL > third_derivative(REAL s) const;
\end{lstlisting}
{\footnotesize Evaluate the third analytic rational derivative.\par}
{\footnotesize\textit{Returns:} Third derivative with respect to native s.\par}
\end{minipage}\par\vspace{1pt}
\par\noindent\begin{minipage}{\linewidth}
\begin{lstlisting}[style=apireference]
vector3< REAL > tangent(REAL s, REAL tangent_tolerance=vector3< REAL >::default_tolerance()) const;
\end{lstlisting}
{\footnotesize Evaluate a unit tangent.\par}
{\footnotesize\textit{Returns:} Normalized first derivative.\par}
\end{minipage}\par\vspace{1pt}
\par\noindent\begin{minipage}{\linewidth}
\begin{lstlisting}[style=apireference]
REAL curvature(REAL s, REAL tangent_tolerance=vector3< REAL >::default_tolerance()) const;
\end{lstlisting}
{\footnotesize Evaluate unsigned geometric curvature in world space.\par}
{\footnotesize\textit{Returns:} Nonnegative curvature magnitude in inverse world units.\par}
\end{minipage}\par\vspace{1pt}
\par\noindent\begin{minipage}{\linewidth}
\begin{lstlisting}[style=apireference]
REAL get_polygon_length() const;
\end{lstlisting}
{\footnotesize Get the total length of the current world-space control polygon.\par}
{\footnotesize\textit{Returns:} Control-polygon length in world units.\par}
\end{minipage}\par\vspace{1pt}
\par\noindent\begin{minipage}{\linewidth}
\begin{lstlisting}[style=apireference]
REAL approximate_arc_length(std::size_t segment_count=512) const;
\end{lstlisting}
{\footnotesize Estimate arc length using a uniform-native-parameter polyline.\par}
{\footnotesize\textit{Returns:} Approximate world-space arc length.\par}
\end{minipage}\par\vspace{1pt}
\par\noindent\begin{minipage}{\linewidth}
\begin{lstlisting}[style=apireference]
void adopt_to_points(const std::vector< point3< REAL > > &samples, const std::vector< REAL > &arc_length_parameters, std::size_t requested_degree=3, REAL tolerance=REAL(1e-10));
\end{lstlisting}
{\footnotesize Replace this spline with an open interpolant through samples.\par}
\end{minipage}\par\vspace{1pt}
\par\noindent\begin{minipage}{\linewidth}
\begin{lstlisting}[style=apireference]
void adopt_to_points(const std::vector< point3< REAL > > &samples, const std::vector< REAL > &arc_length_parameters, std::size_t requested_degree, bool closed, REAL tolerance=REAL(1e-10));
\end{lstlisting}
{\footnotesize Replace this spline with an open or closed interpolant.\par}
\end{minipage}\par\vspace{1pt}
\par\noindent\begin{minipage}{\linewidth}
\begin{lstlisting}[style=apireference]
nurbs_spline3< REAL > to_nurbs_spline() const;
\end{lstlisting}
{\footnotesize Create an independent ordinary 3D NURBS spline.\par}
{\footnotesize\textit{Returns:} Validated ordinary spline with identical geometry and metadata.\par}
\end{minipage}\par\vspace{1pt}
\par\noindent\begin{minipage}{\linewidth}
\begin{lstlisting}[style=apireference]
static nurbs_arr_spline3 interpolate(const std::vector< point3< REAL > > &samples, const std::vector< REAL > &arc_length_parameters, std::size_t requested_degree=3, REAL tolerance=REAL(1e-10));
\end{lstlisting}
{\footnotesize Build an open unit-weight spline interpolating measured samples.\par}
{\footnotesize\textit{Returns:} Self-contained interpolating array spline.\par}
\end{minipage}\par\vspace{1pt}
\par\noindent\begin{minipage}{\linewidth}
\begin{lstlisting}[style=apireference]
static nurbs_arr_spline3 interpolate(const std::vector< point3< REAL > > &samples, const std::vector< REAL > &arc_length_parameters, std::size_t requested_degree, bool closed, REAL tolerance=REAL(1e-10));
\end{lstlisting}
{\footnotesize Build an open or closed unit-weight interpolating array spline.\par}
{\footnotesize\textit{Returns:} Self-contained interpolating array spline.\par}
\end{minipage}\par\vspace{1pt}
\Needspace{9\baselineskip}
\section{nurbspath::nurbs\_spline2}
Rational B-spline curve in the independent 2D world.
\par\noindent\begin{minipage}{\linewidth}
\begin{lstlisting}[style=apireference]
nurbs_spline2(const point2< REAL > &start, const point2< REAL > &end, std::size_t control_point_count, std::size_t degree, REAL tolerance);
\end{lstlisting}
{\footnotesize Construct a standard open 2D spline with explicit tolerance.\par}
\end{minipage}\par\vspace{1pt}
\par\noindent\begin{minipage}{\linewidth}
\begin{lstlisting}[style=apireference]
template <std::same_as< bool > CLOSED = bool>
nurbs_spline2(const point2< REAL > &start, const point2< REAL > &end, std::size_t control_point_count, std::size_t degree, CLOSED closed=false, REAL tolerance=REAL(64) *std::numeric_limits< REAL >::epsilon());
\end{lstlisting}
{\footnotesize Construct a standard open or closed 2D spline between two points.\par}
\end{minipage}\par\vspace{1pt}
\par\noindent\begin{minipage}{\linewidth}
\begin{lstlisting}[style=apireference]
nurbs_spline2(std::vector< point2< REAL > > control_points, std::vector< REAL > weights, std::vector< REAL > knots, std::size_t degree, REAL tolerance=REAL(64) *std::numeric_limits< REAL >::epsilon());
\end{lstlisting}
{\footnotesize Construct a 2D NURBS curve from its complete definition.\par}
\end{minipage}\par\vspace{1pt}
\par\noindent\begin{minipage}{\linewidth}
\begin{lstlisting}[style=apireference]
nurbs_spline2(std::vector< point2< REAL > > control_points, std::vector< REAL > weights, std::vector< REAL > knots, std::size_t degree, bool closed, REAL tolerance=REAL(64) *std::numeric_limits< REAL >::epsilon());
\end{lstlisting}
{\footnotesize Construct an open or closed 2D NURBS curve.\par}
\end{minipage}\par\vspace{1pt}
\par\noindent\begin{minipage}{\linewidth}
\begin{lstlisting}[style=apireference]
nurbs_spline2(const nurbs_arr_spline2< REAL > &spline);
\end{lstlisting}
{\footnotesize Copy an array-backed 2D NURBS into ordinary vector storage.\par}
\end{minipage}\par\vspace{1pt}
\par\noindent\begin{minipage}{\linewidth}
\begin{lstlisting}[style=apireference]
const std::vector< point2< REAL > > & get_control_points() const noexcept;
\end{lstlisting}
{\footnotesize Get all 2D control points.\par}
{\footnotesize\textit{Returns:} Constant control-point vector reference.\par}
\end{minipage}\par\vspace{1pt}
\par\noindent\begin{minipage}{\linewidth}
\begin{lstlisting}[style=apireference]
const std::vector< REAL > & get_weights() const noexcept;
\end{lstlisting}
{\footnotesize Get all rational weights.\par}
{\footnotesize\textit{Returns:} Constant weight vector reference.\par}
\end{minipage}\par\vspace{1pt}
\par\noindent\begin{minipage}{\linewidth}
\begin{lstlisting}[style=apireference]
const std::vector< REAL > & get_knots() const noexcept;
\end{lstlisting}
{\footnotesize Get the complete knot vector.\par}
{\footnotesize\textit{Returns:} Constant knot vector reference.\par}
\end{minipage}\par\vspace{1pt}
\par\noindent\begin{minipage}{\linewidth}
\begin{lstlisting}[style=apireference]
std::size_t degree() const noexcept;
\end{lstlisting}
{\footnotesize Get the polynomial degree.\par}
{\footnotesize\textit{Returns:} Spline degree.\par}
\end{minipage}\par\vspace{1pt}
\par\noindent\begin{minipage}{\linewidth}
\begin{lstlisting}[style=apireference]
REAL tolerance() const noexcept;
\end{lstlisting}
{\footnotesize Get the spline definition tolerance.\par}
{\footnotesize\textit{Returns:} Positive tolerance.\par}
\end{minipage}\par\vspace{1pt}
\par\noindent\begin{minipage}{\linewidth}
\begin{lstlisting}[style=apireference]
bool is_closed() const noexcept;
\end{lstlisting}
{\footnotesize Report whether the spline has a closed seam.\par}
{\footnotesize\textit{Returns:} True for closed curves.\par}
\end{minipage}\par\vspace{1pt}
\par\noindent\begin{minipage}{\linewidth}
\begin{lstlisting}[style=apireference]
std::size_t number_of_parameters() const noexcept;
\end{lstlisting}
{\footnotesize Get the number of editable scalar parameters.\par}
{\footnotesize\textit{Returns:} Current number of indexed scalar parameters.\par}
\end{minipage}\par\vspace{1pt}
\par\noindent\begin{minipage}{\linewidth}
\begin{lstlisting}[style=apireference]
REAL get_parameter(std::size_t index) const;
\end{lstlisting}
{\footnotesize Get one editable scalar parameter by its flattened index.\par}
{\footnotesize\textit{Returns:} Current parameter value.\par}
\end{minipage}\par\vspace{1pt}
\par\noindent\begin{minipage}{\linewidth}
\begin{lstlisting}[style=apireference]
std::string parameter_name(std::size_t index) const;
\end{lstlisting}
{\footnotesize Get the descriptive name of one scalar parameter in the current definition.\par}
{\footnotesize\textit{Returns:} Parameter name in the documented flattened order.\par}
\end{minipage}\par\vspace{1pt}
\par\noindent\begin{minipage}{\linewidth}
\begin{lstlisting}[style=apireference]
void set_parameters(const REAL *parameter_values);
\end{lstlisting}
{\footnotesize Atomically replace all scalar parameters from a raw sequence.\par}
\end{minipage}\par\vspace{1pt}
\par\noindent\begin{minipage}{\linewidth}
\begin{lstlisting}[style=apireference]
bool set_parameter(std::size_t index, REAL value);
\end{lstlisting}
{\footnotesize Validate and atomically replace one scalar parameter.\par}
{\footnotesize\textit{Returns:} True when the complete candidate is valid and committed; false for an out-of-range index or rejected spline definition.\par}
\end{minipage}\par\vspace{1pt}
\par\noindent\begin{minipage}{\linewidth}
\begin{lstlisting}[style=apireference]
const point2< REAL > & get_control_point(std::size_t index) const;
\end{lstlisting}
{\footnotesize Get one control point by index.\par}
{\footnotesize\textit{Returns:} Constant reference to the selected control point.\par}
\end{minipage}\par\vspace{1pt}
\par\noindent\begin{minipage}{\linewidth}
\begin{lstlisting}[style=apireference]
REAL get_weight(std::size_t index) const;
\end{lstlisting}
{\footnotesize Get one rational weight by index.\par}
{\footnotesize\textit{Returns:} Selected rational weight.\par}
\end{minipage}\par\vspace{1pt}
\par\noindent\begin{minipage}{\linewidth}
\begin{lstlisting}[style=apireference]
REAL get_knot(std::size_t index) const;
\end{lstlisting}
{\footnotesize Get one knot by index.\par}
{\footnotesize\textit{Returns:} Selected knot value.\par}
\end{minipage}\par\vspace{1pt}
\par\noindent\begin{minipage}{\linewidth}
\begin{lstlisting}[style=apireference]
void set_control_point(std::size_t index, const point2< REAL > &point);
\end{lstlisting}
{\footnotesize Replace one control point through a checked atomic update.\par}
\end{minipage}\par\vspace{1pt}
\par\noindent\begin{minipage}{\linewidth}
\begin{lstlisting}[style=apireference]
void set_weight(std::size_t index, REAL weight);
\end{lstlisting}
{\footnotesize Replace one rational weight through a checked atomic update.\par}
\end{minipage}\par\vspace{1pt}
\par\noindent\begin{minipage}{\linewidth}
\begin{lstlisting}[style=apireference]
void set_knot(std::size_t index, REAL knot);
\end{lstlisting}
{\footnotesize Replace one knot through a checked atomic update.\par}
\end{minipage}\par\vspace{1pt}
\par\noindent\begin{minipage}{\linewidth}
\begin{lstlisting}[style=apireference]
void set_control_points(std::vector< point2< REAL > > control_points);
\end{lstlisting}
{\footnotesize Replace the complete control-point vector while retaining other fields.\par}
\end{minipage}\par\vspace{1pt}
\par\noindent\begin{minipage}{\linewidth}
\begin{lstlisting}[style=apireference]
void set_weights(std::vector< REAL > weights);
\end{lstlisting}
{\footnotesize Replace the complete weight vector while retaining other fields.\par}
\end{minipage}\par\vspace{1pt}
\par\noindent\begin{minipage}{\linewidth}
\begin{lstlisting}[style=apireference]
void set_knots(std::vector< REAL > knots);
\end{lstlisting}
{\footnotesize Replace the complete knot vector while retaining other fields.\par}
\end{minipage}\par\vspace{1pt}
\par\noindent\begin{minipage}{\linewidth}
\begin{lstlisting}[style=apireference]
void set_standard_knots(REAL s_end);
\end{lstlisting}
{\footnotesize Set standard open-clamped knots on the native domain [0, s\_end].\par}
\end{minipage}\par\vspace{1pt}
\par\noindent\begin{minipage}{\linewidth}
\begin{lstlisting}[style=apireference]
void set_standard_knots(REAL s_start, REAL s_end);
\end{lstlisting}
{\footnotesize Set standard open-clamped knots on [s\_start, s\_end].\par}
\end{minipage}\par\vspace{1pt}
\par\noindent\begin{minipage}{\linewidth}
\begin{lstlisting}[style=apireference]
void set_tolerance(REAL tolerance);
\end{lstlisting}
{\footnotesize Change the validation and parameter-boundary tolerance atomically.\par}
\end{minipage}\par\vspace{1pt}
\par\noindent\begin{minipage}{\linewidth}
\begin{lstlisting}[style=apireference]
void set_closed(bool closed);
\end{lstlisting}
{\footnotesize Change the closed-seam requirement atomically.\par}
\end{minipage}\par\vspace{1pt}
\par\noindent\begin{minipage}{\linewidth}
\begin{lstlisting}[style=apireference]
void set_definition(std::vector< point2< REAL > > control_points, std::vector< REAL > weights, std::vector< REAL > knots, bool closed, REAL tolerance);
\end{lstlisting}
{\footnotesize Atomically replace every definition field except the degree.\par}
\end{minipage}\par\vspace{1pt}
\par\noindent\begin{minipage}{\linewidth}
\begin{lstlisting}[style=apireference]
const point2< REAL > & get_start() const noexcept;
\end{lstlisting}
{\footnotesize Get the cached point at the start of the active domain.\par}
{\footnotesize\textit{Returns:} Constant reference to the point at s\_min() without evaluation.\par}
\end{minipage}\par\vspace{1pt}
\par\noindent\begin{minipage}{\linewidth}
\begin{lstlisting}[style=apireference]
const point2< REAL > & get_end() const noexcept;
\end{lstlisting}
{\footnotesize Get the cached point at the end of the active domain.\par}
{\footnotesize\textit{Returns:} Constant reference to the point at s\_max() without evaluation.\par}
\end{minipage}\par\vspace{1pt}
\par\noindent\begin{minipage}{\linewidth}
\begin{lstlisting}[style=apireference]
REAL s_min() const noexcept;
\end{lstlisting}
{\footnotesize Get the lower active parameter bound.\par}
{\footnotesize\textit{Returns:} Minimum valid s.\par}
\end{minipage}\par\vspace{1pt}
\par\noindent\begin{minipage}{\linewidth}
\begin{lstlisting}[style=apireference]
REAL s_max() const noexcept;
\end{lstlisting}
{\footnotesize Get the upper active parameter bound.\par}
{\footnotesize\textit{Returns:} Maximum valid s.\par}
\end{minipage}\par\vspace{1pt}
\par\noindent\begin{minipage}{\linewidth}
\begin{lstlisting}[style=apireference]
spline_derivatives2< REAL > derivatives_at(REAL s) const;
\end{lstlisting}
{\footnotesize Evaluate position and the first two analytic derivatives.\par}
{\footnotesize\textit{Returns:} Position, first derivative, and second derivative in 2D.\par}
\end{minipage}\par\vspace{1pt}
\par\noindent\begin{minipage}{\linewidth}
\begin{lstlisting}[style=apireference]
point2< REAL > evaluate(REAL s) const;
\end{lstlisting}
{\footnotesize Evaluate the curve position.\par}
{\footnotesize\textit{Returns:} Curve point in the 2D world.\par}
\end{minipage}\par\vspace{1pt}
\par\noindent\begin{minipage}{\linewidth}
\begin{lstlisting}[style=apireference]
point2< REAL > point_at(REAL s) const;
\end{lstlisting}
{\footnotesize Evaluate the curve position.\par}
{\footnotesize\textit{Returns:} Curve point in the 2D world.\par}
\end{minipage}\par\vspace{1pt}
\par\noindent\begin{minipage}{\linewidth}
\begin{lstlisting}[style=apireference]
vector2< REAL > first_derivative(REAL s) const;
\end{lstlisting}
{\footnotesize Evaluate the first derivative.\par}
{\footnotesize\textit{Returns:} Derivative with respect to s in 2D.\par}
\end{minipage}\par\vspace{1pt}
\par\noindent\begin{minipage}{\linewidth}
\begin{lstlisting}[style=apireference]
vector2< REAL > second_derivative(REAL s) const;
\end{lstlisting}
{\footnotesize Evaluate the second derivative.\par}
{\footnotesize\textit{Returns:} Second derivative with respect to s in 2D.\par}
\end{minipage}\par\vspace{1pt}
\par\noindent\begin{minipage}{\linewidth}
\begin{lstlisting}[style=apireference]
vector2< REAL > third_derivative(REAL s) const;
\end{lstlisting}
{\footnotesize Evaluate the third analytic rational derivative.\par}
{\footnotesize\textit{Returns:} Third derivative with respect to s in 2D.\par}
\end{minipage}\par\vspace{1pt}
\par\noindent\begin{minipage}{\linewidth}
\begin{lstlisting}[style=apireference]
vector2< REAL > tangent(REAL s, REAL tolerance=vector2< REAL >::default_tolerance()) const;
\end{lstlisting}
{\footnotesize Evaluate a unit tangent.\par}
{\footnotesize\textit{Returns:} Normalized first derivative.\par}
\end{minipage}\par\vspace{1pt}
\par\noindent\begin{minipage}{\linewidth}
\begin{lstlisting}[style=apireference]
REAL curvature(REAL s, REAL tangent_tolerance=vector2< REAL >::default_tolerance()) const;
\end{lstlisting}
{\footnotesize Evaluate the unsigned curvature magnitude.\par}
{\footnotesize\textit{Returns:} Nonnegative curvature magnitude at s.\par}
\end{minipage}\par\vspace{1pt}
\par\noindent\begin{minipage}{\linewidth}
\begin{lstlisting}[style=apireference]
REAL get_polygon_length() const noexcept;
\end{lstlisting}
{\footnotesize Get the total length of the 2D control polygon.\par}
{\footnotesize\textit{Returns:} Control-polygon length in 2D world units.\par}
\end{minipage}\par\vspace{1pt}
\par\noindent\begin{minipage}{\linewidth}
\begin{lstlisting}[style=apireference]
REAL approximate_arc_length(std::size_t segment_count=512) const;
\end{lstlisting}
{\footnotesize Estimate 2D arc length using a uniform-s polyline.\par}
{\footnotesize\textit{Returns:} Approximate length in 2D world units.\par}
\end{minipage}\par\vspace{1pt}
\par\noindent\begin{minipage}{\linewidth}
\begin{lstlisting}[style=apireference]
void adopt_to_points(const std::vector< point2< REAL > > &samples, const std::vector< REAL > &arc_length_parameters, std::size_t requested_degree=3, REAL tolerance=REAL(1e-10));
\end{lstlisting}
{\footnotesize Replace this curve with an interpolant through measured samples.\par}
\end{minipage}\par\vspace{1pt}
\par\noindent\begin{minipage}{\linewidth}
\begin{lstlisting}[style=apireference]
void adopt_to_points(const std::vector< point2< REAL > > &samples, const std::vector< REAL > &arc_length_parameters, std::size_t requested_degree, bool closed, REAL tolerance=REAL(1e-10));
\end{lstlisting}
{\footnotesize Replace this curve with an open or closed interpolant.\par}
\end{minipage}\par\vspace{1pt}
\par\noindent\begin{minipage}{\linewidth}
\begin{lstlisting}[style=apireference]
static nurbs_spline2 interpolate(const std::vector< point2< REAL > > &samples, const std::vector< REAL > &arc_length_parameters, std::size_t requested_degree=3, REAL tolerance=REAL(1e-10));
\end{lstlisting}
{\footnotesize Build a unit-weight B-spline interpolating measured 2D samples.\par}
{\footnotesize\textit{Returns:} Interpolating 2D NURBS curve with unit weights.\par}
\end{minipage}\par\vspace{1pt}
\par\noindent\begin{minipage}{\linewidth}
\begin{lstlisting}[style=apireference]
static nurbs_spline2 interpolate(const std::vector< point2< REAL > > &samples, const std::vector< REAL > &arc_length_parameters, std::size_t requested_degree, bool closed, REAL tolerance=REAL(1e-10));
\end{lstlisting}
{\footnotesize Build an open or closed unit-weight interpolating 2D spline.\par}
{\footnotesize\textit{Returns:} Interpolating 2D NURBS curve with unit weights.\par}
\end{minipage}\par\vspace{1pt}
\Needspace{9\baselineskip}
\section{nurbspath::nurbs\_spline3}
Rational B-spline curve embedded in the shared 3D world system.
\par\noindent\begin{minipage}{\linewidth}
\begin{lstlisting}[style=apireference]
nurbs_spline3(const point3< REAL > &start, const point3< REAL > &end, std::size_t control_point_count, std::size_t degree, REAL tolerance);
\end{lstlisting}
{\footnotesize Construct a standard open 3D spline with explicit tolerance.\par}
\end{minipage}\par\vspace{1pt}
\par\noindent\begin{minipage}{\linewidth}
\begin{lstlisting}[style=apireference]
template <std::same_as< bool > CLOSED = bool>
nurbs_spline3(const point3< REAL > &start, const point3< REAL > &end, std::size_t control_point_count, std::size_t degree, CLOSED closed=false, REAL tolerance=REAL(64) *std::numeric_limits< REAL >::epsilon());
\end{lstlisting}
{\footnotesize Construct a standard open or closed spline between two 3D points.\par}
\end{minipage}\par\vspace{1pt}
\par\noindent\begin{minipage}{\linewidth}
\begin{lstlisting}[style=apireference]
nurbs_spline3(std::vector< point3< REAL > > control_points, std::vector< REAL > weights, std::vector< REAL > knots, std::size_t degree, REAL tolerance=REAL(64) *std::numeric_limits< REAL >::epsilon());
\end{lstlisting}
{\footnotesize Construct a NURBS curve from its complete definition.\par}
\end{minipage}\par\vspace{1pt}
\par\noindent\begin{minipage}{\linewidth}
\begin{lstlisting}[style=apireference]
nurbs_spline3(std::vector< point3< REAL > > control_points, std::vector< REAL > weights, std::vector< REAL > knots, std::size_t degree, bool closed, REAL tolerance=REAL(64) *std::numeric_limits< REAL >::epsilon());
\end{lstlisting}
{\footnotesize Construct an open or closed 3D NURBS curve.\par}
\end{minipage}\par\vspace{1pt}
\par\noindent\begin{minipage}{\linewidth}
\begin{lstlisting}[style=apireference]
nurbs_spline3(const nurbs_arr_spline3< REAL > &spline);
\end{lstlisting}
{\footnotesize Copy an array-backed 3D NURBS into ordinary vector storage.\par}
\end{minipage}\par\vspace{1pt}
\par\noindent\begin{minipage}{\linewidth}
\begin{lstlisting}[style=apireference]
const std::vector< point3< REAL > > & get_control_points() const noexcept;
\end{lstlisting}
{\footnotesize Get all control points.\par}
{\footnotesize\textit{Returns:} Constant control-point vector reference.\par}
\end{minipage}\par\vspace{1pt}
\par\noindent\begin{minipage}{\linewidth}
\begin{lstlisting}[style=apireference]
const std::vector< REAL > & get_weights() const noexcept;
\end{lstlisting}
{\footnotesize Get all rational weights.\par}
{\footnotesize\textit{Returns:} Constant weight vector reference.\par}
\end{minipage}\par\vspace{1pt}
\par\noindent\begin{minipage}{\linewidth}
\begin{lstlisting}[style=apireference]
const std::vector< REAL > & get_knots() const noexcept;
\end{lstlisting}
{\footnotesize Get the full knot vector.\par}
{\footnotesize\textit{Returns:} Constant knot vector reference.\par}
\end{minipage}\par\vspace{1pt}
\par\noindent\begin{minipage}{\linewidth}
\begin{lstlisting}[style=apireference]
std::size_t degree() const noexcept;
\end{lstlisting}
{\footnotesize Get the polynomial degree.\par}
{\footnotesize\textit{Returns:} Spline degree.\par}
\end{minipage}\par\vspace{1pt}
\par\noindent\begin{minipage}{\linewidth}
\begin{lstlisting}[style=apireference]
REAL tolerance() const noexcept;
\end{lstlisting}
{\footnotesize Get the spline definition tolerance.\par}
{\footnotesize\textit{Returns:} Positive tolerance.\par}
\end{minipage}\par\vspace{1pt}
\par\noindent\begin{minipage}{\linewidth}
\begin{lstlisting}[style=apireference]
bool is_closed() const noexcept;
\end{lstlisting}
{\footnotesize Report whether the spline has a closed seam.\par}
{\footnotesize\textit{Returns:} True for closed curves.\par}
\end{minipage}\par\vspace{1pt}
\par\noindent\begin{minipage}{\linewidth}
\begin{lstlisting}[style=apireference]
std::size_t number_of_parameters() const noexcept;
\end{lstlisting}
{\footnotesize Get the number of editable scalar parameters.\par}
{\footnotesize\textit{Returns:} Current number of indexed scalar parameters.\par}
\end{minipage}\par\vspace{1pt}
\par\noindent\begin{minipage}{\linewidth}
\begin{lstlisting}[style=apireference]
REAL get_parameter(std::size_t index) const;
\end{lstlisting}
{\footnotesize Get one editable scalar parameter by its flattened index.\par}
{\footnotesize\textit{Returns:} Current parameter value.\par}
\end{minipage}\par\vspace{1pt}
\par\noindent\begin{minipage}{\linewidth}
\begin{lstlisting}[style=apireference]
std::string parameter_name(std::size_t index) const;
\end{lstlisting}
{\footnotesize Get the descriptive name of one scalar parameter in the current definition.\par}
{\footnotesize\textit{Returns:} Parameter name in the documented flattened order.\par}
\end{minipage}\par\vspace{1pt}
\par\noindent\begin{minipage}{\linewidth}
\begin{lstlisting}[style=apireference]
void set_parameters(const REAL *parameter_values);
\end{lstlisting}
{\footnotesize Atomically replace all scalar parameters from a raw sequence.\par}
\end{minipage}\par\vspace{1pt}
\par\noindent\begin{minipage}{\linewidth}
\begin{lstlisting}[style=apireference]
bool set_parameter(std::size_t index, REAL value);
\end{lstlisting}
{\footnotesize Validate and atomically replace one scalar parameter.\par}
{\footnotesize\textit{Returns:} True when the complete candidate is valid and committed; false for an out-of-range index or rejected spline definition.\par}
\end{minipage}\par\vspace{1pt}
\par\noindent\begin{minipage}{\linewidth}
\begin{lstlisting}[style=apireference]
const point3< REAL > & get_control_point(std::size_t index) const;
\end{lstlisting}
{\footnotesize Get one control point by index.\par}
{\footnotesize\textit{Returns:} Constant reference to the selected control point.\par}
\end{minipage}\par\vspace{1pt}
\par\noindent\begin{minipage}{\linewidth}
\begin{lstlisting}[style=apireference]
REAL get_weight(std::size_t index) const;
\end{lstlisting}
{\footnotesize Get one rational weight by index.\par}
{\footnotesize\textit{Returns:} Selected rational weight.\par}
\end{minipage}\par\vspace{1pt}
\par\noindent\begin{minipage}{\linewidth}
\begin{lstlisting}[style=apireference]
REAL get_knot(std::size_t index) const;
\end{lstlisting}
{\footnotesize Get one knot by index.\par}
{\footnotesize\textit{Returns:} Selected knot value.\par}
\end{minipage}\par\vspace{1pt}
\par\noindent\begin{minipage}{\linewidth}
\begin{lstlisting}[style=apireference]
void set_control_point(std::size_t index, const point3< REAL > &point);
\end{lstlisting}
{\footnotesize Replace one control point through a checked atomic update.\par}
\end{minipage}\par\vspace{1pt}
\par\noindent\begin{minipage}{\linewidth}
\begin{lstlisting}[style=apireference]
void set_weight(std::size_t index, REAL weight);
\end{lstlisting}
{\footnotesize Replace one rational weight through a checked atomic update.\par}
\end{minipage}\par\vspace{1pt}
\par\noindent\begin{minipage}{\linewidth}
\begin{lstlisting}[style=apireference]
void set_knot(std::size_t index, REAL knot);
\end{lstlisting}
{\footnotesize Replace one knot through a checked atomic update.\par}
\end{minipage}\par\vspace{1pt}
\par\noindent\begin{minipage}{\linewidth}
\begin{lstlisting}[style=apireference]
void set_control_points(std::vector< point3< REAL > > control_points);
\end{lstlisting}
{\footnotesize Replace the complete control-point vector while retaining other fields.\par}
\end{minipage}\par\vspace{1pt}
\par\noindent\begin{minipage}{\linewidth}
\begin{lstlisting}[style=apireference]
void set_weights(std::vector< REAL > weights);
\end{lstlisting}
{\footnotesize Replace the complete weight vector while retaining other fields.\par}
\end{minipage}\par\vspace{1pt}
\par\noindent\begin{minipage}{\linewidth}
\begin{lstlisting}[style=apireference]
void set_knots(std::vector< REAL > knots);
\end{lstlisting}
{\footnotesize Replace the complete knot vector while retaining other fields.\par}
\end{minipage}\par\vspace{1pt}
\par\noindent\begin{minipage}{\linewidth}
\begin{lstlisting}[style=apireference]
void set_standard_knots(REAL s_end);
\end{lstlisting}
{\footnotesize Set standard open-clamped knots on the native domain [0, s\_end].\par}
\end{minipage}\par\vspace{1pt}
\par\noindent\begin{minipage}{\linewidth}
\begin{lstlisting}[style=apireference]
void set_standard_knots(REAL s_start, REAL s_end);
\end{lstlisting}
{\footnotesize Set standard open-clamped knots on [s\_start, s\_end].\par}
\end{minipage}\par\vspace{1pt}
\par\noindent\begin{minipage}{\linewidth}
\begin{lstlisting}[style=apireference]
void set_tolerance(REAL tolerance);
\end{lstlisting}
{\footnotesize Change the validation and parameter-boundary tolerance atomically.\par}
\end{minipage}\par\vspace{1pt}
\par\noindent\begin{minipage}{\linewidth}
\begin{lstlisting}[style=apireference]
void set_closed(bool closed);
\end{lstlisting}
{\footnotesize Change the closed-seam requirement atomically.\par}
\end{minipage}\par\vspace{1pt}
\par\noindent\begin{minipage}{\linewidth}
\begin{lstlisting}[style=apireference]
void set_definition(std::vector< point3< REAL > > control_points, std::vector< REAL > weights, std::vector< REAL > knots, bool closed, REAL tolerance);
\end{lstlisting}
{\footnotesize Atomically replace every definition field except the degree.\par}
\end{minipage}\par\vspace{1pt}
\par\noindent\begin{minipage}{\linewidth}
\begin{lstlisting}[style=apireference]
std::ostream & tag_write(std::size_t tag, std::ostream &output) const;
\end{lstlisting}
{\footnotesize Write one version-1 tagged spline3 text record.\par}
{\footnotesize\textit{Returns:} Reference to output after attempting to write one newline-terminated row.\par}
\end{minipage}\par\vspace{1pt}
\par\noindent\begin{minipage}{\linewidth}
\begin{lstlisting}[style=apireference]
const point3< REAL > & get_start() const noexcept;
\end{lstlisting}
{\footnotesize Get the cached point at the start of the active domain.\par}
{\footnotesize\textit{Returns:} Constant reference to the point at s\_min() without evaluation.\par}
\end{minipage}\par\vspace{1pt}
\par\noindent\begin{minipage}{\linewidth}
\begin{lstlisting}[style=apireference]
const point3< REAL > & get_end() const noexcept;
\end{lstlisting}
{\footnotesize Get the cached point at the end of the active domain.\par}
{\footnotesize\textit{Returns:} Constant reference to the point at s\_max() without evaluation.\par}
\end{minipage}\par\vspace{1pt}
\par\noindent\begin{minipage}{\linewidth}
\begin{lstlisting}[style=apireference]
REAL s_min() const noexcept;
\end{lstlisting}
{\footnotesize Get the lower active parameter bound.\par}
{\footnotesize\textit{Returns:} Minimum valid s.\par}
\end{minipage}\par\vspace{1pt}
\par\noindent\begin{minipage}{\linewidth}
\begin{lstlisting}[style=apireference]
REAL s_max() const noexcept;
\end{lstlisting}
{\footnotesize Get the upper active parameter bound.\par}
{\footnotesize\textit{Returns:} Maximum valid s.\par}
\end{minipage}\par\vspace{1pt}
\par\noindent\begin{minipage}{\linewidth}
\begin{lstlisting}[style=apireference]
spline_derivatives3< REAL > derivatives_at(REAL s) const;
\end{lstlisting}
{\footnotesize Evaluate position and the first two derivatives together.\par}
{\footnotesize\textit{Returns:} Position, first derivative, and second derivative.\par}
\end{minipage}\par\vspace{1pt}
\par\noindent\begin{minipage}{\linewidth}
\begin{lstlisting}[style=apireference]
point3< REAL > evaluate(REAL s) const;
\end{lstlisting}
{\footnotesize Evaluate the curve position.\par}
{\footnotesize\textit{Returns:} World-space curve point.\par}
\end{minipage}\par\vspace{1pt}
\par\noindent\begin{minipage}{\linewidth}
\begin{lstlisting}[style=apireference]
point3< REAL > point_at(REAL s) const;
\end{lstlisting}
{\footnotesize Evaluate the curve position.\par}
{\footnotesize\textit{Returns:} World-space curve point.\par}
\end{minipage}\par\vspace{1pt}
\par\noindent\begin{minipage}{\linewidth}
\begin{lstlisting}[style=apireference]
vector3< REAL > first_derivative(REAL s) const;
\end{lstlisting}
{\footnotesize Evaluate the first derivative.\par}
{\footnotesize\textit{Returns:} Derivative with respect to s.\par}
\end{minipage}\par\vspace{1pt}
\par\noindent\begin{minipage}{\linewidth}
\begin{lstlisting}[style=apireference]
vector3< REAL > second_derivative(REAL s) const;
\end{lstlisting}
{\footnotesize Evaluate the second derivative.\par}
{\footnotesize\textit{Returns:} Second derivative with respect to s.\par}
\end{minipage}\par\vspace{1pt}
\par\noindent\begin{minipage}{\linewidth}
\begin{lstlisting}[style=apireference]
vector3< REAL > third_derivative(REAL s) const;
\end{lstlisting}
{\footnotesize Evaluate the third analytic rational derivative.\par}
{\footnotesize\textit{Returns:} Third derivative with respect to s.\par}
\end{minipage}\par\vspace{1pt}
\par\noindent\begin{minipage}{\linewidth}
\begin{lstlisting}[style=apireference]
vector3< REAL > tangent(REAL s, REAL tolerance=vector3< REAL >::default_tolerance()) const;
\end{lstlisting}
{\footnotesize Evaluate a unit tangent.\par}
{\footnotesize\textit{Returns:} Normalized first derivative.\par}
\end{minipage}\par\vspace{1pt}
\par\noindent\begin{minipage}{\linewidth}
\begin{lstlisting}[style=apireference]
REAL curvature(REAL s, REAL tangent_tolerance=vector3< REAL >::default_tolerance()) const;
\end{lstlisting}
{\footnotesize Evaluate the unsigned curvature magnitude.\par}
{\footnotesize\textit{Returns:} Nonnegative curvature magnitude at s.\par}
\end{minipage}\par\vspace{1pt}
\par\noindent\begin{minipage}{\linewidth}
\begin{lstlisting}[style=apireference]
REAL get_polygon_length() const noexcept;
\end{lstlisting}
{\footnotesize Get the total length of the world-space control polygon.\par}
{\footnotesize\textit{Returns:} Control-polygon length in world units.\par}
\end{minipage}\par\vspace{1pt}
\par\noindent\begin{minipage}{\linewidth}
\begin{lstlisting}[style=apireference]
REAL approximate_arc_length(std::size_t segment_count=512) const;
\end{lstlisting}
{\footnotesize Estimate total arc length using a uniform-s polyline.\par}
{\footnotesize\textit{Returns:} Approximate world-space arc length.\par}
\end{minipage}\par\vspace{1pt}
\par\noindent\begin{minipage}{\linewidth}
\begin{lstlisting}[style=apireference]
void adopt_to_points(const std::vector< point3< REAL > > &samples, const std::vector< REAL > &arc_length_parameters, std::size_t requested_degree=3, REAL tolerance=REAL(1e-10));
\end{lstlisting}
{\footnotesize Replace this curve with an interpolant through measured samples.\par}
\end{minipage}\par\vspace{1pt}
\par\noindent\begin{minipage}{\linewidth}
\begin{lstlisting}[style=apireference]
void adopt_to_points(const std::vector< point3< REAL > > &samples, const std::vector< REAL > &arc_length_parameters, std::size_t requested_degree, bool closed, REAL tolerance=REAL(1e-10));
\end{lstlisting}
{\footnotesize Replace this curve with an open or closed interpolant.\par}
\end{minipage}\par\vspace{1pt}
\par\noindent\begin{minipage}{\linewidth}
\begin{lstlisting}[style=apireference]
static std::optional< tagged_read_result< nurbs_spline3 > > tag_read(std::istream &input);
\end{lstlisting}
{\footnotesize Read and allocate one version-1 tagged spline3 text record.\par}
{\footnotesize\textit{Returns:} Tag and non-null shared spline after success, or std::nullopt at EOF, another read failure, or after a malformed row.\par}
\end{minipage}\par\vspace{1pt}
\par\noindent\begin{minipage}{\linewidth}
\begin{lstlisting}[style=apireference]
static nurbs_spline3 interpolate(const std::vector< point3< REAL > > &samples, const std::vector< REAL > &arc_length_parameters, std::size_t requested_degree=3, REAL tolerance=REAL(1e-10));
\end{lstlisting}
{\footnotesize Build a unit-weight B-spline interpolating measured samples.\par}
{\footnotesize\textit{Returns:} Interpolating NURBS curve with unit weights.\par}
\end{minipage}\par\vspace{1pt}
\par\noindent\begin{minipage}{\linewidth}
\begin{lstlisting}[style=apireference]
static nurbs_spline3 interpolate(const std::vector< point3< REAL > > &samples, const std::vector< REAL > &arc_length_parameters, std::size_t requested_degree, bool closed, REAL tolerance=REAL(1e-10));
\end{lstlisting}
{\footnotesize Build an open or closed unit-weight interpolating 3D spline.\par}
{\footnotesize\textit{Returns:} Interpolating 3D NURBS curve with unit weights.\par}
\end{minipage}\par\vspace{1pt}
\Needspace{9\baselineskip}
\section{nurbspath::plane3}
Infinite plane with an orthonormal (u,v) coordinate frame.
\par\noindent\begin{minipage}{\linewidth}
\begin{lstlisting}[style=apireference]
plane3(const vector3< REAL > &normal, REAL signed_distance, REAL tolerance=vector3< REAL >::default_tolerance());
\end{lstlisting}
{\footnotesize Construct a plane in Hessian normal form.\par}
\end{minipage}\par\vspace{1pt}
\par\noindent\begin{minipage}{\linewidth}
\begin{lstlisting}[style=apireference]
plane3(const point3< REAL > &origin, const vector3< REAL > &normal, REAL tolerance=vector3< REAL >::default_tolerance());
\end{lstlisting}
{\footnotesize Construct a plane from a parameter origin and normal.\par}
\end{minipage}\par\vspace{1pt}
\par\noindent\begin{minipage}{\linewidth}
\begin{lstlisting}[style=apireference]
plane3(const point3< REAL > &origin, const vector3< REAL > &normal, const vector3< REAL > &u_hint, REAL tolerance=vector3< REAL >::default_tolerance());
\end{lstlisting}
{\footnotesize Construct a plane with a preferred positive-u direction.\par}
\end{minipage}\par\vspace{1pt}
\par\noindent\begin{minipage}{\linewidth}
\begin{lstlisting}[style=apireference]
const point3< REAL > & origin() const noexcept;
\end{lstlisting}
{\footnotesize Get the plane parameter origin.\par}
{\footnotesize\textit{Returns:} Constant origin reference.\par}
\end{minipage}\par\vspace{1pt}
\par\noindent\begin{minipage}{\linewidth}
\begin{lstlisting}[style=apireference]
const vector3< REAL > & normal() const noexcept;
\end{lstlisting}
{\footnotesize Get the unit normal.\par}
{\footnotesize\textit{Returns:} Constant normal reference.\par}
\end{minipage}\par\vspace{1pt}
\par\noindent\begin{minipage}{\linewidth}
\begin{lstlisting}[style=apireference]
const vector3< REAL > & u_direction() const noexcept;
\end{lstlisting}
{\footnotesize Get the positive unit-u direction.\par}
{\footnotesize\textit{Returns:} Constant u direction reference.\par}
\end{minipage}\par\vspace{1pt}
\par\noindent\begin{minipage}{\linewidth}
\begin{lstlisting}[style=apireference]
const vector3< REAL > & v_direction() const noexcept;
\end{lstlisting}
{\footnotesize Get the positive unit-v direction.\par}
{\footnotesize\textit{Returns:} Constant v direction reference.\par}
\end{minipage}\par\vspace{1pt}
\par\noindent\begin{minipage}{\linewidth}
\begin{lstlisting}[style=apireference]
const point3< REAL > & get_origin() const noexcept;
\end{lstlisting}
{\footnotesize Get the plane parameter origin.\par}
{\footnotesize\textit{Returns:} Constant world-space origin reference.\par}
\end{minipage}\par\vspace{1pt}
\par\noindent\begin{minipage}{\linewidth}
\begin{lstlisting}[style=apireference]
const vector3< REAL > & get_normal() const noexcept;
\end{lstlisting}
{\footnotesize Get the unit normal.\par}
{\footnotesize\textit{Returns:} Constant normal reference.\par}
\end{minipage}\par\vspace{1pt}
\par\noindent\begin{minipage}{\linewidth}
\begin{lstlisting}[style=apireference]
const vector3< REAL > & get_u_direction() const noexcept;
\end{lstlisting}
{\footnotesize Get the positive unit-u direction.\par}
{\footnotesize\textit{Returns:} Constant u-direction reference.\par}
\end{minipage}\par\vspace{1pt}
\par\noindent\begin{minipage}{\linewidth}
\begin{lstlisting}[style=apireference]
const vector3< REAL > & get_v_direction() const noexcept;
\end{lstlisting}
{\footnotesize Get the positive unit-v direction.\par}
{\footnotesize\textit{Returns:} Constant v-direction reference.\par}
\end{minipage}\par\vspace{1pt}
\par\noindent\begin{minipage}{\linewidth}
\begin{lstlisting}[style=apireference]
void set_origin(const point3< REAL > &origin_value) noexcept;
\end{lstlisting}
{\footnotesize Replace the parameter origin while preserving the frame directions.\par}
\end{minipage}\par\vspace{1pt}
\par\noindent\begin{minipage}{\linewidth}
\begin{lstlisting}[style=apireference]
template <typename VECTOR>
void set_origin(const VECTOR &origin_value) noexcept;
\end{lstlisting}
{\footnotesize Replace the parameter origin using vector components.\par}
\end{minipage}\par\vspace{1pt}
\par\noindent\begin{minipage}{\linewidth}
\begin{lstlisting}[style=apireference]
void set_normal(const vector3< REAL > &normal_value, REAL tolerance=vector3< REAL >::default_tolerance());
\end{lstlisting}
{\footnotesize Replace the normal and generate a stable orthonormal plane frame.\par}
\end{minipage}\par\vspace{1pt}
\par\noindent\begin{minipage}{\linewidth}
\begin{lstlisting}[style=apireference]
void set_u_direction(const vector3< REAL > &u_hint, REAL tolerance=vector3< REAL >::default_tolerance());
\end{lstlisting}
{\footnotesize Replace positive u while preserving the origin and normal.\par}
\end{minipage}\par\vspace{1pt}
\par\noindent\begin{minipage}{\linewidth}
\begin{lstlisting}[style=apireference]
void set_definition(const vector3< REAL > &normal_value, REAL signed_distance, REAL tolerance=vector3< REAL >::default_tolerance());
\end{lstlisting}
{\footnotesize Atomically replace the plane from Hessian normal form.\par}
\end{minipage}\par\vspace{1pt}
\par\noindent\begin{minipage}{\linewidth}
\begin{lstlisting}[style=apireference]
void set_definition(const point3< REAL > &origin_value, const vector3< REAL > &normal_value, REAL tolerance=vector3< REAL >::default_tolerance());
\end{lstlisting}
{\footnotesize Atomically replace the origin, normal, and generated plane frame.\par}
\end{minipage}\par\vspace{1pt}
\par\noindent\begin{minipage}{\linewidth}
\begin{lstlisting}[style=apireference]
template <typename VECTOR>
void set_definition(const VECTOR &origin_value, const vector3< REAL > &normal_value, REAL tolerance=vector3< REAL >::default_tolerance());
\end{lstlisting}
{\footnotesize Atomically replace the origin from vector components and the frame.\par}
\end{minipage}\par\vspace{1pt}
\par\noindent\begin{minipage}{\linewidth}
\begin{lstlisting}[style=apireference]
void set_definition(const point3< REAL > &origin_value, const vector3< REAL > &normal_value, const vector3< REAL > &u_hint, REAL tolerance=vector3< REAL >::default_tolerance());
\end{lstlisting}
{\footnotesize Atomically replace the complete oriented plane definition.\par}
\end{minipage}\par\vspace{1pt}
\par\noindent\begin{minipage}{\linewidth}
\begin{lstlisting}[style=apireference]
template <typename VECTOR>
void set_definition(const VECTOR &origin_value, const vector3< REAL > &normal_value, const vector3< REAL > &u_hint, REAL tolerance=vector3< REAL >::default_tolerance());
\end{lstlisting}
{\footnotesize Atomically replace an oriented plane using a vector origin.\par}
\end{minipage}\par\vspace{1pt}
\par\noindent\begin{minipage}{\linewidth}
\begin{lstlisting}[style=apireference]
REAL signed_distance_from_origin() const noexcept;
\end{lstlisting}
{\footnotesize Get the equivalent Hessian signed distance.\par}
{\footnotesize\textit{Returns:} d for the equation unit\_normal dot x = d.\par}
\end{minipage}\par\vspace{1pt}
\par\noindent\begin{minipage}{\linewidth}
\begin{lstlisting}[style=apireference]
REAL get_signed_distance_from_origin() const noexcept;
\end{lstlisting}
{\footnotesize Get the equivalent Hessian signed distance.\par}
{\footnotesize\textit{Returns:} d for the equation unit\_normal dot x = d.\par}
\end{minipage}\par\vspace{1pt}
\par\noindent\begin{minipage}{\linewidth}
\begin{lstlisting}[style=apireference]
point3< REAL > point_at(REAL u, REAL v) const noexcept;
\end{lstlisting}
{\footnotesize Evaluate plane parameters.\par}
{\footnotesize\textit{Returns:} World-space plane point.\par}
\end{minipage}\par\vspace{1pt}
\par\noindent\begin{minipage}{\linewidth}
\begin{lstlisting}[style=apireference]
REAL signed_distance_to(const point3< REAL > &point) const noexcept;
\end{lstlisting}
{\footnotesize Compute oriented perpendicular distance to the plane.\par}
{\footnotesize\textit{Returns:} Signed distance, positive in the normal direction.\par}
\end{minipage}\par\vspace{1pt}
\par\noindent\begin{minipage}{\linewidth}
\begin{lstlisting}[style=apireference]
point3< REAL > project(const point3< REAL > &point) const noexcept;
\end{lstlisting}
{\footnotesize Orthogonally project a point onto the plane.\par}
{\footnotesize\textit{Returns:} Closest point on the plane.\par}
\end{minipage}\par\vspace{1pt}
\par\noindent\begin{minipage}{\linewidth}
\begin{lstlisting}[style=apireference]
std::pair< REAL, REAL > parameters_of(const point3< REAL > &point) const noexcept;
\end{lstlisting}
{\footnotesize Recover plane coordinates after orthogonal projection.\par}
{\footnotesize\textit{Returns:} Pair (u,v) in the plane frame.\par}
\end{minipage}\par\vspace{1pt}
\Needspace{9\baselineskip}
\section{nurbspath::point2}
Position in the independent two-dimensional Cartesian world.
\par\noindent\begin{minipage}{\linewidth}
\begin{lstlisting}[style=apireference]
using value_type = REAL;
\end{lstlisting}
{\footnotesize Floating-point scalar type.\par}
\end{minipage}\par\vspace{1pt}
\par\noindent\begin{minipage}{\linewidth}
\begin{lstlisting}[style=apireference]
REAL x = REAL(0);
\end{lstlisting}
{\footnotesize X coordinate in the 2D world.\par}
\end{minipage}\par\vspace{1pt}
\par\noindent\begin{minipage}{\linewidth}
\begin{lstlisting}[style=apireference]
REAL y = REAL(0);
\end{lstlisting}
{\footnotesize Y coordinate in the 2D world.\par}
\end{minipage}\par\vspace{1pt}
\par\noindent\begin{minipage}{\linewidth}
\begin{lstlisting}[style=apireference]
bool operator==(const point2 &other) const noexcept=default;
\end{lstlisting}
{\footnotesize Compare coordinates exactly.\par}
{\footnotesize\textit{Returns:} True when both coordinates are exactly equal.\par}
\end{minipage}\par\vspace{1pt}
\par\noindent\begin{minipage}{\linewidth}
\begin{lstlisting}[style=apireference]
template <typename VECTOR>
point2 & operator=(const VECTOR &vector_value) noexcept;
\end{lstlisting}
{\footnotesize Replace this point's coordinates with vector components.\par}
{\footnotesize\textit{Returns:} Reference to this updated point.\par}
\end{minipage}\par\vspace{1pt}
\par\noindent\begin{minipage}{\linewidth}
\begin{lstlisting}[style=apireference]
std::ostream & csv_write(std::ostream &output, text_format format=text_format::csv, std::streamsize decimal_places=6) const;
\end{lstlisting}
{\footnotesize Write the coordinates in CSV, TSV, or whitespace-delimited text.\par}
{\footnotesize\textit{Returns:} Reference to output.\par}
\end{minipage}\par\vspace{1pt}
\par\noindent\begin{minipage}{\linewidth}
\begin{lstlisting}[style=apireference]
std::istream & csv_read(std::istream &input, text_format format=text_format::csv);
\end{lstlisting}
{\footnotesize Read the coordinates from CSV, TSV, or whitespace-delimited text.\par}
{\footnotesize\textit{Returns:} Reference to input. This point changes only after a complete record.\par}
\end{minipage}\par\vspace{1pt}
\par\noindent\begin{minipage}{\linewidth}
\begin{lstlisting}[style=apireference]
std::ostream & write(std::ostream &output) const;
\end{lstlisting}
{\footnotesize Write two native REAL values in X/Y order as binary data.\par}
{\footnotesize\textit{Returns:} Reference to output.\par}
\end{minipage}\par\vspace{1pt}
\par\noindent\begin{minipage}{\linewidth}
\begin{lstlisting}[style=apireference]
std::istream & read(std::istream &input);
\end{lstlisting}
{\footnotesize Read two native REAL values in X/Y order from binary data.\par}
{\footnotesize\textit{Returns:} Reference to input. This point changes only after a complete record.\par}
\end{minipage}\par\vspace{1pt}
\par\noindent\begin{minipage}{\linewidth}
\begin{lstlisting}[style=apireference]
REAL & operator[](std::size_t index);
\end{lstlisting}
{\footnotesize Access a coordinate by index.\par}
{\footnotesize\textit{Returns:} Mutable reference to the selected coordinate.\par}
\end{minipage}\par\vspace{1pt}
\par\noindent\begin{minipage}{\linewidth}
\begin{lstlisting}[style=apireference]
const REAL & operator[](std::size_t index) const;
\end{lstlisting}
{\footnotesize Access a coordinate by index.\par}
{\footnotesize\textit{Returns:} Constant reference to the selected coordinate.\par}
\end{minipage}\par\vspace{1pt}
\par\noindent\begin{minipage}{\linewidth}
\begin{lstlisting}[style=apireference]
REAL magnitude() const noexcept;
\end{lstlisting}
{\footnotesize Compute distance from the Cartesian 2D origin.\par}
{\footnotesize\textit{Returns:} Euclidean magnitude in the coordinate units.\par}
\end{minipage}\par\vspace{1pt}
\par\noindent\begin{minipage}{\linewidth}
\begin{lstlisting}[style=apireference]
REAL manhattan_distance() const noexcept;
\end{lstlisting}
{\footnotesize Compute Manhattan distance from the Cartesian 2D origin.\par}
{\footnotesize\textit{Returns:} Sum of the absolute coordinate values in the coordinate units.\par}
\end{minipage}\par\vspace{1pt}
\par\noindent\begin{minipage}{\linewidth}
\begin{lstlisting}[style=apireference]
point2 & operator+=(const vector2< REAL > &offset) noexcept;
\end{lstlisting}
{\footnotesize Translate this point in place.\par}
{\footnotesize\textit{Returns:} Reference to this point.\par}
\end{minipage}\par\vspace{1pt}
\par\noindent\begin{minipage}{\linewidth}
\begin{lstlisting}[style=apireference]
point2 & operator-=(const vector2< REAL > &offset) noexcept;
\end{lstlisting}
{\footnotesize Translate this point by the negative of a vector.\par}
{\footnotesize\textit{Returns:} Reference to this point.\par}
\end{minipage}\par\vspace{1pt}
\par\noindent\begin{minipage}{\linewidth}
\begin{lstlisting}[style=apireference]
bool approximately_equal(const point2 &other, REAL tolerance=vector2< REAL >::default_tolerance()) const noexcept;
\end{lstlisting}
{\footnotesize Compare points using Euclidean distance.\par}
{\footnotesize\textit{Returns:} True when positions differ by at most tolerance.\par}
\end{minipage}\par\vspace{1pt}
\par\noindent\begin{minipage}{\linewidth}
\begin{lstlisting}[style=apireference]
static point2 origin() noexcept;
\end{lstlisting}
{\footnotesize Construct the Cartesian 2D origin.\par}
{\footnotesize\textit{Returns:} (0,0).\par}
\end{minipage}\par\vspace{1pt}
\Needspace{9\baselineskip}
\section{nurbspath::point3}
Position in the shared Cartesian world coordinate system.
\par\noindent\begin{minipage}{\linewidth}
\begin{lstlisting}[style=apireference]
using value_type = REAL;
\end{lstlisting}
{\footnotesize Floating-point scalar type.\par}
\end{minipage}\par\vspace{1pt}
\par\noindent\begin{minipage}{\linewidth}
\begin{lstlisting}[style=apireference]
REAL x = REAL(0);
\end{lstlisting}
{\footnotesize X coordinate in the 3D world.\par}
\end{minipage}\par\vspace{1pt}
\par\noindent\begin{minipage}{\linewidth}
\begin{lstlisting}[style=apireference]
REAL y = REAL(0);
\end{lstlisting}
{\footnotesize Y coordinate in the 3D world.\par}
\end{minipage}\par\vspace{1pt}
\par\noindent\begin{minipage}{\linewidth}
\begin{lstlisting}[style=apireference]
REAL z = REAL(0);
\end{lstlisting}
{\footnotesize Z coordinate in the 3D world.\par}
\end{minipage}\par\vspace{1pt}
\par\noindent\begin{minipage}{\linewidth}
\begin{lstlisting}[style=apireference]
bool operator==(const point3 &other) const noexcept=default;
\end{lstlisting}
{\footnotesize Compare coordinates exactly.\par}
{\footnotesize\textit{Returns:} True when all coordinates are exactly equal.\par}
\end{minipage}\par\vspace{1pt}
\par\noindent\begin{minipage}{\linewidth}
\begin{lstlisting}[style=apireference]
template <typename VECTOR>
point3 & operator=(const VECTOR &vector) noexcept;
\end{lstlisting}
{\footnotesize Replace this point's coordinates with vector components.\par}
{\footnotesize\textit{Returns:} Reference to this updated point.\par}
\end{minipage}\par\vspace{1pt}
\par\noindent\begin{minipage}{\linewidth}
\begin{lstlisting}[style=apireference]
std::ostream & csv_write(std::ostream &output, text_format format=text_format::csv, std::streamsize decimal_places=6) const;
\end{lstlisting}
{\footnotesize Write the coordinates in CSV, TSV, or whitespace-delimited text.\par}
{\footnotesize\textit{Returns:} Reference to output.\par}
\end{minipage}\par\vspace{1pt}
\par\noindent\begin{minipage}{\linewidth}
\begin{lstlisting}[style=apireference]
std::istream & csv_read(std::istream &input, text_format format=text_format::csv);
\end{lstlisting}
{\footnotesize Read the coordinates from CSV, TSV, or whitespace-delimited text.\par}
{\footnotesize\textit{Returns:} Reference to input. This point changes only after a complete record.\par}
\end{minipage}\par\vspace{1pt}
\par\noindent\begin{minipage}{\linewidth}
\begin{lstlisting}[style=apireference]
std::ostream & tag_write(std::size_t tag, std::ostream &output) const;
\end{lstlisting}
{\footnotesize Write one version-1 tagged point3 text record.\par}
{\footnotesize\textit{Returns:} Reference to output after attempting to write one newline-terminated row.\par}
\end{minipage}\par\vspace{1pt}
\par\noindent\begin{minipage}{\linewidth}
\begin{lstlisting}[style=apireference]
std::ostream & write(std::ostream &output) const;
\end{lstlisting}
{\footnotesize Write three native REAL values in X/Y/Z order as binary data.\par}
{\footnotesize\textit{Returns:} Reference to output.\par}
\end{minipage}\par\vspace{1pt}
\par\noindent\begin{minipage}{\linewidth}
\begin{lstlisting}[style=apireference]
std::istream & read(std::istream &input);
\end{lstlisting}
{\footnotesize Read three native REAL values in X/Y/Z order from binary data.\par}
{\footnotesize\textit{Returns:} Reference to input. This point changes only after a complete record.\par}
\end{minipage}\par\vspace{1pt}
\par\noindent\begin{minipage}{\linewidth}
\begin{lstlisting}[style=apireference]
REAL & operator[](std::size_t index);
\end{lstlisting}
{\footnotesize Access a coordinate by index.\par}
{\footnotesize\textit{Returns:} Mutable reference to the selected coordinate.\par}
\end{minipage}\par\vspace{1pt}
\par\noindent\begin{minipage}{\linewidth}
\begin{lstlisting}[style=apireference]
const REAL & operator[](std::size_t index) const;
\end{lstlisting}
{\footnotesize Access a coordinate by index.\par}
{\footnotesize\textit{Returns:} Constant reference to the selected coordinate.\par}
\end{minipage}\par\vspace{1pt}
\par\noindent\begin{minipage}{\linewidth}
\begin{lstlisting}[style=apireference]
REAL magnitude() const noexcept;
\end{lstlisting}
{\footnotesize Compute distance from the Cartesian world origin.\par}
{\footnotesize\textit{Returns:} Euclidean magnitude in the coordinate units.\par}
\end{minipage}\par\vspace{1pt}
\par\noindent\begin{minipage}{\linewidth}
\begin{lstlisting}[style=apireference]
REAL manhattan_distance() const noexcept;
\end{lstlisting}
{\footnotesize Compute Manhattan distance from the Cartesian world origin.\par}
{\footnotesize\textit{Returns:} Sum of the absolute coordinate values in the coordinate units.\par}
\end{minipage}\par\vspace{1pt}
\par\noindent\begin{minipage}{\linewidth}
\begin{lstlisting}[style=apireference]
point3 & operator+=(const vector3< REAL > &offset) noexcept;
\end{lstlisting}
{\footnotesize Translate this point in place.\par}
{\footnotesize\textit{Returns:} Reference to this point.\par}
\end{minipage}\par\vspace{1pt}
\par\noindent\begin{minipage}{\linewidth}
\begin{lstlisting}[style=apireference]
point3 & operator-=(const vector3< REAL > &offset) noexcept;
\end{lstlisting}
{\footnotesize Translate this point by the negative of a vector.\par}
{\footnotesize\textit{Returns:} Reference to this point.\par}
\end{minipage}\par\vspace{1pt}
\par\noindent\begin{minipage}{\linewidth}
\begin{lstlisting}[style=apireference]
bool approximately_equal(const point3 &other, REAL tolerance=vector3< REAL >::default_tolerance()) const noexcept;
\end{lstlisting}
{\footnotesize Compare points using Euclidean distance.\par}
{\footnotesize\textit{Returns:} True when positions differ by at most tolerance.\par}
\end{minipage}\par\vspace{1pt}
\par\noindent\begin{minipage}{\linewidth}
\begin{lstlisting}[style=apireference]
static std::optional< tagged_read_result< point3 > > tag_read(std::istream &input);
\end{lstlisting}
{\footnotesize Read and allocate one version-1 tagged point3 text record.\par}
{\footnotesize\textit{Returns:} Tag and non-null shared point after success, or std::nullopt at EOF, another read failure, or after a malformed row.\par}
\end{minipage}\par\vspace{1pt}
\par\noindent\begin{minipage}{\linewidth}
\begin{lstlisting}[style=apireference]
static point3 origin() noexcept;
\end{lstlisting}
{\footnotesize Construct the Cartesian world origin.\par}
{\footnotesize\textit{Returns:} (0,0,0).\par}
\end{minipage}\par\vspace{1pt}
\Needspace{9\baselineskip}
\section{nurbspath::point\_circle\_distance2}
Detailed 2D point-to-circle distance result.
\par\noindent\begin{minipage}{\linewidth}
\begin{lstlisting}[style=apireference]
REAL distance = REAL(0);
\end{lstlisting}
{\footnotesize Nonnegative tolerance-adjusted 2D distance.\par}
\end{minipage}\par\vspace{1pt}
\par\noindent\begin{minipage}{\linewidth}
\begin{lstlisting}[style=apireference]
point2< REAL > closest_point;
\end{lstlisting}
{\footnotesize Closest point on the 2D circle.\par}
\end{minipage}\par\vspace{1pt}
\par\noindent\begin{minipage}{\linewidth}
\begin{lstlisting}[style=apireference]
REAL u = REAL(0);
\end{lstlisting}
{\footnotesize Circle angle at closest\_point in [0,2*pi).\par}
\end{minipage}\par\vspace{1pt}
\Needspace{9\baselineskip}
\section{nurbspath::point\_spline\_distance2}
Detailed 2D point-to-spline distance result.
\par\noindent\begin{minipage}{\linewidth}
\begin{lstlisting}[style=apireference]
REAL distance = REAL(0);
\end{lstlisting}
{\footnotesize Nonnegative tolerance-adjusted 2D distance.\par}
\end{minipage}\par\vspace{1pt}
\par\noindent\begin{minipage}{\linewidth}
\begin{lstlisting}[style=apireference]
point2< REAL > closest_point;
\end{lstlisting}
{\footnotesize Closest sampled-and-refined spline point.\par}
\end{minipage}\par\vspace{1pt}
\par\noindent\begin{minipage}{\linewidth}
\begin{lstlisting}[style=apireference]
REAL s = REAL(0);
\end{lstlisting}
{\footnotesize Spline parameter at closest\_point.\par}
\end{minipage}\par\vspace{1pt}
\Needspace{9\baselineskip}
\section{nurbspath::point\_spline\_distance3}
Detailed point-to-spline distance result.
\par\noindent\begin{minipage}{\linewidth}
\begin{lstlisting}[style=apireference]
REAL distance = REAL(0);
\end{lstlisting}
{\footnotesize Nonnegative tolerance-adjusted distance.\par}
\end{minipage}\par\vspace{1pt}
\par\noindent\begin{minipage}{\linewidth}
\begin{lstlisting}[style=apireference]
point3< REAL > closest_point;
\end{lstlisting}
{\footnotesize Closest point found on the spline.\par}
\end{minipage}\par\vspace{1pt}
\par\noindent\begin{minipage}{\linewidth}
\begin{lstlisting}[style=apireference]
REAL s = REAL(0);
\end{lstlisting}
{\footnotesize Spline parameter at closest\_point.\par}
\end{minipage}\par\vspace{1pt}
\Needspace{9\baselineskip}
\section{nurbspath::point\_surface\_distance3}
Detailed point-to-surface distance result.
\par\noindent\begin{minipage}{\linewidth}
\begin{lstlisting}[style=apireference]
REAL distance = REAL(0);
\end{lstlisting}
{\footnotesize Nonnegative tolerance-adjusted distance.\par}
\end{minipage}\par\vspace{1pt}
\par\noindent\begin{minipage}{\linewidth}
\begin{lstlisting}[style=apireference]
point3< REAL > closest_point;
\end{lstlisting}
{\footnotesize Closest point on the queried surface.\par}
\end{minipage}\par\vspace{1pt}
\par\noindent\begin{minipage}{\linewidth}
\begin{lstlisting}[style=apireference]
REAL u = REAL(0);
\end{lstlisting}
{\footnotesize Surface u parameter at closest\_point.\par}
\end{minipage}\par\vspace{1pt}
\par\noindent\begin{minipage}{\linewidth}
\begin{lstlisting}[style=apireference]
REAL v = REAL(0);
\end{lstlisting}
{\footnotesize Surface v parameter at closest\_point.\par}
\end{minipage}\par\vspace{1pt}
\Needspace{9\baselineskip}
\section{nurbspath::ray2}
2D half-line parameterized as origin + s * direction for s >= 0.
\par\noindent\begin{minipage}{\linewidth}
\begin{lstlisting}[style=apireference]
ray2(const point2< REAL > &origin_value, const vector2< REAL > &direction_value, REAL tolerance=vector2< REAL >::default_tolerance());
\end{lstlisting}
{\footnotesize Construct a forward 2D ray.\par}
\end{minipage}\par\vspace{1pt}
\par\noindent\begin{minipage}{\linewidth}
\begin{lstlisting}[style=apireference]
const point2< REAL > & origin() const noexcept;
\end{lstlisting}
{\footnotesize Get the base point.\par}
{\footnotesize\textit{Returns:} Constant origin reference.\par}
\end{minipage}\par\vspace{1pt}
\par\noindent\begin{minipage}{\linewidth}
\begin{lstlisting}[style=apireference]
const vector2< REAL > & direction() const noexcept;
\end{lstlisting}
{\footnotesize Get the unnormalized direction.\par}
{\footnotesize\textit{Returns:} Constant direction reference.\par}
\end{minipage}\par\vspace{1pt}
\par\noindent\begin{minipage}{\linewidth}
\begin{lstlisting}[style=apireference]
const point2< REAL > & get_origin() const noexcept;
\end{lstlisting}
{\footnotesize Get the base point.\par}
{\footnotesize\textit{Returns:} Constant reference to the origin in the 2D world.\par}
\end{minipage}\par\vspace{1pt}
\par\noindent\begin{minipage}{\linewidth}
\begin{lstlisting}[style=apireference]
const vector2< REAL > & get_direction() const noexcept;
\end{lstlisting}
{\footnotesize Get the unnormalized direction.\par}
{\footnotesize\textit{Returns:} Constant reference to the parameter direction.\par}
\end{minipage}\par\vspace{1pt}
\par\noindent\begin{minipage}{\linewidth}
\begin{lstlisting}[style=apireference]
void set_origin(const point2< REAL > &origin_value) noexcept;
\end{lstlisting}
{\footnotesize Replace the origin while preserving the direction.\par}
\end{minipage}\par\vspace{1pt}
\par\noindent\begin{minipage}{\linewidth}
\begin{lstlisting}[style=apireference]
template <typename VECTOR>
void set_origin(const VECTOR &origin_value) noexcept;
\end{lstlisting}
{\footnotesize Replace the origin using vector components while preserving direction.\par}
\end{minipage}\par\vspace{1pt}
\par\noindent\begin{minipage}{\linewidth}
\begin{lstlisting}[style=apireference]
void set_direction(const vector2< REAL > &direction_value, REAL tolerance=vector2< REAL >::default_tolerance());
\end{lstlisting}
{\footnotesize Replace the unnormalized parameter direction.\par}
\end{minipage}\par\vspace{1pt}
\par\noindent\begin{minipage}{\linewidth}
\begin{lstlisting}[style=apireference]
void set_origin_and_direction(const point2< REAL > &origin_value, const vector2< REAL > &direction_value, REAL tolerance=vector2< REAL >::default_tolerance());
\end{lstlisting}
{\footnotesize Replace the origin and direction together.\par}
\end{minipage}\par\vspace{1pt}
\par\noindent\begin{minipage}{\linewidth}
\begin{lstlisting}[style=apireference]
template <typename VECTOR>
void set_origin_and_direction(const VECTOR &origin_value, const vector2< REAL > &direction_value, REAL tolerance=vector2< REAL >::default_tolerance());
\end{lstlisting}
{\footnotesize Replace the origin from vector components and replace the direction.\par}
\end{minipage}\par\vspace{1pt}
\par\noindent\begin{minipage}{\linewidth}
\begin{lstlisting}[style=apireference]
point2< REAL > point_at(REAL s) const noexcept;
\end{lstlisting}
{\footnotesize Evaluate the ray parameterization.\par}
{\footnotesize\textit{Returns:} origin + s * direction in the 2D world.\par}
\end{minipage}\par\vspace{1pt}
\par\noindent\begin{minipage}{\linewidth}
\begin{lstlisting}[style=apireference]
point2< REAL > evaluate(REAL s) const noexcept;
\end{lstlisting}
{\footnotesize Evaluate the ray parameterization.\par}
{\footnotesize\textit{Returns:} origin + s * direction in the 2D world.\par}
\end{minipage}\par\vspace{1pt}
\par\noindent\begin{minipage}{\linewidth}
\begin{lstlisting}[style=apireference]
vector2< REAL > tangent() const;
\end{lstlisting}
{\footnotesize Get the unit positive parameter direction.\par}
{\footnotesize\textit{Returns:} Normalized direction.\par}
\end{minipage}\par\vspace{1pt}
\Needspace{9\baselineskip}
\section{nurbspath::ray3}
Half-line parameterized as origin + s * direction for s >= 0.
\par\noindent\begin{minipage}{\linewidth}
\begin{lstlisting}[style=apireference]
ray3(const point3< REAL > &origin_value, const vector3< REAL > &direction_value, REAL tolerance=vector3< REAL >::default_tolerance());
\end{lstlisting}
{\footnotesize Construct a forward parametric ray.\par}
\end{minipage}\par\vspace{1pt}
\par\noindent\begin{minipage}{\linewidth}
\begin{lstlisting}[style=apireference]
const point3< REAL > & origin() const noexcept;
\end{lstlisting}
{\footnotesize Get the base point.\par}
{\footnotesize\textit{Returns:} Constant reference to the origin.\par}
\end{minipage}\par\vspace{1pt}
\par\noindent\begin{minipage}{\linewidth}
\begin{lstlisting}[style=apireference]
const vector3< REAL > & direction() const noexcept;
\end{lstlisting}
{\footnotesize Get the unnormalized parameter direction.\par}
{\footnotesize\textit{Returns:} Ray direction.\par}
\end{minipage}\par\vspace{1pt}
\par\noindent\begin{minipage}{\linewidth}
\begin{lstlisting}[style=apireference]
const point3< REAL > & get_origin() const noexcept;
\end{lstlisting}
{\footnotesize Get the base point.\par}
{\footnotesize\textit{Returns:} Constant reference to the origin.\par}
\end{minipage}\par\vspace{1pt}
\par\noindent\begin{minipage}{\linewidth}
\begin{lstlisting}[style=apireference]
const vector3< REAL > & get_direction() const noexcept;
\end{lstlisting}
{\footnotesize Get the unnormalized parameter direction.\par}
{\footnotesize\textit{Returns:} Constant reference to the ray direction.\par}
\end{minipage}\par\vspace{1pt}
\par\noindent\begin{minipage}{\linewidth}
\begin{lstlisting}[style=apireference]
void set_origin(const point3< REAL > &origin_value) noexcept;
\end{lstlisting}
{\footnotesize Replace the base point while preserving the direction.\par}
\end{minipage}\par\vspace{1pt}
\par\noindent\begin{minipage}{\linewidth}
\begin{lstlisting}[style=apireference]
template <typename VECTOR>
void set_origin(const VECTOR &origin_value) noexcept;
\end{lstlisting}
{\footnotesize Replace the base point using vector components.\par}
\end{minipage}\par\vspace{1pt}
\par\noindent\begin{minipage}{\linewidth}
\begin{lstlisting}[style=apireference]
void set_direction(const vector3< REAL > &direction_value, REAL tolerance=vector3< REAL >::default_tolerance());
\end{lstlisting}
{\footnotesize Replace the unnormalized parameter direction.\par}
\end{minipage}\par\vspace{1pt}
\par\noindent\begin{minipage}{\linewidth}
\begin{lstlisting}[style=apireference]
void set_origin_and_direction(const point3< REAL > &origin_value, const vector3< REAL > &direction_value, REAL tolerance=vector3< REAL >::default_tolerance());
\end{lstlisting}
{\footnotesize Atomically replace the base point and parameter direction.\par}
\end{minipage}\par\vspace{1pt}
\par\noindent\begin{minipage}{\linewidth}
\begin{lstlisting}[style=apireference]
template <typename VECTOR>
void set_origin_and_direction(const VECTOR &origin_value, const vector3< REAL > &direction_value, REAL tolerance=vector3< REAL >::default_tolerance());
\end{lstlisting}
{\footnotesize Atomically replace the base point from vector components and the parameter direction.\par}
\end{minipage}\par\vspace{1pt}
\par\noindent\begin{minipage}{\linewidth}
\begin{lstlisting}[style=apireference]
point3< REAL > point_at(REAL s) const noexcept;
\end{lstlisting}
{\footnotesize Evaluate the ray parameterization.\par}
{\footnotesize\textit{Returns:} origin + s * direction.\par}
\end{minipage}\par\vspace{1pt}
\par\noindent\begin{minipage}{\linewidth}
\begin{lstlisting}[style=apireference]
point3< REAL > evaluate(REAL s) const noexcept;
\end{lstlisting}
{\footnotesize Evaluate the ray parameterization.\par}
{\footnotesize\textit{Returns:} origin + s * direction.\par}
\end{minipage}\par\vspace{1pt}
\par\noindent\begin{minipage}{\linewidth}
\begin{lstlisting}[style=apireference]
vector3< REAL > tangent() const;
\end{lstlisting}
{\footnotesize Get the unit positive parameter direction.\par}
{\footnotesize\textit{Returns:} Normalized ray direction.\par}
\end{minipage}\par\vspace{1pt}
\Needspace{9\baselineskip}
\section{nurbspath::sphere3}
Sphere in world coordinates.
\par\noindent\begin{minipage}{\linewidth}
\begin{lstlisting}[style=apireference]
sphere3(const point3< REAL > &center_value, REAL radius_value);
\end{lstlisting}
{\footnotesize Construct a sphere.\par}
\end{minipage}\par\vspace{1pt}
\par\noindent\begin{minipage}{\linewidth}
\begin{lstlisting}[style=apireference]
const point3< REAL > & center() const noexcept;
\end{lstlisting}
{\footnotesize Get the world-space center.\par}
{\footnotesize\textit{Returns:} Constant center reference.\par}
\end{minipage}\par\vspace{1pt}
\par\noindent\begin{minipage}{\linewidth}
\begin{lstlisting}[style=apireference]
REAL radius() const noexcept;
\end{lstlisting}
{\footnotesize Get the radius.\par}
{\footnotesize\textit{Returns:} Radius in world units.\par}
\end{minipage}\par\vspace{1pt}
\par\noindent\begin{minipage}{\linewidth}
\begin{lstlisting}[style=apireference]
const point3< REAL > & get_center() const noexcept;
\end{lstlisting}
{\footnotesize Get the world-space center.\par}
{\footnotesize\textit{Returns:} Constant center reference.\par}
\end{minipage}\par\vspace{1pt}
\par\noindent\begin{minipage}{\linewidth}
\begin{lstlisting}[style=apireference]
REAL get_radius() const noexcept;
\end{lstlisting}
{\footnotesize Get the radius.\par}
{\footnotesize\textit{Returns:} Radius in world units.\par}
\end{minipage}\par\vspace{1pt}
\par\noindent\begin{minipage}{\linewidth}
\begin{lstlisting}[style=apireference]
void set_center(const point3< REAL > &center_value) noexcept;
\end{lstlisting}
{\footnotesize Replace the world-space center while preserving the radius.\par}
\end{minipage}\par\vspace{1pt}
\par\noindent\begin{minipage}{\linewidth}
\begin{lstlisting}[style=apireference]
template <typename VECTOR>
void set_center(const VECTOR &center_value) noexcept;
\end{lstlisting}
{\footnotesize Replace the world-space center using vector components.\par}
\end{minipage}\par\vspace{1pt}
\par\noindent\begin{minipage}{\linewidth}
\begin{lstlisting}[style=apireference]
void set_radius(REAL radius_value);
\end{lstlisting}
{\footnotesize Replace the radius while preserving the center.\par}
\end{minipage}\par\vspace{1pt}
\par\noindent\begin{minipage}{\linewidth}
\begin{lstlisting}[style=apireference]
void set_center_and_radius(const point3< REAL > &center_value, REAL radius_value);
\end{lstlisting}
{\footnotesize Atomically replace the center and radius.\par}
\end{minipage}\par\vspace{1pt}
\par\noindent\begin{minipage}{\linewidth}
\begin{lstlisting}[style=apireference]
template <typename VECTOR>
void set_center_and_radius(const VECTOR &center_value, REAL radius_value);
\end{lstlisting}
{\footnotesize Atomically replace the center from vector components and the radius.\par}
\end{minipage}\par\vspace{1pt}
\par\noindent\begin{minipage}{\linewidth}
\begin{lstlisting}[style=apireference]
sphere3 & operator=(const vector3< REAL > &center_vector) noexcept;
\end{lstlisting}
{\footnotesize Replace the center with vector components while preserving radius.\par}
{\footnotesize\textit{Returns:} Reference to this updated sphere.\par}
\end{minipage}\par\vspace{1pt}
\par\noindent\begin{minipage}{\linewidth}
\begin{lstlisting}[style=apireference]
std::ostream & tag_write(std::size_t tag, std::ostream &output) const;
\end{lstlisting}
{\footnotesize Write one version-1 tagged sphere3 text record.\par}
{\footnotesize\textit{Returns:} Reference to output after attempting to write one newline-terminated row.\par}
\end{minipage}\par\vspace{1pt}
\par\noindent\begin{minipage}{\linewidth}
\begin{lstlisting}[style=apireference]
point3< REAL > point_at(REAL u, REAL v) const noexcept;
\end{lstlisting}
{\footnotesize Evaluate spherical surface parameters.\par}
{\footnotesize\textit{Returns:} World-space surface point.\par}
\end{minipage}\par\vspace{1pt}
\par\noindent\begin{minipage}{\linewidth}
\begin{lstlisting}[style=apireference]
vector3< REAL > normal_at(const point3< REAL > &point) const;
\end{lstlisting}
{\footnotesize Compute the outward radial unit normal at a point.\par}
{\footnotesize\textit{Returns:} Outward unit normal after radial projection.\par}
\end{minipage}\par\vspace{1pt}
\par\noindent\begin{minipage}{\linewidth}
\begin{lstlisting}[style=apireference]
std::pair< REAL, REAL > parameters_of(const point3< REAL > &point, REAL tolerance=REAL(0)) const;
\end{lstlisting}
{\footnotesize Recover surface coordinates by radial projection.\par}
{\footnotesize\textit{Returns:} Pair (u,v) containing longitude and latitude in radians.\par}
\end{minipage}\par\vspace{1pt}
\par\noindent\begin{minipage}{\linewidth}
\begin{lstlisting}[style=apireference]
static std::optional< tagged_read_result< sphere3 > > tag_read(std::istream &input);
\end{lstlisting}
{\footnotesize Read and allocate one version-1 tagged sphere3 text record.\par}
{\footnotesize\textit{Returns:} Tag and non-null shared sphere after success, or std::nullopt at EOF, another read failure, or after a malformed row.\par}
\end{minipage}\par\vspace{1pt}
\Needspace{9\baselineskip}
\section{nurbspath::spline\_derivatives2}
Position and first two derivatives at one 2D spline parameter.
\par\noindent\begin{minipage}{\linewidth}
\begin{lstlisting}[style=apireference]
point2< REAL > point;
\end{lstlisting}
{\footnotesize Evaluated position in the 2D world.\par}
\end{minipage}\par\vspace{1pt}
\par\noindent\begin{minipage}{\linewidth}
\begin{lstlisting}[style=apireference]
vector2< REAL > first;
\end{lstlisting}
{\footnotesize First derivative with respect to s.\par}
\end{minipage}\par\vspace{1pt}
\par\noindent\begin{minipage}{\linewidth}
\begin{lstlisting}[style=apireference]
vector2< REAL > second;
\end{lstlisting}
{\footnotesize Second derivative with respect to s.\par}
\end{minipage}\par\vspace{1pt}
\Needspace{9\baselineskip}
\section{nurbspath::spline\_derivatives3}
Position and first two derivatives evaluated at one spline parameter.
\par\noindent\begin{minipage}{\linewidth}
\begin{lstlisting}[style=apireference]
point3< REAL > point;
\end{lstlisting}
{\footnotesize Evaluated world-space position.\par}
\end{minipage}\par\vspace{1pt}
\par\noindent\begin{minipage}{\linewidth}
\begin{lstlisting}[style=apireference]
vector3< REAL > first;
\end{lstlisting}
{\footnotesize First derivative with respect to s.\par}
\end{minipage}\par\vspace{1pt}
\par\noindent\begin{minipage}{\linewidth}
\begin{lstlisting}[style=apireference]
vector3< REAL > second;
\end{lstlisting}
{\footnotesize Second derivative with respect to s.\par}
\end{minipage}\par\vspace{1pt}
\Needspace{9\baselineskip}
\section{nurbspath::svg\_canvas2}
Accumulate native 2D geometry and emit a flat SVG document.
\par\noindent\begin{minipage}{\linewidth}
\begin{lstlisting}[style=apireference]
svg_canvas2(const svg_view2< REAL > &view, const svg_graphics_options2< REAL > &options={});
\end{lstlisting}
{\footnotesize Construct an empty flat canvas containing the 2D XY reference.\par}
\end{minipage}\par\vspace{1pt}
\par\noindent\begin{minipage}{\linewidth}
\begin{lstlisting}[style=apireference]
void add(const point2< REAL > &point);
\end{lstlisting}
{\footnotesize Add a native 2D point as a black dot.\par}
\end{minipage}\par\vspace{1pt}
\par\noindent\begin{minipage}{\linewidth}
\begin{lstlisting}[style=apireference]
void add(const ray2< REAL > &ray);
\end{lstlisting}
{\footnotesize Add a native 2D ray using solid pale-red lines.\par}
\end{minipage}\par\vspace{1pt}
\par\noindent\begin{minipage}{\linewidth}
\begin{lstlisting}[style=apireference]
void add(const circle2< REAL > &circle);
\end{lstlisting}
{\footnotesize Add a native 2D circle and center using solid pale-blue lines.\par}
\end{minipage}\par\vspace{1pt}
\par\noindent\begin{minipage}{\linewidth}
\begin{lstlisting}[style=apireference]
void add(const nurbs_spline2< REAL > &spline);
\end{lstlisting}
{\footnotesize Add a native 2D spline with solid curve and control lines.\par}
\end{minipage}\par\vspace{1pt}
\par\noindent\begin{minipage}{\linewidth}
\begin{lstlisting}[style=apireference]
std::string svg() const;
\end{lstlisting}
{\footnotesize Serialize the flat canvas.\par}
{\footnotesize\textit{Returns:} Self-contained SVG XML string.\par}
\end{minipage}\par\vspace{1pt}
\par\noindent\begin{minipage}{\linewidth}
\begin{lstlisting}[style=apireference]
void write(std::ostream &output) const;
\end{lstlisting}
{\footnotesize Stream the complete flat SVG document.\par}
\end{minipage}\par\vspace{1pt}
\Needspace{9\baselineskip}
\section{nurbspath::svg\_document3}
Accumulate projected geometry and emit a self-contained SVG document.
\par\noindent\begin{minipage}{\linewidth}
\begin{lstlisting}[style=apireference]
svg_document3(const svg_view3< REAL > &view, const svg_graphics_options3< REAL > &options={});
\end{lstlisting}
{\footnotesize Construct an empty SVG scene containing the world XYZ reference.\par}
\end{minipage}\par\vspace{1pt}
\par\noindent\begin{minipage}{\linewidth}
\begin{lstlisting}[style=apireference]
void add(const point3< REAL > &point);
\end{lstlisting}
{\footnotesize Add a point as a black dot of diameter twice the line width.\par}
\end{minipage}\par\vspace{1pt}
\par\noindent\begin{minipage}{\linewidth}
\begin{lstlisting}[style=apireference]
void add(const ray3< REAL > &ray);
\end{lstlisting}
{\footnotesize Add a dashed supporting line, positive unit direction, and origin.\par}
\end{minipage}\par\vspace{1pt}
\par\noindent\begin{minipage}{\linewidth}
\begin{lstlisting}[style=apireference]
void add(const sphere3< REAL > &sphere);
\end{lstlisting}
{\footnotesize Add a sphere silhouette, visible/hidden equator, and center dot.\par}
\end{minipage}\par\vspace{1pt}
\par\noindent\begin{minipage}{\linewidth}
\begin{lstlisting}[style=apireference]
void add(const plane3< REAL > &plane);
\end{lstlisting}
{\footnotesize Add a plane square, positive parameter axes, normal, and origin.\par}
\end{minipage}\par\vspace{1pt}
\par\noindent\begin{minipage}{\linewidth}
\begin{lstlisting}[style=apireference]
void add(const nurbs_spline3< REAL > &spline);
\end{lstlisting}
{\footnotesize Add control geometry, a tessellated curve, and endpoint dots.\par}
\end{minipage}\par\vspace{1pt}
\par\noindent\begin{minipage}{\linewidth}
\begin{lstlisting}[style=apireference]
void add(const plane3< REAL > &plane, const point2< REAL > &point);
\end{lstlisting}
{\footnotesize Project and add a 2D point as a black point marker.\par}
\end{minipage}\par\vspace{1pt}
\par\noindent\begin{minipage}{\linewidth}
\begin{lstlisting}[style=apireference]
void add(const plane3< REAL > &plane, const ray2< REAL > &ray);
\end{lstlisting}
{\footnotesize Project and add a 2D ray with pale-red patterned strokes.\par}
\end{minipage}\par\vspace{1pt}
\par\noindent\begin{minipage}{\linewidth}
\begin{lstlisting}[style=apireference]
void add(const plane3< REAL > &plane, const circle2< REAL > &circle);
\end{lstlisting}
{\footnotesize Project and add a 2D circle as a patterned center/radius sphere.\par}
\end{minipage}\par\vspace{1pt}
\par\noindent\begin{minipage}{\linewidth}
\begin{lstlisting}[style=apireference]
void add(const plane3< REAL > &plane, const nurbs_spline2< REAL > &spline);
\end{lstlisting}
{\footnotesize Project and add a 2D spline using patterned black/violet strokes.\par}
\end{minipage}\par\vspace{1pt}
\par\noindent\begin{minipage}{\linewidth}
\begin{lstlisting}[style=apireference]
std::string svg() const;
\end{lstlisting}
{\footnotesize Serialize the complete scene.\par}
{\footnotesize\textit{Returns:} Self-contained SVG XML string.\par}
\end{minipage}\par\vspace{1pt}
\par\noindent\begin{minipage}{\linewidth}
\begin{lstlisting}[style=apireference]
void write(std::ostream &output) const;
\end{lstlisting}
{\footnotesize Stream the complete SVG document.\par}
\end{minipage}\par\vspace{1pt}
\Needspace{9\baselineskip}
\section{nurbspath::svg\_graphics\_options2}
Styling and tessellation controls for the flat 2D SVG canvas.
\par\noindent\begin{minipage}{\linewidth}
\begin{lstlisting}[style=apireference]
REAL line_width = REAL(1.5);
\end{lstlisting}
\end{minipage}\par\vspace{1pt}
\par\noindent\begin{minipage}{\linewidth}
\begin{lstlisting}[style=apireference]
std::size_t spline_segment_count = 128;
\end{lstlisting}
{\footnotesize Number of solid straight elements used to approximate a NURBS path.\par}
\end{minipage}\par\vspace{1pt}
\par\noindent\begin{minipage}{\linewidth}
\begin{lstlisting}[style=apireference]
void validate() const;
\end{lstlisting}
{\footnotesize Validate all flat-canvas graphics controls.\par}
\end{minipage}\par\vspace{1pt}
\Needspace{9\baselineskip}
\section{nurbspath::svg\_graphics\_options3}
Styling and tessellation controls for three-dimensional SVG diagnostics.
\par\noindent\begin{minipage}{\linewidth}
\begin{lstlisting}[style=apireference]
REAL line_width = REAL(1.5);
\end{lstlisting}
\end{minipage}\par\vspace{1pt}
\par\noindent\begin{minipage}{\linewidth}
\begin{lstlisting}[style=apireference]
std::size_t spline_segment_count = 128;
\end{lstlisting}
{\footnotesize Number of straight elements used to approximate a complete NURBS path.\par}
\end{minipage}\par\vspace{1pt}
\par\noindent\begin{minipage}{\linewidth}
\begin{lstlisting}[style=apireference]
std::size_t sphere_segment_count = 96;
\end{lstlisting}
{\footnotesize Number of straight elements used for the sphere outline and equator.\par}
\end{minipage}\par\vspace{1pt}
\par\noindent\begin{minipage}{\linewidth}
\begin{lstlisting}[style=apireference]
REAL plane_square_side_length = REAL(0.25);
\end{lstlisting}
\end{minipage}\par\vspace{1pt}
\par\noindent\begin{minipage}{\linewidth}
\begin{lstlisting}[style=apireference]
REAL plane_parameter_axis_length = REAL(0.5);
\end{lstlisting}
{\footnotesize Length of the positive-u and positive-v plane direction markers.\par}
\end{minipage}\par\vspace{1pt}
\par\noindent\begin{minipage}{\linewidth}
\begin{lstlisting}[style=apireference]
void validate() const;
\end{lstlisting}
{\footnotesize Validate all graphics controls.\par}
\end{minipage}\par\vspace{1pt}
\Needspace{9\baselineskip}
\section{nurbspath::svg\_view2}
Flat viewport definition for the two-dimensional SVG canvas.
\par\noindent\begin{minipage}{\linewidth}
\begin{lstlisting}[style=apireference]
svg_view2(const point2< REAL > &center, REAL viewport_width, REAL viewport_height, std::size_t image_width=800, std::size_t image_height=600);
\end{lstlisting}
{\footnotesize Construct a flat 2D viewport.\par}
\end{minipage}\par\vspace{1pt}
\par\noindent\begin{minipage}{\linewidth}
\begin{lstlisting}[style=apireference]
const point2< REAL > & center() const noexcept;
\end{lstlisting}
{\footnotesize Get the 2D world point at image center.\par}
{\footnotesize\textit{Returns:} Center reference.\par}
\end{minipage}\par\vspace{1pt}
\par\noindent\begin{minipage}{\linewidth}
\begin{lstlisting}[style=apireference]
REAL viewport_width() const noexcept;
\end{lstlisting}
{\footnotesize Get the minimum visible world width.\par}
{\footnotesize\textit{Returns:} Positive width.\par}
\end{minipage}\par\vspace{1pt}
\par\noindent\begin{minipage}{\linewidth}
\begin{lstlisting}[style=apireference]
REAL viewport_height() const noexcept;
\end{lstlisting}
{\footnotesize Get the minimum visible world height.\par}
{\footnotesize\textit{Returns:} Positive height.\par}
\end{minipage}\par\vspace{1pt}
\par\noindent\begin{minipage}{\linewidth}
\begin{lstlisting}[style=apireference]
std::size_t image_width() const noexcept;
\end{lstlisting}
{\footnotesize Get the SVG viewBox width.\par}
{\footnotesize\textit{Returns:} Positive width.\par}
\end{minipage}\par\vspace{1pt}
\par\noindent\begin{minipage}{\linewidth}
\begin{lstlisting}[style=apireference]
std::size_t image_height() const noexcept;
\end{lstlisting}
{\footnotesize Get the SVG viewBox height.\par}
{\footnotesize\textit{Returns:} Positive height.\par}
\end{minipage}\par\vspace{1pt}
\Needspace{9\baselineskip}
\section{nurbspath::svg\_view3}
Camera and viewport definition for a three-dimensional SVG view.
\par\noindent\begin{minipage}{\linewidth}
\begin{lstlisting}[style=apireference]
static svg_view3 orthographic(const point3< REAL > &eyepoint, const point3< REAL > &lookatpoint, REAL viewport_width, REAL viewport_height, std::size_t image_width=800, std::size_t image_height=600);
\end{lstlisting}
{\footnotesize Construct an orthographic camera.\par}
{\footnotesize\textit{Returns:} Validated orthographic view definition.\par}
\end{minipage}\par\vspace{1pt}
\par\noindent\begin{minipage}{\linewidth}
\begin{lstlisting}[style=apireference]
static svg_view3 perspective(const point3< REAL > &eyepoint, const point3< REAL > &lookatpoint, REAL vertical_view_angle, std::size_t image_width=800, std::size_t image_height=600, REAL near_distance=REAL(1e-4));
\end{lstlisting}
{\footnotesize Construct a perspective camera.\par}
{\footnotesize\textit{Returns:} Validated perspective view definition.\par}
\end{minipage}\par\vspace{1pt}
\par\noindent\begin{minipage}{\linewidth}
\begin{lstlisting}[style=apireference]
const point3< REAL > & eyepoint() const noexcept;
\end{lstlisting}
{\footnotesize Get camera location.\par}
{\footnotesize\textit{Returns:} Constant eyepoint reference.\par}
\end{minipage}\par\vspace{1pt}
\par\noindent\begin{minipage}{\linewidth}
\begin{lstlisting}[style=apireference]
const point3< REAL > & lookatpoint() const noexcept;
\end{lstlisting}
{\footnotesize Get center-of-view target.\par}
{\footnotesize\textit{Returns:} Constant look-at reference.\par}
\end{minipage}\par\vspace{1pt}
\par\noindent\begin{minipage}{\linewidth}
\begin{lstlisting}[style=apireference]
svg_projection3 projection() const noexcept;
\end{lstlisting}
{\footnotesize Get projection mode.\par}
{\footnotesize\textit{Returns:} Orthographic or perspective mode.\par}
\end{minipage}\par\vspace{1pt}
\par\noindent\begin{minipage}{\linewidth}
\begin{lstlisting}[style=apireference]
REAL viewport_width() const noexcept;
\end{lstlisting}
{\footnotesize Get orthographic world width.\par}
{\footnotesize\textit{Returns:} Width, or zero for perspective.\par}
\end{minipage}\par\vspace{1pt}
\par\noindent\begin{minipage}{\linewidth}
\begin{lstlisting}[style=apireference]
REAL viewport_height() const noexcept;
\end{lstlisting}
{\footnotesize Get orthographic world height.\par}
{\footnotesize\textit{Returns:} Height, or zero for perspective.\par}
\end{minipage}\par\vspace{1pt}
\par\noindent\begin{minipage}{\linewidth}
\begin{lstlisting}[style=apireference]
REAL vertical_view_angle() const noexcept;
\end{lstlisting}
{\footnotesize Get vertical perspective field angle.\par}
{\footnotesize\textit{Returns:} Radians, or zero for orthographic.\par}
\end{minipage}\par\vspace{1pt}
\par\noindent\begin{minipage}{\linewidth}
\begin{lstlisting}[style=apireference]
REAL near_distance() const noexcept;
\end{lstlisting}
{\footnotesize Get perspective near distance.\par}
{\footnotesize\textit{Returns:} Distance, or zero for orthographic.\par}
\end{minipage}\par\vspace{1pt}
\par\noindent\begin{minipage}{\linewidth}
\begin{lstlisting}[style=apireference]
std::size_t image_width() const noexcept;
\end{lstlisting}
{\footnotesize Get SVG viewBox width.\par}
{\footnotesize\textit{Returns:} Positive width.\par}
\end{minipage}\par\vspace{1pt}
\par\noindent\begin{minipage}{\linewidth}
\begin{lstlisting}[style=apireference]
std::size_t image_height() const noexcept;
\end{lstlisting}
{\footnotesize Get SVG viewBox height.\par}
{\footnotesize\textit{Returns:} Positive height.\par}
\end{minipage}\par\vspace{1pt}
\Needspace{9\baselineskip}
\section{nurbspath::tagged\_entity3}
Application tag and dynamically selected concrete 3D entity.
\par\noindent\begin{minipage}{\linewidth}
\begin{lstlisting}[style=apireference]
std::size_t tag = 0;
\end{lstlisting}
\end{minipage}\par\vspace{1pt}
\par\noindent\begin{minipage}{\linewidth}
\begin{lstlisting}[style=apireference]
entity3_ptr< REAL > entity;
\end{lstlisting}
{\footnotesize Entity pointer; successful reads always select a non-null shared owner.\par}
\end{minipage}\par\vspace{1pt}
\Needspace{9\baselineskip}
\section{nurbspath::tagged\_read\_result}
Result of reading one tagged row for a known concrete entity type.
\par\noindent\begin{minipage}{\linewidth}
\begin{lstlisting}[style=apireference]
std::size_t tag = 0;
\end{lstlisting}
\end{minipage}\par\vspace{1pt}
\par\noindent\begin{minipage}{\linewidth}
\begin{lstlisting}[style=apireference]
std::shared_ptr< ENTITY > entity;
\end{lstlisting}
{\footnotesize Shared owner; successful reads always provide a non-null pointer.\par}
\end{minipage}\par\vspace{1pt}
\Needspace{9\baselineskip}
\section{nurbspath::tagged\_record}
Application tag and one geometry or optional presentation record.
\par\noindent\begin{minipage}{\linewidth}
\begin{lstlisting}[style=apireference]
std::size_t tag = 0;
\end{lstlisting}
{\footnotesize Application-defined tag associating related rows.\par}
\end{minipage}\par\vspace{1pt}
\par\noindent\begin{minipage}{\linewidth}
\begin{lstlisting}[style=apireference]
tagged_record_ptr< REAL > value;
\end{lstlisting}
{\footnotesize Non-null owner after a successful read.\par}
\end{minipage}\par\vspace{1pt}
\Needspace{9\baselineskip}
\section{nurbspath::text\_label}
Optional sized text associated with an application tag.
\par\noindent\begin{minipage}{\linewidth}
\begin{lstlisting}[style=apireference]
REAL size = REAL(12);
\end{lstlisting}
{\footnotesize Positive text size in consumer display units.\par}
\end{minipage}\par\vspace{1pt}
\par\noindent\begin{minipage}{\linewidth}
\begin{lstlisting}[style=apireference]
std::string value;
\end{lstlisting}
{\footnotesize Label text, optionally including escaped line breaks.\par}
\end{minipage}\par\vspace{1pt}
\par\noindent\begin{minipage}{\linewidth}
\begin{lstlisting}[style=apireference]
std::ostream & tag_write(std::size_t tag, std::ostream &output) const;
\end{lstlisting}
{\footnotesize Write one version-1 text row with a quoted, escaped string.\par}
{\footnotesize\textit{Returns:} The destination stream, with its formatting preserved.\par}
\end{minipage}\par\vspace{1pt}
\par\noindent\begin{minipage}{\linewidth}
\begin{lstlisting}[style=apireference]
static std::optional< tagged_read_result< text_label > > tag_read(std::istream &input);
\end{lstlisting}
{\footnotesize Read one text v1 row and allocate a validated text label. Advisory application and diagnostic rows are skipped without decoding.\par}
{\footnotesize\textit{Returns:} Preserved tag and non-null label owner, or nullopt on EOF/failure. Invalid quoting, escapes, fields, or values set failbit after that row.\par}
\end{minipage}\par\vspace{1pt}
\Needspace{9\baselineskip}
\section{nurbspath::vector2}
Cartesian vector in the independent two-dimensional world.
\par\noindent\begin{minipage}{\linewidth}
\begin{lstlisting}[style=apireference]
using value_type = REAL;
\end{lstlisting}
{\footnotesize Floating-point scalar type.\par}
\end{minipage}\par\vspace{1pt}
\par\noindent\begin{minipage}{\linewidth}
\begin{lstlisting}[style=apireference]
REAL x = REAL(0);
\end{lstlisting}
{\footnotesize X component in the 2D world.\par}
\end{minipage}\par\vspace{1pt}
\par\noindent\begin{minipage}{\linewidth}
\begin{lstlisting}[style=apireference]
REAL y = REAL(0);
\end{lstlisting}
{\footnotesize Y component in the 2D world.\par}
\end{minipage}\par\vspace{1pt}
\par\noindent\begin{minipage}{\linewidth}
\begin{lstlisting}[style=apireference]
constexpr vector2() noexcept=default;
\end{lstlisting}
{\footnotesize Construct the zero vector.\par}
\end{minipage}\par\vspace{1pt}
\par\noindent\begin{minipage}{\linewidth}
\begin{lstlisting}[style=apireference]
constexpr vector2(REAL x_value, REAL y_value) noexcept;
\end{lstlisting}
{\footnotesize Construct a vector from Cartesian components.\par}
\end{minipage}\par\vspace{1pt}
\par\noindent\begin{minipage}{\linewidth}
\begin{lstlisting}[style=apireference]
constexpr vector2(const point2< REAL > &point_value) noexcept;
\end{lstlisting}
{\footnotesize Construct the displacement from the 2D origin to a point.\par}
\end{minipage}\par\vspace{1pt}
\par\noindent\begin{minipage}{\linewidth}
\begin{lstlisting}[style=apireference]
vector2(const circle2< REAL > &circle_value) noexcept;
\end{lstlisting}
{\footnotesize Construct the displacement from the 2D origin to a circle center.\par}
\end{minipage}\par\vspace{1pt}
\par\noindent\begin{minipage}{\linewidth}
\begin{lstlisting}[style=apireference]
vector2(const ray2< REAL > &ray_value) noexcept;
\end{lstlisting}
{\footnotesize Construct the displacement from the 2D origin to a ray origin.\par}
\end{minipage}\par\vspace{1pt}
\par\noindent\begin{minipage}{\linewidth}
\begin{lstlisting}[style=apireference]
bool operator==(const vector2 &other) const noexcept=default;
\end{lstlisting}
{\footnotesize Compare components exactly.\par}
{\footnotesize\textit{Returns:} True when both components are exactly equal.\par}
\end{minipage}\par\vspace{1pt}
\par\noindent\begin{minipage}{\linewidth}
\begin{lstlisting}[style=apireference]
std::ostream & csv_write(std::ostream &output, text_format format=text_format::csv, std::streamsize decimal_places=6) const;
\end{lstlisting}
{\footnotesize Write the components in CSV, TSV, or whitespace-delimited text.\par}
{\footnotesize\textit{Returns:} Reference to output.\par}
\end{minipage}\par\vspace{1pt}
\par\noindent\begin{minipage}{\linewidth}
\begin{lstlisting}[style=apireference]
std::istream & csv_read(std::istream &input, text_format format=text_format::csv);
\end{lstlisting}
{\footnotesize Read the components from CSV, TSV, or whitespace-delimited text.\par}
{\footnotesize\textit{Returns:} Reference to input. This vector changes only after a complete record.\par}
\end{minipage}\par\vspace{1pt}
\par\noindent\begin{minipage}{\linewidth}
\begin{lstlisting}[style=apireference]
std::ostream & write(std::ostream &output) const;
\end{lstlisting}
{\footnotesize Write two native REAL values in X/Y order as binary data.\par}
{\footnotesize\textit{Returns:} Reference to output.\par}
\end{minipage}\par\vspace{1pt}
\par\noindent\begin{minipage}{\linewidth}
\begin{lstlisting}[style=apireference]
std::istream & read(std::istream &input);
\end{lstlisting}
{\footnotesize Read two native REAL values in X/Y order from binary data.\par}
{\footnotesize\textit{Returns:} Reference to input. This vector changes only after a complete record.\par}
\end{minipage}\par\vspace{1pt}
\par\noindent\begin{minipage}{\linewidth}
\begin{lstlisting}[style=apireference]
REAL & operator[](std::size_t index);
\end{lstlisting}
{\footnotesize Access a component by index.\par}
{\footnotesize\textit{Returns:} Mutable reference to the selected component.\par}
\end{minipage}\par\vspace{1pt}
\par\noindent\begin{minipage}{\linewidth}
\begin{lstlisting}[style=apireference]
const REAL & operator[](std::size_t index) const;
\end{lstlisting}
{\footnotesize Access a component by index.\par}
{\footnotesize\textit{Returns:} Constant reference to the selected component.\par}
\end{minipage}\par\vspace{1pt}
\par\noindent\begin{minipage}{\linewidth}
\begin{lstlisting}[style=apireference]
vector2 operator+() const noexcept;
\end{lstlisting}
{\footnotesize Return this vector unchanged.\par}
{\footnotesize\textit{Returns:} A copy of this vector.\par}
\end{minipage}\par\vspace{1pt}
\par\noindent\begin{minipage}{\linewidth}
\begin{lstlisting}[style=apireference]
vector2 operator-() const noexcept;
\end{lstlisting}
{\footnotesize Negate every component.\par}
{\footnotesize\textit{Returns:} Negated vector.\par}
\end{minipage}\par\vspace{1pt}
\par\noindent\begin{minipage}{\linewidth}
\begin{lstlisting}[style=apireference]
vector2 & operator+=(const vector2 &other) noexcept;
\end{lstlisting}
{\footnotesize Add another vector in place.\par}
{\footnotesize\textit{Returns:} Reference to this vector.\par}
\end{minipage}\par\vspace{1pt}
\par\noindent\begin{minipage}{\linewidth}
\begin{lstlisting}[style=apireference]
vector2 & operator-=(const vector2 &other) noexcept;
\end{lstlisting}
{\footnotesize Subtract another vector in place.\par}
{\footnotesize\textit{Returns:} Reference to this vector.\par}
\end{minipage}\par\vspace{1pt}
\par\noindent\begin{minipage}{\linewidth}
\begin{lstlisting}[style=apireference]
vector2 & operator*=(REAL scalar) noexcept;
\end{lstlisting}
{\footnotesize Multiply every component in place.\par}
{\footnotesize\textit{Returns:} Reference to this vector.\par}
\end{minipage}\par\vspace{1pt}
\par\noindent\begin{minipage}{\linewidth}
\begin{lstlisting}[style=apireference]
vector2 & operator/=(REAL scalar);
\end{lstlisting}
{\footnotesize Divide every component in place.\par}
{\footnotesize\textit{Returns:} Reference to this vector.\par}
\end{minipage}\par\vspace{1pt}
\par\noindent\begin{minipage}{\linewidth}
\begin{lstlisting}[style=apireference]
REAL dot(const vector2 &other) const noexcept;
\end{lstlisting}
{\footnotesize Compute the scalar dot product.\par}
{\footnotesize\textit{Returns:} Dot product.\par}
\end{minipage}\par\vspace{1pt}
\par\noindent\begin{minipage}{\linewidth}
\begin{lstlisting}[style=apireference]
REAL cross(const vector2 &other) const noexcept;
\end{lstlisting}
{\footnotesize Compute the signed scalar 2D cross product.\par}
{\footnotesize\textit{Returns:} Z component of the corresponding embedded 3D cross product.\par}
\end{minipage}\par\vspace{1pt}
\par\noindent\begin{minipage}{\linewidth}
\begin{lstlisting}[style=apireference]
REAL length_squared() const noexcept;
\end{lstlisting}
{\footnotesize Compute squared Euclidean length.\par}
{\footnotesize\textit{Returns:} Squared length.\par}
\end{minipage}\par\vspace{1pt}
\par\noindent\begin{minipage}{\linewidth}
\begin{lstlisting}[style=apireference]
REAL length() const noexcept;
\end{lstlisting}
{\footnotesize Compute Euclidean length.\par}
{\footnotesize\textit{Returns:} Vector length.\par}
\end{minipage}\par\vspace{1pt}
\par\noindent\begin{minipage}{\linewidth}
\begin{lstlisting}[style=apireference]
bool is_near_zero(REAL tolerance=default_tolerance()) const noexcept;
\end{lstlisting}
{\footnotesize Test whether vector length is within a tolerance of zero.\par}
{\footnotesize\textit{Returns:} True when squared length is at most tolerance squared.\par}
\end{minipage}\par\vspace{1pt}
\par\noindent\begin{minipage}{\linewidth}
\begin{lstlisting}[style=apireference]
vector2 normalized(REAL tolerance=default_tolerance()) const;
\end{lstlisting}
{\footnotesize Return a unit vector in the same direction.\par}
{\footnotesize\textit{Returns:} Normalized vector.\par}
\end{minipage}\par\vspace{1pt}
\par\noindent\begin{minipage}{\linewidth}
\begin{lstlisting}[style=apireference]
void normalize();
\end{lstlisting}
{\footnotesize Normalize this vector in place using the default vector tolerance.\par}
\end{minipage}\par\vspace{1pt}
\par\noindent\begin{minipage}{\linewidth}
\begin{lstlisting}[style=apireference]
vector2 perpendicular_left() const noexcept;
\end{lstlisting}
{\footnotesize Rotate this vector counterclockwise by ninety degrees.\par}
{\footnotesize\textit{Returns:} Left perpendicular vector (-y,x).\par}
\end{minipage}\par\vspace{1pt}
\par\noindent\begin{minipage}{\linewidth}
\begin{lstlisting}[style=apireference]
vector2 perpendicular_right() const noexcept;
\end{lstlisting}
{\footnotesize Rotate this vector clockwise by ninety degrees.\par}
{\footnotesize\textit{Returns:} Right perpendicular vector (y,-x).\par}
\end{minipage}\par\vspace{1pt}
\par\noindent\begin{minipage}{\linewidth}
\begin{lstlisting}[style=apireference]
vector2 projected_onto(const vector2 &axis, REAL tolerance=default_tolerance()) const;
\end{lstlisting}
{\footnotesize Project this vector onto an axis.\par}
{\footnotesize\textit{Returns:} Component parallel to axis.\par}
\end{minipage}\par\vspace{1pt}
\par\noindent\begin{minipage}{\linewidth}
\begin{lstlisting}[style=apireference]
vector2 rejected_from(const vector2 &axis, REAL tolerance=default_tolerance()) const;
\end{lstlisting}
{\footnotesize Remove the component parallel to an axis.\par}
{\footnotesize\textit{Returns:} Component perpendicular to axis.\par}
\end{minipage}\par\vspace{1pt}
\par\noindent\begin{minipage}{\linewidth}
\begin{lstlisting}[style=apireference]
vector2 reflected(const vector2 &normal, REAL tolerance=default_tolerance()) const;
\end{lstlisting}
{\footnotesize Reflect this vector across a line with a supplied normal.\par}
{\footnotesize\textit{Returns:} Reflected vector.\par}
\end{minipage}\par\vspace{1pt}
\par\noindent\begin{minipage}{\linewidth}
\begin{lstlisting}[style=apireference]
REAL angle_to(const vector2 &other, REAL tolerance=default_tolerance()) const;
\end{lstlisting}
{\footnotesize Compute the unsigned angle to another vector.\par}
{\footnotesize\textit{Returns:} Angle in radians in [0,pi].\par}
\end{minipage}\par\vspace{1pt}
\par\noindent\begin{minipage}{\linewidth}
\begin{lstlisting}[style=apireference]
REAL signed_angle_to(const vector2 &other, REAL tolerance=default_tolerance()) const;
\end{lstlisting}
{\footnotesize Compute the signed counterclockwise angle to another vector.\par}
{\footnotesize\textit{Returns:} Angle in radians in [-pi,pi].\par}
\end{minipage}\par\vspace{1pt}
\par\noindent\begin{minipage}{\linewidth}
\begin{lstlisting}[style=apireference]
vector2 component_product(const vector2 &other) const noexcept;
\end{lstlisting}
{\footnotesize Multiply corresponding components.\par}
{\footnotesize\textit{Returns:} Component-wise product.\par}
\end{minipage}\par\vspace{1pt}
\par\noindent\begin{minipage}{\linewidth}
\begin{lstlisting}[style=apireference]
REAL max_component() const noexcept;
\end{lstlisting}
{\footnotesize Get the largest component.\par}
{\footnotesize\textit{Returns:} Largest component value.\par}
\end{minipage}\par\vspace{1pt}
\par\noindent\begin{minipage}{\linewidth}
\begin{lstlisting}[style=apireference]
REAL min_component() const noexcept;
\end{lstlisting}
{\footnotesize Get the smallest component.\par}
{\footnotesize\textit{Returns:} Smallest component value.\par}
\end{minipage}\par\vspace{1pt}
\par\noindent\begin{minipage}{\linewidth}
\begin{lstlisting}[style=apireference]
bool approximately_equal(const vector2 &other, REAL tolerance=default_tolerance()) const noexcept;
\end{lstlisting}
{\footnotesize Compare vectors using Euclidean distance.\par}
{\footnotesize\textit{Returns:} True when vectors differ by at most tolerance.\par}
\end{minipage}\par\vspace{1pt}
\par\noindent\begin{minipage}{\linewidth}
\begin{lstlisting}[style=apireference]
static vector2 zero() noexcept;
\end{lstlisting}
{\footnotesize Construct the zero vector.\par}
{\footnotesize\textit{Returns:} (0,0).\par}
\end{minipage}\par\vspace{1pt}
\par\noindent\begin{minipage}{\linewidth}
\begin{lstlisting}[style=apireference]
static vector2 unit_x() noexcept;
\end{lstlisting}
{\footnotesize Construct the positive X unit vector.\par}
{\footnotesize\textit{Returns:} (1,0).\par}
\end{minipage}\par\vspace{1pt}
\par\noindent\begin{minipage}{\linewidth}
\begin{lstlisting}[style=apireference]
static vector2 unit_y() noexcept;
\end{lstlisting}
{\footnotesize Construct the positive Y unit vector.\par}
{\footnotesize\textit{Returns:} (0,1).\par}
\end{minipage}\par\vspace{1pt}
\par\noindent\begin{minipage}{\linewidth}
\begin{lstlisting}[style=apireference]
static REAL default_tolerance() noexcept;
\end{lstlisting}
{\footnotesize Get the default absolute vector tolerance.\par}
{\footnotesize\textit{Returns:} 64 times machine epsilon for REAL.\par}
\end{minipage}\par\vspace{1pt}
\Needspace{9\baselineskip}
\section{nurbspath::vector3}
Cartesian three-dimensional vector.
\par\noindent\begin{minipage}{\linewidth}
\begin{lstlisting}[style=apireference]
using value_type = REAL;
\end{lstlisting}
{\footnotesize Floating-point scalar type.\par}
\end{minipage}\par\vspace{1pt}
\par\noindent\begin{minipage}{\linewidth}
\begin{lstlisting}[style=apireference]
REAL x = REAL(0);
\end{lstlisting}
{\footnotesize X component.\par}
\end{minipage}\par\vspace{1pt}
\par\noindent\begin{minipage}{\linewidth}
\begin{lstlisting}[style=apireference]
REAL y = REAL(0);
\end{lstlisting}
{\footnotesize Y component.\par}
\end{minipage}\par\vspace{1pt}
\par\noindent\begin{minipage}{\linewidth}
\begin{lstlisting}[style=apireference]
REAL z = REAL(0);
\end{lstlisting}
{\footnotesize Z component.\par}
\end{minipage}\par\vspace{1pt}
\par\noindent\begin{minipage}{\linewidth}
\begin{lstlisting}[style=apireference]
constexpr vector3() noexcept=default;
\end{lstlisting}
{\footnotesize Construct the zero vector.\par}
\end{minipage}\par\vspace{1pt}
\par\noindent\begin{minipage}{\linewidth}
\begin{lstlisting}[style=apireference]
constexpr vector3(REAL x_value, REAL y_value, REAL z_value) noexcept;
\end{lstlisting}
{\footnotesize Construct a vector from X, Y, and Z components.\par}
\end{minipage}\par\vspace{1pt}
\par\noindent\begin{minipage}{\linewidth}
\begin{lstlisting}[style=apireference]
constexpr vector3(const point3< REAL > &point_value) noexcept;
\end{lstlisting}
{\footnotesize Construct a vector containing a point's coordinates.\par}
\end{minipage}\par\vspace{1pt}
\par\noindent\begin{minipage}{\linewidth}
\begin{lstlisting}[style=apireference]
vector3(const sphere3< REAL > &sphere_value) noexcept;
\end{lstlisting}
{\footnotesize Construct a vector containing a sphere center's coordinates.\par}
\end{minipage}\par\vspace{1pt}
\par\noindent\begin{minipage}{\linewidth}
\begin{lstlisting}[style=apireference]
vector3(const plane3< REAL > &plane_value) noexcept;
\end{lstlisting}
{\footnotesize Construct a vector containing a plane origin's coordinates.\par}
\end{minipage}\par\vspace{1pt}
\par\noindent\begin{minipage}{\linewidth}
\begin{lstlisting}[style=apireference]
vector3(const ray3< REAL > &ray_value) noexcept;
\end{lstlisting}
{\footnotesize Construct a vector containing a ray origin's coordinates.\par}
\end{minipage}\par\vspace{1pt}
\par\noindent\begin{minipage}{\linewidth}
\begin{lstlisting}[style=apireference]
bool operator==(const vector3 &other) const noexcept=default;
\end{lstlisting}
{\footnotesize Compare components exactly.\par}
{\footnotesize\textit{Returns:} True when all components are exactly equal.\par}
\end{minipage}\par\vspace{1pt}
\par\noindent\begin{minipage}{\linewidth}
\begin{lstlisting}[style=apireference]
std::ostream & csv_write(std::ostream &output, text_format format=text_format::csv, std::streamsize decimal_places=6) const;
\end{lstlisting}
{\footnotesize Write the components in CSV, TSV, or whitespace-delimited text.\par}
{\footnotesize\textit{Returns:} Reference to output.\par}
\end{minipage}\par\vspace{1pt}
\par\noindent\begin{minipage}{\linewidth}
\begin{lstlisting}[style=apireference]
std::istream & csv_read(std::istream &input, text_format format=text_format::csv);
\end{lstlisting}
{\footnotesize Read the components from CSV, TSV, or whitespace-delimited text.\par}
{\footnotesize\textit{Returns:} Reference to input. This vector changes only after a complete record.\par}
\end{minipage}\par\vspace{1pt}
\par\noindent\begin{minipage}{\linewidth}
\begin{lstlisting}[style=apireference]
std::ostream & tag_write(std::size_t tag, std::ostream &output) const;
\end{lstlisting}
{\footnotesize Write one version-1 tagged vector3 text record.\par}
{\footnotesize\textit{Returns:} Reference to output after attempting to write one newline-terminated row.\par}
\end{minipage}\par\vspace{1pt}
\par\noindent\begin{minipage}{\linewidth}
\begin{lstlisting}[style=apireference]
std::ostream & write(std::ostream &output) const;
\end{lstlisting}
{\footnotesize Write three native REAL values in X/Y/Z order as binary data.\par}
{\footnotesize\textit{Returns:} Reference to output.\par}
\end{minipage}\par\vspace{1pt}
\par\noindent\begin{minipage}{\linewidth}
\begin{lstlisting}[style=apireference]
std::istream & read(std::istream &input);
\end{lstlisting}
{\footnotesize Read three native REAL values in X/Y/Z order from binary data.\par}
{\footnotesize\textit{Returns:} Reference to input. This vector changes only after a complete record.\par}
\end{minipage}\par\vspace{1pt}
\par\noindent\begin{minipage}{\linewidth}
\begin{lstlisting}[style=apireference]
REAL & operator[](std::size_t index);
\end{lstlisting}
{\footnotesize Access a component by index.\par}
{\footnotesize\textit{Returns:} Mutable reference to the selected component.\par}
\end{minipage}\par\vspace{1pt}
\par\noindent\begin{minipage}{\linewidth}
\begin{lstlisting}[style=apireference]
const REAL & operator[](std::size_t index) const;
\end{lstlisting}
{\footnotesize Access a component by index.\par}
{\footnotesize\textit{Returns:} Constant reference to the selected component.\par}
\end{minipage}\par\vspace{1pt}
\par\noindent\begin{minipage}{\linewidth}
\begin{lstlisting}[style=apireference]
vector3 operator+() const noexcept;
\end{lstlisting}
{\footnotesize Return this vector unchanged.\par}
{\footnotesize\textit{Returns:} A copy of this vector.\par}
\end{minipage}\par\vspace{1pt}
\par\noindent\begin{minipage}{\linewidth}
\begin{lstlisting}[style=apireference]
vector3 operator-() const noexcept;
\end{lstlisting}
{\footnotesize Negate every component.\par}
{\footnotesize\textit{Returns:} Negated vector.\par}
\end{minipage}\par\vspace{1pt}
\par\noindent\begin{minipage}{\linewidth}
\begin{lstlisting}[style=apireference]
vector3 & operator+=(const vector3 &other) noexcept;
\end{lstlisting}
{\footnotesize Add another vector in place.\par}
{\footnotesize\textit{Returns:} Reference to this vector.\par}
\end{minipage}\par\vspace{1pt}
\par\noindent\begin{minipage}{\linewidth}
\begin{lstlisting}[style=apireference]
vector3 & operator-=(const vector3 &other) noexcept;
\end{lstlisting}
{\footnotesize Subtract another vector in place.\par}
{\footnotesize\textit{Returns:} Reference to this vector.\par}
\end{minipage}\par\vspace{1pt}
\par\noindent\begin{minipage}{\linewidth}
\begin{lstlisting}[style=apireference]
vector3 & operator*=(REAL scalar) noexcept;
\end{lstlisting}
{\footnotesize Multiply every component in place.\par}
{\footnotesize\textit{Returns:} Reference to this vector.\par}
\end{minipage}\par\vspace{1pt}
\par\noindent\begin{minipage}{\linewidth}
\begin{lstlisting}[style=apireference]
vector3 & operator/=(REAL scalar);
\end{lstlisting}
{\footnotesize Divide every component in place.\par}
{\footnotesize\textit{Returns:} Reference to this vector.\par}
\end{minipage}\par\vspace{1pt}
\par\noindent\begin{minipage}{\linewidth}
\begin{lstlisting}[style=apireference]
REAL dot(const vector3 &other) const noexcept;
\end{lstlisting}
{\footnotesize Compute the scalar dot product.\par}
{\footnotesize\textit{Returns:} Dot product.\par}
\end{minipage}\par\vspace{1pt}
\par\noindent\begin{minipage}{\linewidth}
\begin{lstlisting}[style=apireference]
vector3 cross(const vector3 &other) const noexcept;
\end{lstlisting}
{\footnotesize Compute the right-handed cross product.\par}
{\footnotesize\textit{Returns:} Vector perpendicular to both operands.\par}
\end{minipage}\par\vspace{1pt}
\par\noindent\begin{minipage}{\linewidth}
\begin{lstlisting}[style=apireference]
REAL length_squared() const noexcept;
\end{lstlisting}
{\footnotesize Compute squared Euclidean length.\par}
{\footnotesize\textit{Returns:} Squared length.\par}
\end{minipage}\par\vspace{1pt}
\par\noindent\begin{minipage}{\linewidth}
\begin{lstlisting}[style=apireference]
REAL length() const noexcept;
\end{lstlisting}
{\footnotesize Compute Euclidean length.\par}
{\footnotesize\textit{Returns:} Vector length.\par}
\end{minipage}\par\vspace{1pt}
\par\noindent\begin{minipage}{\linewidth}
\begin{lstlisting}[style=apireference]
bool is_near_zero(REAL tolerance=default_tolerance()) const noexcept;
\end{lstlisting}
{\footnotesize Test whether vector length is within a tolerance of zero.\par}
{\footnotesize\textit{Returns:} True when squared length is at most tolerance squared.\par}
\end{minipage}\par\vspace{1pt}
\par\noindent\begin{minipage}{\linewidth}
\begin{lstlisting}[style=apireference]
vector3 normalized(REAL tolerance=default_tolerance()) const;
\end{lstlisting}
{\footnotesize Return a unit vector in the same direction.\par}
{\footnotesize\textit{Returns:} Normalized vector.\par}
\end{minipage}\par\vspace{1pt}
\par\noindent\begin{minipage}{\linewidth}
\begin{lstlisting}[style=apireference]
void normalize();
\end{lstlisting}
{\footnotesize Normalize this vector in place using the default vector tolerance.\par}
\end{minipage}\par\vspace{1pt}
\par\noindent\begin{minipage}{\linewidth}
\begin{lstlisting}[style=apireference]
vector3 projected_onto(const vector3 &axis, REAL tolerance=default_tolerance()) const;
\end{lstlisting}
{\footnotesize Project this vector onto an axis.\par}
{\footnotesize\textit{Returns:} Component parallel to axis.\par}
\end{minipage}\par\vspace{1pt}
\par\noindent\begin{minipage}{\linewidth}
\begin{lstlisting}[style=apireference]
vector3 rejected_from(const vector3 &axis, REAL tolerance=default_tolerance()) const;
\end{lstlisting}
{\footnotesize Remove the component parallel to an axis.\par}
{\footnotesize\textit{Returns:} Component perpendicular to axis.\par}
\end{minipage}\par\vspace{1pt}
\par\noindent\begin{minipage}{\linewidth}
\begin{lstlisting}[style=apireference]
vector3 reflected(const vector3 &normal, REAL tolerance=default_tolerance()) const;
\end{lstlisting}
{\footnotesize Reflect this vector across a plane normal.\par}
{\footnotesize\textit{Returns:} Reflected vector.\par}
\end{minipage}\par\vspace{1pt}
\par\noindent\begin{minipage}{\linewidth}
\begin{lstlisting}[style=apireference]
REAL angle_to(const vector3 &other, REAL tolerance=default_tolerance()) const;
\end{lstlisting}
{\footnotesize Compute the unsigned angle to another vector.\par}
{\footnotesize\textit{Returns:} Angle in radians in the range [0, pi].\par}
\end{minipage}\par\vspace{1pt}
\par\noindent\begin{minipage}{\linewidth}
\begin{lstlisting}[style=apireference]
vector3 component_product(const vector3 &other) const noexcept;
\end{lstlisting}
{\footnotesize Multiply corresponding components.\par}
{\footnotesize\textit{Returns:} Component-wise product.\par}
\end{minipage}\par\vspace{1pt}
\par\noindent\begin{minipage}{\linewidth}
\begin{lstlisting}[style=apireference]
REAL max_component() const noexcept;
\end{lstlisting}
{\footnotesize Get the largest component.\par}
{\footnotesize\textit{Returns:} Largest component value.\par}
\end{minipage}\par\vspace{1pt}
\par\noindent\begin{minipage}{\linewidth}
\begin{lstlisting}[style=apireference]
REAL min_component() const noexcept;
\end{lstlisting}
{\footnotesize Get the smallest component.\par}
{\footnotesize\textit{Returns:} Smallest component value.\par}
\end{minipage}\par\vspace{1pt}
\par\noindent\begin{minipage}{\linewidth}
\begin{lstlisting}[style=apireference]
bool approximately_equal(const vector3 &other, REAL tolerance=default_tolerance()) const noexcept;
\end{lstlisting}
{\footnotesize Compare vectors using Euclidean distance.\par}
{\footnotesize\textit{Returns:} True when vectors differ by at most tolerance.\par}
\end{minipage}\par\vspace{1pt}
\par\noindent\begin{minipage}{\linewidth}
\begin{lstlisting}[style=apireference]
static std::optional< tagged_read_result< vector3 > > tag_read(std::istream &input);
\end{lstlisting}
{\footnotesize Read and allocate one version-1 tagged vector3 text record.\par}
{\footnotesize\textit{Returns:} Tag and non-null shared vector after success, or std::nullopt at EOF, another read failure, or after a malformed row.\par}
\end{minipage}\par\vspace{1pt}
\par\noindent\begin{minipage}{\linewidth}
\begin{lstlisting}[style=apireference]
static vector3 zero() noexcept;
\end{lstlisting}
{\footnotesize Construct the zero vector.\par}
{\footnotesize\textit{Returns:} (0,0,0).\par}
\end{minipage}\par\vspace{1pt}
\par\noindent\begin{minipage}{\linewidth}
\begin{lstlisting}[style=apireference]
static vector3 unit_x() noexcept;
\end{lstlisting}
{\footnotesize Construct the positive X unit vector.\par}
{\footnotesize\textit{Returns:} (1,0,0).\par}
\end{minipage}\par\vspace{1pt}
\par\noindent\begin{minipage}{\linewidth}
\begin{lstlisting}[style=apireference]
static vector3 unit_y() noexcept;
\end{lstlisting}
{\footnotesize Construct the positive Y unit vector.\par}
{\footnotesize\textit{Returns:} (0,1,0).\par}
\end{minipage}\par\vspace{1pt}
\par\noindent\begin{minipage}{\linewidth}
\begin{lstlisting}[style=apireference]
static vector3 unit_z() noexcept;
\end{lstlisting}
{\footnotesize Construct the positive Z unit vector.\par}
{\footnotesize\textit{Returns:} (0,0,1).\par}
\end{minipage}\par\vspace{1pt}
\par\noindent\begin{minipage}{\linewidth}
\begin{lstlisting}[style=apireference]
static REAL default_tolerance() noexcept;
\end{lstlisting}
{\footnotesize Get the default absolute vector tolerance.\par}
{\footnotesize\textit{Returns:} 64 times machine epsilon for REAL.\par}
\end{minipage}\par\vspace{1pt}
\section{Namespace functions, aliases, and result enums}
\Needspace{9\baselineskip}
\subsection{creators.hpp}
\par\noindent\begin{minipage}{\linewidth}
\begin{lstlisting}[style=apireference]
template <std::floating_point REAL>
std::shared_ptr< vector2< REAL > > make_vector2();
\end{lstlisting}
{\footnotesize Allocate a zero 2D vector with shared ownership.\par}
{\footnotesize\textit{Returns:} Shared smart pointer owning the new vector.\par}
\end{minipage}\par\vspace{1pt}
\par\noindent\begin{minipage}{\linewidth}
\begin{lstlisting}[style=apireference]
template <std::floating_point REAL>
std::shared_ptr< vector2< REAL > > make_vector2(REAL x, REAL y);
\end{lstlisting}
{\footnotesize Allocate a 2D vector with shared ownership.\par}
{\footnotesize\textit{Returns:} Shared smart pointer owning the new vector.\par}
\end{minipage}\par\vspace{1pt}
\par\noindent\begin{minipage}{\linewidth}
\begin{lstlisting}[style=apireference]
template <std::floating_point REAL>
std::shared_ptr< vector2< REAL > > make_vector2(const point2< REAL > &point_value);
\end{lstlisting}
{\footnotesize Allocate a 2D vector from a point with shared ownership.\par}
{\footnotesize\textit{Returns:} Shared smart pointer owning the new vector.\par}
\end{minipage}\par\vspace{1pt}
\par\noindent\begin{minipage}{\linewidth}
\begin{lstlisting}[style=apireference]
template <std::floating_point REAL>
std::shared_ptr< vector2< REAL > > make_vector2(const circle2< REAL > &circle_value);
\end{lstlisting}
{\footnotesize Allocate a 2D vector from a circle center with shared ownership.\par}
{\footnotesize\textit{Returns:} Shared smart pointer owning the new vector.\par}
\end{minipage}\par\vspace{1pt}
\par\noindent\begin{minipage}{\linewidth}
\begin{lstlisting}[style=apireference]
template <std::floating_point REAL>
std::shared_ptr< vector2< REAL > > make_vector2(const ray2< REAL > &ray_value);
\end{lstlisting}
{\footnotesize Allocate a 2D vector from a ray origin with shared ownership.\par}
{\footnotesize\textit{Returns:} Shared smart pointer owning the new vector.\par}
\end{minipage}\par\vspace{1pt}
\par\noindent\begin{minipage}{\linewidth}
\begin{lstlisting}[style=apireference]
template <std::floating_point REAL>
std::shared_ptr< point2< REAL > > make_point2();
\end{lstlisting}
{\footnotesize Allocate the 2D origin with shared ownership.\par}
{\footnotesize\textit{Returns:} Shared smart pointer owning the new point.\par}
\end{minipage}\par\vspace{1pt}
\par\noindent\begin{minipage}{\linewidth}
\begin{lstlisting}[style=apireference]
template <std::floating_point REAL>
std::shared_ptr< point2< REAL > > make_point2(REAL x, REAL y);
\end{lstlisting}
{\footnotesize Allocate a 2D point with shared ownership.\par}
{\footnotesize\textit{Returns:} Shared smart pointer owning the new point.\par}
\end{minipage}\par\vspace{1pt}
\par\noindent\begin{minipage}{\linewidth}
\begin{lstlisting}[style=apireference]
template <std::floating_point REAL>
std::shared_ptr< ray2< REAL > > make_ray2(const point2< REAL > &origin, const vector2< REAL > &direction, REAL tolerance=vector2< REAL >::default_tolerance());
\end{lstlisting}
{\footnotesize Allocate a 2D ray with shared ownership.\par}
{\footnotesize\textit{Returns:} Shared smart pointer owning the new ray.\par}
\end{minipage}\par\vspace{1pt}
\par\noindent\begin{minipage}{\linewidth}
\begin{lstlisting}[style=apireference]
template <std::floating_point REAL>
std::shared_ptr< ray2< REAL > > make_ray2_from_points(const point2< REAL > &origin, const point2< REAL > &through_point, REAL tolerance=vector2< REAL >::default_tolerance());
\end{lstlisting}
{\footnotesize Allocate a 2D ray defined by two points with shared ownership.\par}
{\footnotesize\textit{Returns:} Shared smart pointer owning the new ray. Its unnormalized direction is through\_point - origin.\par}
\end{minipage}\par\vspace{1pt}
\par\noindent\begin{minipage}{\linewidth}
\begin{lstlisting}[style=apireference]
template <std::floating_point REAL>
std::shared_ptr< circle2< REAL > > make_circle2(const point2< REAL > &center, REAL radius);
\end{lstlisting}
{\footnotesize Allocate a 2D circle with shared ownership.\par}
{\footnotesize\textit{Returns:} Shared smart pointer owning the new circle.\par}
\end{minipage}\par\vspace{1pt}
\par\noindent\begin{minipage}{\linewidth}
\begin{lstlisting}[style=apireference]
template <std::floating_point REAL>
std::shared_ptr< nurbs_spline2< REAL > > make_nurbs_spline2(const point2< REAL > &start, const point2< REAL > &end, std::size_t control_point_count, std::size_t degree, std::type_identity_t< REAL > tolerance);
\end{lstlisting}
{\footnotesize Allocate a standard open 2D NURBS with explicit tolerance.\par}
{\footnotesize\textit{Returns:} Shared smart pointer owning the new 2D spline.\par}
\end{minipage}\par\vspace{1pt}
\par\noindent\begin{minipage}{\linewidth}
\begin{lstlisting}[style=apireference]
template <std::floating_point REAL, std::same_as< bool > CLOSED = bool>
std::shared_ptr< nurbs_spline2< REAL > > make_nurbs_spline2(const point2< REAL > &start, const point2< REAL > &end, std::size_t control_point_count, std::size_t degree, CLOSED closed=false, std::type_identity_t< REAL > tolerance=REAL(64) *std::numeric_limits< REAL >::epsilon());
\end{lstlisting}
{\footnotesize Allocate a standard open or closed 2D NURBS between two points.\par}
{\footnotesize\textit{Returns:} Shared smart pointer owning the new 2D spline.\par}
\end{minipage}\par\vspace{1pt}
\par\noindent\begin{minipage}{\linewidth}
\begin{lstlisting}[style=apireference]
template <std::floating_point REAL>
std::shared_ptr< nurbs_spline2< REAL > > make_nurbs_spline2(std::vector< point2< REAL > > control_points, std::vector< REAL > weights, std::vector< REAL > knots, std::size_t degree, REAL tolerance=REAL(64) *std::numeric_limits< REAL >::epsilon());
\end{lstlisting}
{\footnotesize Allocate a complete 2D NURBS definition with shared ownership.\par}
{\footnotesize\textit{Returns:} Shared smart pointer owning the new 2D spline.\par}
\end{minipage}\par\vspace{1pt}
\par\noindent\begin{minipage}{\linewidth}
\begin{lstlisting}[style=apireference]
template <std::floating_point REAL>
std::shared_ptr< nurbs_spline2< REAL > > make_nurbs_spline2(std::vector< point2< REAL > > control_points, std::vector< REAL > weights, std::vector< REAL > knots, std::size_t degree, bool closed, REAL tolerance=REAL(64) *std::numeric_limits< REAL >::epsilon());
\end{lstlisting}
{\footnotesize Allocate an open or closed complete 2D NURBS definition.\par}
{\footnotesize\textit{Returns:} Shared smart pointer owning the new 2D spline.\par}
\end{minipage}\par\vspace{1pt}
\par\noindent\begin{minipage}{\linewidth}
\begin{lstlisting}[style=apireference]
template <std::floating_point REAL, typename control_point_array>
std::shared_ptr< nurbs_spline2< REAL > > make_nurbs_spline2(control_point_array &&control_points, const std::valarray< REAL > &weights, const std::valarray< REAL > &knots, std::size_t degree, REAL tolerance=REAL(64) *std::numeric_limits< REAL >::epsilon());
\end{lstlisting}
{\footnotesize Allocate a complete 2D NURBS definition from valarrays.\par}
{\footnotesize\textit{Returns:} Shared smart pointer owning the new 2D spline.\par}
\end{minipage}\par\vspace{1pt}
\par\noindent\begin{minipage}{\linewidth}
\begin{lstlisting}[style=apireference]
template <std::floating_point REAL, typename control_point_array>
std::shared_ptr< nurbs_spline2< REAL > > make_nurbs_spline2(control_point_array &&control_points, const std::valarray< REAL > &weights, const std::valarray< REAL > &knots, std::size_t degree, bool closed, REAL tolerance=REAL(64) *std::numeric_limits< REAL >::epsilon());
\end{lstlisting}
{\footnotesize Allocate an open or closed 2D NURBS definition from valarrays.\par}
{\footnotesize\textit{Returns:} Shared smart pointer owning the new 2D spline.\par}
\end{minipage}\par\vspace{1pt}
\par\noindent\begin{minipage}{\linewidth}
\begin{lstlisting}[style=apireference]
template <std::floating_point REAL>
std::shared_ptr< nurbs_spline2< REAL > > make_nurbs_spline2(const nurbs_arr_spline2< REAL > &spline);
\end{lstlisting}
{\footnotesize Allocate an ordinary 2D NURBS by copying an array-backed spline.\par}
{\footnotesize\textit{Returns:} Shared smart pointer owning an independent ordinary 2D spline.\par}
\end{minipage}\par\vspace{1pt}
\par\noindent\begin{minipage}{\linewidth}
\begin{lstlisting}[style=apireference]
template <std::floating_point REAL>
std::shared_ptr< nurbs_arr_spline2< REAL > > make_nurbs_arr_spline2(const point2< REAL > &start, const point2< REAL > &end, std::size_t control_point_count, std::size_t degree, std::type_identity_t< REAL > tolerance);
\end{lstlisting}
{\footnotesize Allocate a standard open array-backed 2D NURBS with explicit tolerance.\par}
{\footnotesize\textit{Returns:} Shared smart pointer owning the new array-backed spline.\par}
\end{minipage}\par\vspace{1pt}
\par\noindent\begin{minipage}{\linewidth}
\begin{lstlisting}[style=apireference]
template <std::floating_point REAL, std::same_as< bool > CLOSED = bool>
std::shared_ptr< nurbs_arr_spline2< REAL > > make_nurbs_arr_spline2(const point2< REAL > &start, const point2< REAL > &end, std::size_t control_point_count, std::size_t degree, CLOSED closed=false, std::type_identity_t< REAL > tolerance=REAL(64) *std::numeric_limits< REAL >::epsilon());
\end{lstlisting}
{\footnotesize Allocate a standard open or closed array-backed 2D NURBS.\par}
{\footnotesize\textit{Returns:} Shared smart pointer owning the new array-backed spline.\par}
\end{minipage}\par\vspace{1pt}
\par\noindent\begin{minipage}{\linewidth}
\begin{lstlisting}[style=apireference]
template <std::floating_point REAL>
std::shared_ptr< nurbs_arr_spline2< REAL > > make_nurbs_arr_spline2(std::valarray< REAL > parameters, std::size_t degree, std::type_identity_t< REAL > tolerance=REAL(64) *std::numeric_limits< REAL >::epsilon());
\end{lstlisting}
{\footnotesize Allocate an open array-backed 2D NURBS from flattened scalars.\par}
{\footnotesize\textit{Returns:} Shared smart pointer owning the new array-backed spline.\par}
\end{minipage}\par\vspace{1pt}
\par\noindent\begin{minipage}{\linewidth}
\begin{lstlisting}[style=apireference]
template <std::floating_point REAL, std::same_as< bool > CLOSED>
std::shared_ptr< nurbs_arr_spline2< REAL > > make_nurbs_arr_spline2(std::valarray< REAL > parameters, std::size_t degree, CLOSED closed, std::type_identity_t< REAL > tolerance=REAL(64) *std::numeric_limits< REAL >::epsilon());
\end{lstlisting}
{\footnotesize Allocate an open or closed array-backed 2D NURBS from scalars.\par}
{\footnotesize\textit{Returns:} Shared smart pointer owning the new array-backed spline.\par}
\end{minipage}\par\vspace{1pt}
\par\noindent\begin{minipage}{\linewidth}
\begin{lstlisting}[style=apireference]
template <std::floating_point REAL>
std::shared_ptr< nurbs_arr_spline2< REAL > > make_nurbs_arr_spline2(std::vector< point2< REAL > > control_points, std::vector< REAL > weights, std::vector< REAL > knots, std::size_t degree, std::type_identity_t< REAL > tolerance=REAL(64) *std::numeric_limits< REAL >::epsilon());
\end{lstlisting}
{\footnotesize Allocate an open array-backed 2D NURBS from vector collections.\par}
{\footnotesize\textit{Returns:} Shared smart pointer owning the new array-backed spline.\par}
\end{minipage}\par\vspace{1pt}
\par\noindent\begin{minipage}{\linewidth}
\begin{lstlisting}[style=apireference]
template <std::floating_point REAL, std::same_as< bool > CLOSED>
std::shared_ptr< nurbs_arr_spline2< REAL > > make_nurbs_arr_spline2(std::vector< point2< REAL > > control_points, std::vector< REAL > weights, std::vector< REAL > knots, std::size_t degree, CLOSED closed, std::type_identity_t< REAL > tolerance=REAL(64) *std::numeric_limits< REAL >::epsilon());
\end{lstlisting}
{\footnotesize Allocate an open or closed array-backed 2D NURBS from vectors.\par}
{\footnotesize\textit{Returns:} Shared smart pointer owning the new array-backed spline.\par}
\end{minipage}\par\vspace{1pt}
\par\noindent\begin{minipage}{\linewidth}
\begin{lstlisting}[style=apireference]
template <std::floating_point REAL>
std::shared_ptr< nurbs_arr_spline2< REAL > > make_nurbs_arr_spline2(const nurbs_spline2< REAL > &spline);
\end{lstlisting}
{\footnotesize Allocate an array-backed copy of an ordinary 2D NURBS spline.\par}
{\footnotesize\textit{Returns:} Shared smart pointer owning the independent array-backed spline.\par}
\end{minipage}\par\vspace{1pt}
\par\noindent\begin{minipage}{\linewidth}
\begin{lstlisting}[style=apireference]
template <std::floating_point REAL>
std::shared_ptr< vector3< REAL > > make_vector3();
\end{lstlisting}
{\footnotesize Allocate a zero 3D vector with shared ownership.\par}
{\footnotesize\textit{Returns:} Shared smart pointer owning the new vector.\par}
\end{minipage}\par\vspace{1pt}
\par\noindent\begin{minipage}{\linewidth}
\begin{lstlisting}[style=apireference]
template <std::floating_point REAL>
std::shared_ptr< vector3< REAL > > make_vector3(REAL x, REAL y, REAL z);
\end{lstlisting}
{\footnotesize Allocate a 3D vector with shared ownership.\par}
{\footnotesize\textit{Returns:} Shared smart pointer owning the new vector.\par}
\end{minipage}\par\vspace{1pt}
\par\noindent\begin{minipage}{\linewidth}
\begin{lstlisting}[style=apireference]
template <std::floating_point REAL>
std::shared_ptr< vector3< REAL > > make_vector3(const point3< REAL > &point_value);
\end{lstlisting}
{\footnotesize Allocate a 3D vector from a point with shared ownership.\par}
{\footnotesize\textit{Returns:} Shared smart pointer owning the new vector.\par}
\end{minipage}\par\vspace{1pt}
\par\noindent\begin{minipage}{\linewidth}
\begin{lstlisting}[style=apireference]
template <std::floating_point REAL>
std::shared_ptr< vector3< REAL > > make_vector3(const sphere3< REAL > &sphere_value);
\end{lstlisting}
{\footnotesize Allocate a 3D vector from a sphere center with shared ownership.\par}
{\footnotesize\textit{Returns:} Shared smart pointer owning the new vector.\par}
\end{minipage}\par\vspace{1pt}
\par\noindent\begin{minipage}{\linewidth}
\begin{lstlisting}[style=apireference]
template <std::floating_point REAL>
std::shared_ptr< vector3< REAL > > make_vector3(const plane3< REAL > &plane_value);
\end{lstlisting}
{\footnotesize Allocate a 3D vector from a plane origin with shared ownership.\par}
{\footnotesize\textit{Returns:} Shared smart pointer owning the new vector.\par}
\end{minipage}\par\vspace{1pt}
\par\noindent\begin{minipage}{\linewidth}
\begin{lstlisting}[style=apireference]
template <std::floating_point REAL>
std::shared_ptr< vector3< REAL > > make_vector3(const ray3< REAL > &ray_value);
\end{lstlisting}
{\footnotesize Allocate a 3D vector from a ray origin with shared ownership.\par}
{\footnotesize\textit{Returns:} Shared smart pointer owning the new vector.\par}
\end{minipage}\par\vspace{1pt}
\par\noindent\begin{minipage}{\linewidth}
\begin{lstlisting}[style=apireference]
template <std::floating_point REAL>
std::shared_ptr< point3< REAL > > make_point3();
\end{lstlisting}
{\footnotesize Allocate the 3D world origin with shared ownership.\par}
{\footnotesize\textit{Returns:} Shared smart pointer owning the new point.\par}
\end{minipage}\par\vspace{1pt}
\par\noindent\begin{minipage}{\linewidth}
\begin{lstlisting}[style=apireference]
template <std::floating_point REAL>
std::shared_ptr< point3< REAL > > make_point3(REAL x, REAL y, REAL z);
\end{lstlisting}
{\footnotesize Allocate a 3D point with shared ownership.\par}
{\footnotesize\textit{Returns:} Shared smart pointer owning the new point.\par}
\end{minipage}\par\vspace{1pt}
\par\noindent\begin{minipage}{\linewidth}
\begin{lstlisting}[style=apireference]
template <std::floating_point REAL>
std::shared_ptr< ray3< REAL > > make_ray3(const point3< REAL > &origin, const vector3< REAL > &direction, REAL tolerance=vector3< REAL >::default_tolerance());
\end{lstlisting}
{\footnotesize Allocate a 3D ray with shared ownership.\par}
{\footnotesize\textit{Returns:} Shared smart pointer owning the new ray.\par}
\end{minipage}\par\vspace{1pt}
\par\noindent\begin{minipage}{\linewidth}
\begin{lstlisting}[style=apireference]
template <std::floating_point REAL>
std::shared_ptr< ray3< REAL > > make_ray3_from_points(const point3< REAL > &origin, const point3< REAL > &through_point, REAL tolerance=vector3< REAL >::default_tolerance());
\end{lstlisting}
{\footnotesize Allocate a 3D ray defined by two world points with shared ownership.\par}
{\footnotesize\textit{Returns:} Shared smart pointer owning the new ray. Its unnormalized direction is through\_point - origin.\par}
\end{minipage}\par\vspace{1pt}
\par\noindent\begin{minipage}{\linewidth}
\begin{lstlisting}[style=apireference]
template <std::floating_point REAL>
std::shared_ptr< sphere3< REAL > > make_sphere3(const point3< REAL > &center, REAL radius);
\end{lstlisting}
{\footnotesize Allocate a 3D sphere with shared ownership.\par}
{\footnotesize\textit{Returns:} Shared smart pointer owning the new sphere.\par}
\end{minipage}\par\vspace{1pt}
\par\noindent\begin{minipage}{\linewidth}
\begin{lstlisting}[style=apireference]
template <std::floating_point REAL>
std::shared_ptr< plane3< REAL > > make_plane3(const vector3< REAL > &normal, REAL signed_distance, REAL tolerance=vector3< REAL >::default_tolerance());
\end{lstlisting}
{\footnotesize Allocate a Hessian normal-form plane with shared ownership.\par}
{\footnotesize\textit{Returns:} Shared smart pointer owning the new plane.\par}
\end{minipage}\par\vspace{1pt}
\par\noindent\begin{minipage}{\linewidth}
\begin{lstlisting}[style=apireference]
template <std::floating_point REAL>
std::shared_ptr< plane3< REAL > > make_plane3(const point3< REAL > &origin, const vector3< REAL > &normal, REAL tolerance=vector3< REAL >::default_tolerance());
\end{lstlisting}
{\footnotesize Allocate a point-and-normal plane with shared ownership.\par}
{\footnotesize\textit{Returns:} Shared smart pointer owning the new plane.\par}
\end{minipage}\par\vspace{1pt}
\par\noindent\begin{minipage}{\linewidth}
\begin{lstlisting}[style=apireference]
template <std::floating_point REAL>
std::shared_ptr< plane3< REAL > > make_plane3(const point3< REAL > &origin, const vector3< REAL > &normal, const vector3< REAL > &u_hint, REAL tolerance=vector3< REAL >::default_tolerance());
\end{lstlisting}
{\footnotesize Allocate a plane with a preferred positive-u direction.\par}
{\footnotesize\textit{Returns:} Shared smart pointer owning the new plane.\par}
\end{minipage}\par\vspace{1pt}
\par\noindent\begin{minipage}{\linewidth}
\begin{lstlisting}[style=apireference]
template <std::floating_point REAL>
std::shared_ptr< plane3< REAL > > make_plane3_from_u_direction(const point3< REAL > &origin, const vector3< REAL > &u_direction, const point3< REAL > &plane_point, REAL tolerance=vector3< REAL >::default_tolerance());
\end{lstlisting}
{\footnotesize Allocate a plane from a parameter origin, positive-u direction, and another plane point with shared ownership.\par}
{\footnotesize\textit{Returns:} Shared smart pointer owning a plane with a right-handed (u,v,normal) frame.\par}
\end{minipage}\par\vspace{1pt}
\par\noindent\begin{minipage}{\linewidth}
\begin{lstlisting}[style=apireference]
template <std::floating_point REAL>
std::shared_ptr< plane3< REAL > > make_plane3_from_points(const point3< REAL > &origin, const point3< REAL > &u_point, const point3< REAL > &plane_point, REAL tolerance=vector3< REAL >::default_tolerance());
\end{lstlisting}
{\footnotesize Allocate a plane from three non-collinear world points with shared ownership.\par}
{\footnotesize\textit{Returns:} Shared smart pointer owning a plane with a right-handed (u,v,normal) frame.\par}
\end{minipage}\par\vspace{1pt}
\par\noindent\begin{minipage}{\linewidth}
\begin{lstlisting}[style=apireference]
template <std::floating_point REAL>
std::shared_ptr< nurbs_spline3< REAL > > make_nurbs_spline3(const point3< REAL > &start, const point3< REAL > &end, std::size_t control_point_count, std::size_t degree, std::type_identity_t< REAL > tolerance);
\end{lstlisting}
{\footnotesize Allocate a standard open 3D NURBS with explicit tolerance.\par}
{\footnotesize\textit{Returns:} Shared smart pointer owning the new 3D spline.\par}
\end{minipage}\par\vspace{1pt}
\par\noindent\begin{minipage}{\linewidth}
\begin{lstlisting}[style=apireference]
template <std::floating_point REAL, std::same_as< bool > CLOSED = bool>
std::shared_ptr< nurbs_spline3< REAL > > make_nurbs_spline3(const point3< REAL > &start, const point3< REAL > &end, std::size_t control_point_count, std::size_t degree, CLOSED closed=false, std::type_identity_t< REAL > tolerance=REAL(64) *std::numeric_limits< REAL >::epsilon());
\end{lstlisting}
{\footnotesize Allocate a standard open or closed 3D NURBS between two points.\par}
{\footnotesize\textit{Returns:} Shared smart pointer owning the new 3D spline.\par}
\end{minipage}\par\vspace{1pt}
\par\noindent\begin{minipage}{\linewidth}
\begin{lstlisting}[style=apireference]
template <std::floating_point REAL>
std::shared_ptr< nurbs_spline3< REAL > > make_nurbs_spline3(std::vector< point3< REAL > > control_points, std::vector< REAL > weights, std::vector< REAL > knots, std::size_t degree, REAL tolerance=REAL(64) *std::numeric_limits< REAL >::epsilon());
\end{lstlisting}
{\footnotesize Allocate a complete 3D NURBS definition with shared ownership.\par}
{\footnotesize\textit{Returns:} Shared smart pointer owning the new 3D spline.\par}
\end{minipage}\par\vspace{1pt}
\par\noindent\begin{minipage}{\linewidth}
\begin{lstlisting}[style=apireference]
template <std::floating_point REAL>
std::shared_ptr< nurbs_spline3< REAL > > make_nurbs_spline3(std::vector< point3< REAL > > control_points, std::vector< REAL > weights, std::vector< REAL > knots, std::size_t degree, bool closed, REAL tolerance=REAL(64) *std::numeric_limits< REAL >::epsilon());
\end{lstlisting}
{\footnotesize Allocate an open or closed complete 3D NURBS definition.\par}
{\footnotesize\textit{Returns:} Shared smart pointer owning the new 3D spline.\par}
\end{minipage}\par\vspace{1pt}
\par\noindent\begin{minipage}{\linewidth}
\begin{lstlisting}[style=apireference]
template <std::floating_point REAL, typename control_point_array>
std::shared_ptr< nurbs_spline3< REAL > > make_nurbs_spline3(control_point_array &&control_points, const std::valarray< REAL > &weights, const std::valarray< REAL > &knots, std::size_t degree, REAL tolerance=REAL(64) *std::numeric_limits< REAL >::epsilon());
\end{lstlisting}
{\footnotesize Allocate a complete 3D NURBS definition from valarrays.\par}
{\footnotesize\textit{Returns:} Shared smart pointer owning the new 3D spline.\par}
\end{minipage}\par\vspace{1pt}
\par\noindent\begin{minipage}{\linewidth}
\begin{lstlisting}[style=apireference]
template <std::floating_point REAL, typename control_point_array>
std::shared_ptr< nurbs_spline3< REAL > > make_nurbs_spline3(control_point_array &&control_points, const std::valarray< REAL > &weights, const std::valarray< REAL > &knots, std::size_t degree, bool closed, REAL tolerance=REAL(64) *std::numeric_limits< REAL >::epsilon());
\end{lstlisting}
{\footnotesize Allocate an open or closed 3D NURBS definition from valarrays.\par}
{\footnotesize\textit{Returns:} Shared smart pointer owning the new 3D spline.\par}
\end{minipage}\par\vspace{1pt}
\par\noindent\begin{minipage}{\linewidth}
\begin{lstlisting}[style=apireference]
template <std::floating_point REAL>
std::shared_ptr< nurbs_spline3< REAL > > make_nurbs_spline3(const nurbs_arr_spline3< REAL > &spline);
\end{lstlisting}
{\footnotesize Allocate an ordinary 3D NURBS by copying an array-backed spline.\par}
{\footnotesize\textit{Returns:} Shared smart pointer owning an independent ordinary 3D spline.\par}
\end{minipage}\par\vspace{1pt}
\par\noindent\begin{minipage}{\linewidth}
\begin{lstlisting}[style=apireference]
template <std::floating_point REAL>
std::shared_ptr< nurbs_arr_spline3< REAL > > make_nurbs_arr_spline3(const point3< REAL > &start, const point3< REAL > &end, std::size_t control_point_count, std::size_t degree, std::type_identity_t< REAL > tolerance);
\end{lstlisting}
{\footnotesize Allocate a standard open array-backed 3D NURBS with explicit tolerance.\par}
{\footnotesize\textit{Returns:} Shared smart pointer owning the new array-backed spline.\par}
\end{minipage}\par\vspace{1pt}
\par\noindent\begin{minipage}{\linewidth}
\begin{lstlisting}[style=apireference]
template <std::floating_point REAL, std::same_as< bool > CLOSED = bool>
std::shared_ptr< nurbs_arr_spline3< REAL > > make_nurbs_arr_spline3(const point3< REAL > &start, const point3< REAL > &end, std::size_t control_point_count, std::size_t degree, CLOSED closed=false, std::type_identity_t< REAL > tolerance=REAL(64) *std::numeric_limits< REAL >::epsilon());
\end{lstlisting}
{\footnotesize Allocate a standard open or closed array-backed 3D NURBS.\par}
{\footnotesize\textit{Returns:} Shared smart pointer owning the new array-backed spline.\par}
\end{minipage}\par\vspace{1pt}
\par\noindent\begin{minipage}{\linewidth}
\begin{lstlisting}[style=apireference]
template <std::floating_point REAL>
std::shared_ptr< nurbs_arr_spline3< REAL > > make_nurbs_arr_spline3(std::valarray< REAL > parameters, std::size_t degree, std::type_identity_t< REAL > tolerance=REAL(64) *std::numeric_limits< REAL >::epsilon());
\end{lstlisting}
{\footnotesize Allocate an open array-backed 3D NURBS from flattened scalars.\par}
{\footnotesize\textit{Returns:} Shared smart pointer owning the new array-backed spline.\par}
\end{minipage}\par\vspace{1pt}
\par\noindent\begin{minipage}{\linewidth}
\begin{lstlisting}[style=apireference]
template <std::floating_point REAL, std::same_as< bool > CLOSED>
std::shared_ptr< nurbs_arr_spline3< REAL > > make_nurbs_arr_spline3(std::valarray< REAL > parameters, std::size_t degree, CLOSED closed, std::type_identity_t< REAL > tolerance=REAL(64) *std::numeric_limits< REAL >::epsilon());
\end{lstlisting}
{\footnotesize Allocate an open or closed array-backed 3D NURBS from scalars.\par}
{\footnotesize\textit{Returns:} Shared smart pointer owning the new array-backed spline.\par}
\end{minipage}\par\vspace{1pt}
\par\noindent\begin{minipage}{\linewidth}
\begin{lstlisting}[style=apireference]
template <std::floating_point REAL>
std::shared_ptr< nurbs_arr_spline3< REAL > > make_nurbs_arr_spline3(std::vector< point3< REAL > > control_points, std::vector< REAL > weights, std::vector< REAL > knots, std::size_t degree, std::type_identity_t< REAL > tolerance=REAL(64) *std::numeric_limits< REAL >::epsilon());
\end{lstlisting}
{\footnotesize Allocate an open array-backed 3D NURBS from vector collections.\par}
{\footnotesize\textit{Returns:} Shared smart pointer owning the new array-backed spline.\par}
\end{minipage}\par\vspace{1pt}
\par\noindent\begin{minipage}{\linewidth}
\begin{lstlisting}[style=apireference]
template <std::floating_point REAL, std::same_as< bool > CLOSED>
std::shared_ptr< nurbs_arr_spline3< REAL > > make_nurbs_arr_spline3(std::vector< point3< REAL > > control_points, std::vector< REAL > weights, std::vector< REAL > knots, std::size_t degree, CLOSED closed, std::type_identity_t< REAL > tolerance=REAL(64) *std::numeric_limits< REAL >::epsilon());
\end{lstlisting}
{\footnotesize Allocate an open or closed array-backed 3D NURBS from vectors.\par}
{\footnotesize\textit{Returns:} Shared smart pointer owning the new array-backed spline.\par}
\end{minipage}\par\vspace{1pt}
\par\noindent\begin{minipage}{\linewidth}
\begin{lstlisting}[style=apireference]
template <std::floating_point REAL>
std::shared_ptr< nurbs_arr_spline3< REAL > > make_nurbs_arr_spline3(const nurbs_spline3< REAL > &spline);
\end{lstlisting}
{\footnotesize Allocate an array-backed copy of an ordinary 3D NURBS spline.\par}
{\footnotesize\textit{Returns:} Shared smart pointer owning the independent array-backed spline.\par}
\end{minipage}\par\vspace{1pt}
\Needspace{9\baselineskip}
\subsection{distance.hpp}
\par\noindent\begin{minipage}{\linewidth}
\begin{lstlisting}[style=apireference]
template <std::floating_point REAL>
point_circle_distance2< REAL > distance_to_circle(const point2< REAL > &point, const circle2< REAL > &circle, REAL tolerance=REAL(1e-8));
\end{lstlisting}
{\footnotesize Compute radial point-to-circle distance information in 2D.\par}
{\footnotesize\textit{Returns:} 2D distance, closest circle point, and circle angle.\par}
\end{minipage}\par\vspace{1pt}
\par\noindent\begin{minipage}{\linewidth}
\begin{lstlisting}[style=apireference]
template <std::floating_point REAL>
point_spline_distance2< REAL > distance_to_spline(const point2< REAL > &point, const nurbs_spline2< REAL > &spline, const numerical_settings< REAL > &settings={});
\end{lstlisting}
{\footnotesize Numerically find the sampled global 2D point-to-spline minimum.\par}
{\footnotesize\textit{Returns:} 2D distance, closest spline point found, and spline parameter.\par}
\end{minipage}\par\vspace{1pt}
\par\noindent\begin{minipage}{\linewidth}
\begin{lstlisting}[style=apireference]
template <std::floating_point REAL>
point_surface_distance3< REAL > distance_to_plane(const point3< REAL > &point, const plane3< REAL > &plane, REAL tolerance=REAL(1e-8));
\end{lstlisting}
{\footnotesize Compute orthogonal point-to-plane distance information.\par}
{\footnotesize\textit{Returns:} Distance, closest point, and plane parameters.\par}
\end{minipage}\par\vspace{1pt}
\par\noindent\begin{minipage}{\linewidth}
\begin{lstlisting}[style=apireference]
template <std::floating_point REAL>
point_surface_distance3< REAL > distance_to_sphere(const point3< REAL > &point, const sphere3< REAL > &sphere, REAL tolerance=REAL(1e-8));
\end{lstlisting}
{\footnotesize Compute radial point-to-sphere distance information.\par}
{\footnotesize\textit{Returns:} Distance, closest point, and spherical parameters.\par}
\end{minipage}\par\vspace{1pt}
\par\noindent\begin{minipage}{\linewidth}
\begin{lstlisting}[style=apireference]
template <std::floating_point REAL>
point_spline_distance3< REAL > distance_to_spline(const point3< REAL > &point, const nurbs_spline3< REAL > &spline, const numerical_settings< REAL > &settings={});
\end{lstlisting}
{\footnotesize Numerically find the sampled global point-to-spline minimum.\par}
{\footnotesize\textit{Returns:} Distance, closest point found, and spline parameter.\par}
\end{minipage}\par\vspace{1pt}
\par\noindent\begin{minipage}{\linewidth}
\begin{lstlisting}[style=apireference]
template <std::floating_point REAL>
REAL distance(const point3< REAL > &point, const plane3< REAL > &plane, REAL tolerance=REAL(1e-8));
\end{lstlisting}
{\footnotesize Compute scalar point-to-plane distance.\par}
{\footnotesize\textit{Returns:} Tolerance-adjusted nonnegative distance.\par}
\end{minipage}\par\vspace{1pt}
\par\noindent\begin{minipage}{\linewidth}
\begin{lstlisting}[style=apireference]
template <std::floating_point REAL>
REAL distance(const point3< REAL > &point, const sphere3< REAL > &sphere, REAL tolerance=REAL(1e-8));
\end{lstlisting}
{\footnotesize Compute scalar point-to-sphere distance.\par}
{\footnotesize\textit{Returns:} Tolerance-adjusted nonnegative distance.\par}
\end{minipage}\par\vspace{1pt}
\par\noindent\begin{minipage}{\linewidth}
\begin{lstlisting}[style=apireference]
template <std::floating_point REAL>
REAL distance(const point3< REAL > &point, const nurbs_spline3< REAL > &spline, const numerical_settings< REAL > &settings={});
\end{lstlisting}
{\footnotesize Compute scalar numerical point-to-spline distance.\par}
{\footnotesize\textit{Returns:} Tolerance-adjusted nonnegative distance.\par}
\end{minipage}\par\vspace{1pt}
\par\noindent\begin{minipage}{\linewidth}
\begin{lstlisting}[style=apireference]
template <std::floating_point REAL>
REAL distance(const point2< REAL > &point, const circle2< REAL > &circle, REAL tolerance=REAL(1e-8));
\end{lstlisting}
{\footnotesize Compute scalar point-to-circle distance entirely in 2D.\par}
{\footnotesize\textit{Returns:} Tolerance-adjusted nonnegative 2D distance.\par}
\end{minipage}\par\vspace{1pt}
\par\noindent\begin{minipage}{\linewidth}
\begin{lstlisting}[style=apireference]
template <std::floating_point REAL>
REAL distance(const point2< REAL > &point, const nurbs_spline2< REAL > &spline, const numerical_settings< REAL > &settings={});
\end{lstlisting}
{\footnotesize Compute scalar numerical point-to-spline distance entirely in 2D.\par}
{\footnotesize\textit{Returns:} Tolerance-adjusted nonnegative 2D distance.\par}
\end{minipage}\par\vspace{1pt}
\Needspace{9\baselineskip}
\subsection{graphics.hpp}
\par\noindent\begin{minipage}{\linewidth}
\begin{lstlisting}[style=apireference]
enum class svg_projection3 { orthographic, perspective };
\end{lstlisting}
{\footnotesize Camera projections supported by the three-dimensional SVG renderer.\par}
\end{minipage}\par\vspace{1pt}
\par\noindent\begin{minipage}{\linewidth}
\begin{lstlisting}[style=apireference]
template <std::floating_point REAL, typename ENTITY>
std::string to_svg(const svg_view2< REAL > &view, const ENTITY &entity, const svg_graphics_options2< REAL > &options={});
\end{lstlisting}
{\footnotesize Render one native 2D entity to a standalone flat SVG document.\par}
{\footnotesize\textit{Returns:} Complete self-contained flat SVG XML string.\par}
\end{minipage}\par\vspace{1pt}
\par\noindent\begin{minipage}{\linewidth}
\begin{lstlisting}[style=apireference]
template <std::floating_point REAL, typename ENTITY>
std::string to_svg(const svg_view3< REAL > &view, const ENTITY &entity, const svg_graphics_options3< REAL > &options={});
\end{lstlisting}
{\footnotesize Render one supported entity to a standalone SVG document.\par}
{\footnotesize\textit{Returns:} Complete self-contained SVG XML string.\par}
\end{minipage}\par\vspace{1pt}
\par\noindent\begin{minipage}{\linewidth}
\begin{lstlisting}[style=apireference]
template <std::floating_point REAL, typename ENTITY>
std::string to_svg(const svg_view3< REAL > &view, const plane3< REAL > &plane, const ENTITY &entity, const svg_graphics_options3< REAL > &options={});
\end{lstlisting}
{\footnotesize Project one supported 2D entity through a plane and render it to SVG.\par}
{\footnotesize\textit{Returns:} Complete self-contained SVG XML string.\par}
\end{minipage}\par\vspace{1pt}
\Needspace{9\baselineskip}
\subsection{intersections.hpp}
\par\noindent\begin{minipage}{\linewidth}
\begin{lstlisting}[style=apireference]
enum class intersection_kind { none, discrete, coincident };
\end{lstlisting}
{\footnotesize Classification of a geometric intersection set.\par}
\end{minipage}\par\vspace{1pt}
\par\noindent\begin{minipage}{\linewidth}
\begin{lstlisting}[style=apireference]
template <std::floating_point REAL>
intersection_result2< REAL > intersect_ray_circle(const ray2< REAL > &ray, const circle2< REAL > &circle, const numerical_settings< REAL > &settings={});
\end{lstlisting}
{\footnotesize Numerically intersect a forward 2D ray with a 2D circle.\par}
{\footnotesize\textit{Returns:} Classified 2D contacts sorted by increasing s.\par}
\end{minipage}\par\vspace{1pt}
\par\noindent\begin{minipage}{\linewidth}
\begin{lstlisting}[style=apireference]
template <std::floating_point REAL>
intersection_result2< REAL > intersect_spline_circle(const nurbs_spline2< REAL > &spline, const circle2< REAL > &circle, const numerical_settings< REAL > &settings={});
\end{lstlisting}
{\footnotesize Numerically intersect a 2D NURBS spline with a 2D circle.\par}
{\footnotesize\textit{Returns:} Classified 2D contacts sorted by increasing s.\par}
\end{minipage}\par\vspace{1pt}
\par\noindent\begin{minipage}{\linewidth}
\begin{lstlisting}[style=apireference]
template <std::floating_point REAL>
intersection_result3< REAL > intersect_ray_sphere(const ray3< REAL > &ray, const sphere3< REAL > &sphere, const numerical_settings< REAL > &settings={});
\end{lstlisting}
{\footnotesize Numerically intersect a forward ray with a sphere.\par}
{\footnotesize\textit{Returns:} Classified contacts sorted by increasing s.\par}
\end{minipage}\par\vspace{1pt}
\par\noindent\begin{minipage}{\linewidth}
\begin{lstlisting}[style=apireference]
template <std::floating_point REAL>
intersection_result3< REAL > intersect_ray_plane(const ray3< REAL > &ray, const plane3< REAL > &plane, const numerical_settings< REAL > &settings={});
\end{lstlisting}
{\footnotesize Numerically intersect a forward ray with an infinite plane.\par}
{\footnotesize\textit{Returns:} None, one isolated contact, or a coincident classification.\par}
\end{minipage}\par\vspace{1pt}
\par\noindent\begin{minipage}{\linewidth}
\begin{lstlisting}[style=apireference]
template <std::floating_point REAL>
intersection_result3< REAL > intersect_spline_sphere(const nurbs_spline3< REAL > &spline, const sphere3< REAL > &sphere, const numerical_settings< REAL > &settings={});
\end{lstlisting}
{\footnotesize Numerically intersect a spline with a sphere over its active domain.\par}
{\footnotesize\textit{Returns:} Classified sampled-and-refined contacts sorted by s.\par}
\end{minipage}\par\vspace{1pt}
\par\noindent\begin{minipage}{\linewidth}
\begin{lstlisting}[style=apireference]
template <std::floating_point REAL>
intersection_result3< REAL > intersect_spline_plane(const nurbs_spline3< REAL > &spline, const plane3< REAL > &plane, const numerical_settings< REAL > &settings={});
\end{lstlisting}
{\footnotesize Numerically intersect a spline with an infinite plane.\par}
{\footnotesize\textit{Returns:} Classified sampled-and-refined contacts sorted by s.\par}
\end{minipage}\par\vspace{1pt}
\par\noindent\begin{minipage}{\linewidth}
\begin{lstlisting}[style=apireference]
template <std::floating_point REAL>
intersection_result3< REAL > intersect(const ray3< REAL > &ray, const sphere3< REAL > &sphere, const numerical_settings< REAL > &settings={});
\end{lstlisting}
{\footnotesize Convenience overload for ray-sphere intersection.\par}
{\footnotesize\textit{Returns:} Ray-sphere intersection result.\par}
\end{minipage}\par\vspace{1pt}
\par\noindent\begin{minipage}{\linewidth}
\begin{lstlisting}[style=apireference]
template <std::floating_point REAL>
intersection_result3< REAL > intersect(const ray3< REAL > &ray, const plane3< REAL > &plane, const numerical_settings< REAL > &settings={});
\end{lstlisting}
{\footnotesize Convenience overload for ray-plane intersection.\par}
{\footnotesize\textit{Returns:} Ray-plane intersection result.\par}
\end{minipage}\par\vspace{1pt}
\par\noindent\begin{minipage}{\linewidth}
\begin{lstlisting}[style=apireference]
template <std::floating_point REAL>
intersection_result3< REAL > intersect(const nurbs_spline3< REAL > &spline, const sphere3< REAL > &sphere, const numerical_settings< REAL > &settings={});
\end{lstlisting}
{\footnotesize Convenience overload for spline-sphere intersection.\par}
{\footnotesize\textit{Returns:} Spline-sphere intersection result.\par}
\end{minipage}\par\vspace{1pt}
\par\noindent\begin{minipage}{\linewidth}
\begin{lstlisting}[style=apireference]
template <std::floating_point REAL>
intersection_result3< REAL > intersect(const nurbs_spline3< REAL > &spline, const plane3< REAL > &plane, const numerical_settings< REAL > &settings={});
\end{lstlisting}
{\footnotesize Convenience overload for spline-plane intersection.\par}
{\footnotesize\textit{Returns:} Spline-plane intersection result.\par}
\end{minipage}\par\vspace{1pt}
\par\noindent\begin{minipage}{\linewidth}
\begin{lstlisting}[style=apireference]
template <std::floating_point REAL>
intersection_result2< REAL > intersect(const ray2< REAL > &ray, const circle2< REAL > &circle, const numerical_settings< REAL > &settings={});
\end{lstlisting}
{\footnotesize Convenience overload for 2D ray-circle intersection.\par}
{\footnotesize\textit{Returns:} Two-dimensional ray-circle intersection result.\par}
\end{minipage}\par\vspace{1pt}
\par\noindent\begin{minipage}{\linewidth}
\begin{lstlisting}[style=apireference]
template <std::floating_point REAL>
intersection_result2< REAL > intersect(const nurbs_spline2< REAL > &spline, const circle2< REAL > &circle, const numerical_settings< REAL > &settings={});
\end{lstlisting}
{\footnotesize Convenience overload for 2D spline-circle intersection.\par}
{\footnotesize\textit{Returns:} Two-dimensional spline-circle intersection result.\par}
\end{minipage}\par\vspace{1pt}
\Needspace{9\baselineskip}
\subsection{point2.hpp}
\par\noindent\begin{minipage}{\linewidth}
\begin{lstlisting}[style=apireference]
template <std::floating_point REAL>
std::ostream & operator<<(std::ostream &output, const point2< REAL > &value);
\end{lstlisting}
{\footnotesize Write a 2D point as whitespace-separated X and Y coordinates.\par}
{\footnotesize\textit{Returns:} Reference to output.\par}
\end{minipage}\par\vspace{1pt}
\par\noindent\begin{minipage}{\linewidth}
\begin{lstlisting}[style=apireference]
template <std::floating_point REAL>
std::istream & operator>>(std::istream &input, point2< REAL > &value);
\end{lstlisting}
{\footnotesize Read whitespace-separated X and Y coordinates into a 2D point.\par}
{\footnotesize\textit{Returns:} Reference to input.\par}
\end{minipage}\par\vspace{1pt}
\par\noindent\begin{minipage}{\linewidth}
\begin{lstlisting}[style=apireference]
template <std::floating_point REAL>
point2< REAL > operator+(point2< REAL > point, const vector2< REAL > &offset) noexcept;
\end{lstlisting}
{\footnotesize Translate a point by a vector.\par}
{\footnotesize\textit{Returns:} Translated point.\par}
\end{minipage}\par\vspace{1pt}
\par\noindent\begin{minipage}{\linewidth}
\begin{lstlisting}[style=apireference]
template <std::floating_point REAL>
point2< REAL > operator+(const vector2< REAL > &offset, point2< REAL > point) noexcept;
\end{lstlisting}
{\footnotesize Translate a point with vector-first syntax.\par}
{\footnotesize\textit{Returns:} Translated point.\par}
\end{minipage}\par\vspace{1pt}
\par\noindent\begin{minipage}{\linewidth}
\begin{lstlisting}[style=apireference]
template <std::floating_point REAL>
point2< REAL > operator-(point2< REAL > point, const vector2< REAL > &offset) noexcept;
\end{lstlisting}
{\footnotesize Translate a point by the negative of a vector.\par}
{\footnotesize\textit{Returns:} Translated point.\par}
\end{minipage}\par\vspace{1pt}
\par\noindent\begin{minipage}{\linewidth}
\begin{lstlisting}[style=apireference]
template <std::floating_point REAL>
vector2< REAL > operator-(const point2< REAL > &left, const point2< REAL > &right) noexcept;
\end{lstlisting}
{\footnotesize Form the displacement from one point to another.\par}
{\footnotesize\textit{Returns:} Vector from right to left.\par}
\end{minipage}\par\vspace{1pt}
\par\noindent\begin{minipage}{\linewidth}
\begin{lstlisting}[style=apireference]
template <std::floating_point REAL>
point2< REAL > lerp(const point2< REAL > &first, const point2< REAL > &second, REAL fraction) noexcept;
\end{lstlisting}
{\footnotesize Linearly interpolate points.\par}
{\footnotesize\textit{Returns:} Interpolated position.\par}
\end{minipage}\par\vspace{1pt}
\par\noindent\begin{minipage}{\linewidth}
\begin{lstlisting}[style=apireference]
template <std::floating_point REAL>
REAL distance_squared(const point2< REAL > &first, const point2< REAL > &second) noexcept;
\end{lstlisting}
{\footnotesize Compute squared Euclidean point distance.\par}
{\footnotesize\textit{Returns:} Squared distance.\par}
\end{minipage}\par\vspace{1pt}
\par\noindent\begin{minipage}{\linewidth}
\begin{lstlisting}[style=apireference]
template <std::floating_point REAL>
REAL distance(const point2< REAL > &first, const point2< REAL > &second) noexcept;
\end{lstlisting}
{\footnotesize Compute Euclidean point distance.\par}
{\footnotesize\textit{Returns:} Distance.\par}
\end{minipage}\par\vspace{1pt}
\Needspace{9\baselineskip}
\subsection{point3.hpp}
\par\noindent\begin{minipage}{\linewidth}
\begin{lstlisting}[style=apireference]
template <std::floating_point REAL>
std::ostream & operator<<(std::ostream &output, const point3< REAL > &value);
\end{lstlisting}
{\footnotesize Write a 3D point as whitespace-separated X, Y, and Z coordinates.\par}
{\footnotesize\textit{Returns:} Reference to output.\par}
\end{minipage}\par\vspace{1pt}
\par\noindent\begin{minipage}{\linewidth}
\begin{lstlisting}[style=apireference]
template <std::floating_point REAL>
std::istream & operator>>(std::istream &input, point3< REAL > &value);
\end{lstlisting}
{\footnotesize Read whitespace-separated X, Y, and Z coordinates into a 3D point.\par}
{\footnotesize\textit{Returns:} Reference to input.\par}
\end{minipage}\par\vspace{1pt}
\par\noindent\begin{minipage}{\linewidth}
\begin{lstlisting}[style=apireference]
template <std::floating_point REAL>
point3< REAL > operator+(point3< REAL > point, const vector3< REAL > &offset) noexcept;
\end{lstlisting}
{\footnotesize Translate a point by a vector.\par}
{\footnotesize\textit{Returns:} Translated point.\par}
\end{minipage}\par\vspace{1pt}
\par\noindent\begin{minipage}{\linewidth}
\begin{lstlisting}[style=apireference]
template <std::floating_point REAL>
point3< REAL > operator+(const vector3< REAL > &offset, point3< REAL > point) noexcept;
\end{lstlisting}
{\footnotesize Translate a point by a vector with vector-first syntax.\par}
{\footnotesize\textit{Returns:} Translated point.\par}
\end{minipage}\par\vspace{1pt}
\par\noindent\begin{minipage}{\linewidth}
\begin{lstlisting}[style=apireference]
template <std::floating_point REAL>
point3< REAL > operator-(point3< REAL > point, const vector3< REAL > &offset) noexcept;
\end{lstlisting}
{\footnotesize Translate a point by the negative of a vector.\par}
{\footnotesize\textit{Returns:} Translated point.\par}
\end{minipage}\par\vspace{1pt}
\par\noindent\begin{minipage}{\linewidth}
\begin{lstlisting}[style=apireference]
template <std::floating_point REAL>
vector3< REAL > operator-(const point3< REAL > &left, const point3< REAL > &right) noexcept;
\end{lstlisting}
{\footnotesize Form the displacement from one point to another.\par}
{\footnotesize\textit{Returns:} Vector from right to left.\par}
\end{minipage}\par\vspace{1pt}
\par\noindent\begin{minipage}{\linewidth}
\begin{lstlisting}[style=apireference]
template <std::floating_point REAL>
point3< REAL > rotate(const vector3< REAL > &axis, REAL angle, const point3< REAL > &point);
\end{lstlisting}
{\footnotesize Rotate a point about an axis through the Cartesian world origin.\par}
{\footnotesize\textit{Returns:} Rotated point in the same Cartesian world coordinate system.\par}
\end{minipage}\par\vspace{1pt}
\par\noindent\begin{minipage}{\linewidth}
\begin{lstlisting}[style=apireference]
template <std::floating_point REAL>
point3< REAL > lerp(const point3< REAL > &first, const point3< REAL > &second, REAL fraction) noexcept;
\end{lstlisting}
{\footnotesize Linearly interpolate two points.\par}
{\footnotesize\textit{Returns:} Interpolated position.\par}
\end{minipage}\par\vspace{1pt}
\par\noindent\begin{minipage}{\linewidth}
\begin{lstlisting}[style=apireference]
template <std::floating_point REAL>
REAL distance_squared(const point3< REAL > &first, const point3< REAL > &second) noexcept;
\end{lstlisting}
{\footnotesize Compute squared Euclidean distance between points.\par}
{\footnotesize\textit{Returns:} Squared distance.\par}
\end{minipage}\par\vspace{1pt}
\par\noindent\begin{minipage}{\linewidth}
\begin{lstlisting}[style=apireference]
template <std::floating_point REAL>
REAL distance(const point3< REAL > &first, const point3< REAL > &second) noexcept;
\end{lstlisting}
{\footnotesize Compute Euclidean distance between points.\par}
{\footnotesize\textit{Returns:} Distance.\par}
\end{minipage}\par\vspace{1pt}
\Needspace{9\baselineskip}
\subsection{projection.hpp}
\par\noindent\begin{minipage}{\linewidth}
\begin{lstlisting}[style=apireference]
template <std::floating_point REAL>
vector3< REAL > project(const plane3< REAL > &plane, const vector2< REAL > &vector) noexcept;
\end{lstlisting}
{\footnotesize Embed a 2D vector in a plane's orthonormal parameter frame.\par}
{\footnotesize\textit{Returns:} vector.x*u\_direction + vector.y*v\_direction in 3D.\par}
\end{minipage}\par\vspace{1pt}
\par\noindent\begin{minipage}{\linewidth}
\begin{lstlisting}[style=apireference]
template <std::floating_point REAL>
point3< REAL > project(const plane3< REAL > &plane, const point2< REAL > &point) noexcept;
\end{lstlisting}
{\footnotesize Embed a 2D point in a plane's parameter frame.\par}
{\footnotesize\textit{Returns:} 3D point plane.point\_at(point.x, point.y).\par}
\end{minipage}\par\vspace{1pt}
\par\noindent\begin{minipage}{\linewidth}
\begin{lstlisting}[style=apireference]
template <std::floating_point REAL>
ray3< REAL > project(const plane3< REAL > &plane, const ray2< REAL > &ray);
\end{lstlisting}
{\footnotesize Embed a 2D ray in a plane's parameter frame.\par}
{\footnotesize\textit{Returns:} 3D ray with projected origin and direction and the same native s.\par}
\end{minipage}\par\vspace{1pt}
\par\noindent\begin{minipage}{\linewidth}
\begin{lstlisting}[style=apireference]
template <std::floating_point REAL>
sphere3< REAL > project(const plane3< REAL > &plane, const circle2< REAL > &circle);
\end{lstlisting}
{\footnotesize Convert a 2D circle's center and radius to a 3D sphere.\par}
{\footnotesize\textit{Returns:} Sphere with projected center and unchanged radius.\par}
\end{minipage}\par\vspace{1pt}
\par\noindent\begin{minipage}{\linewidth}
\begin{lstlisting}[style=apireference]
template <std::floating_point REAL>
nurbs_spline3< REAL > project(const plane3< REAL > &plane, const nurbs_spline2< REAL > &spline);
\end{lstlisting}
{\footnotesize Embed a 2D NURBS spline by projecting its control points.\par}
{\footnotesize\textit{Returns:} 3D NURBS spline with projected control points.\par}
\end{minipage}\par\vspace{1pt}
\Needspace{9\baselineskip}
\subsection{serialization.hpp}
\par\noindent\begin{minipage}{\linewidth}
\begin{lstlisting}[style=apireference]
enum class text_format { csv, tsv, txt };
\end{lstlisting}
{\footnotesize Delimited text format used for Cartesian coordinate input and output.\par}
\end{minipage}\par\vspace{1pt}
\Needspace{9\baselineskip}
\subsection{tagged\_annotations.hpp}
\par\noindent\begin{minipage}{\linewidth}
\begin{lstlisting}[style=apireference]
enum class diagnostic_level { warning, information, error };
\end{lstlisting}
{\footnotesize Severity selecting an optional positioned diagnostic output token.\par}
\end{minipage}\par\vspace{1pt}
\Needspace{9\baselineskip}
\subsection{tagged\_serialization.hpp}
\par\noindent\begin{minipage}{\linewidth}
\begin{lstlisting}[style=apireference]
template <std::floating_point REAL>
using entity3_ptr = std::variant< std::shared_ptr< point3< REAL > >, std::shared_ptr< vector3< REAL > >, std::shared_ptr< sphere3< REAL > >, std::shared_ptr< nurbs_spline3< REAL > > >;
\end{lstlisting}
{\footnotesize Type-safe pointer alternatives supported by tagged 3D input.\par}
\end{minipage}\par\vspace{1pt}
\par\noindent\begin{minipage}{\linewidth}
\begin{lstlisting}[style=apireference]
template <std::floating_point REAL>
using tagged_record_ptr = std::variant< std::shared_ptr< point3< REAL > >, std::shared_ptr< vector3< REAL > >, std::shared_ptr< sphere3< REAL > >, std::shared_ptr< nurbs_spline3< REAL > >, std::shared_ptr< line_color< REAL > >, std::shared_ptr< line_font< REAL > >, std::shared_ptr< text_label< REAL > > >;
\end{lstlisting}
{\footnotesize Type-safe owners for geometry and optional presentation records.\par}
\end{minipage}\par\vspace{1pt}
\par\noindent\begin{minipage}{\linewidth}
\begin{lstlisting}[style=apireference]
template <std::floating_point REAL>
std::optional< tagged_entity3< REAL > > tag_read(std::istream &input);
\end{lstlisting}
{\footnotesize Read one version-1 tagged row and allocate its concrete 3D entity.\par}
{\footnotesize\textit{Returns:} Tag and entity variant selecting a non-null shared owner after success, or std::nullopt at EOF, another read failure, or after a malformed row. Unknown types, invalid definitions, and trailing fields set failbit.\par}
\end{minipage}\par\vspace{1pt}
\par\noindent\begin{minipage}{\linewidth}
\begin{lstlisting}[style=apireference]
template <std::floating_point REAL>
std::optional< tagged_record< REAL > > tag_read_record(std::istream &input);
\end{lstlisting}
{\footnotesize Read one geometry, linecolor, linefont, or text version-1 row.\par}
{\footnotesize\textit{Returns:} Preserved tag and variant holding a non-null owner, or nullopt on EOF/read failure. Invalid values, unknown types/versions, malformed quoting, or extra fields in a non-advisory row set failbit after consuming that row.\par}
\end{minipage}\par\vspace{1pt}
\Needspace{9\baselineskip}
\subsection{utility.hpp}
\par\noindent\begin{minipage}{\linewidth}
\begin{lstlisting}[style=apireference]
std::size_t nurbs_knot_count(std::size_t degree, std::size_t control_point_count);
\end{lstlisting}
{\footnotesize Calculate the knot-vector size required by a valid NURBS spline.\par}
{\footnotesize\textit{Returns:} Required number of values in the complete knot vector.\par}
\end{minipage}\par\vspace{1pt}
\par\noindent\begin{minipage}{\linewidth}
\begin{lstlisting}[style=apireference]
template <std::floating_point REAL>
REAL square(REAL value) noexcept;
\end{lstlisting}
{\footnotesize Square a scalar.\par}
{\footnotesize\textit{Returns:} value * value.\par}
\end{minipage}\par\vspace{1pt}
\par\noindent\begin{minipage}{\linewidth}
\begin{lstlisting}[style=apireference]
template <std::floating_point REAL>
bool approximately_equal(REAL first, REAL second, REAL tolerance) noexcept;
\end{lstlisting}
{\footnotesize Compare two scalar values with absolute tolerance.\par}
{\footnotesize\textit{Returns:} True when absolute difference is at most tolerance.\par}
\end{minipage}\par\vspace{1pt}
\par\noindent\begin{minipage}{\linewidth}
\begin{lstlisting}[style=apireference]
template <std::floating_point REAL>
std::vector< REAL > solve_linear_system(std::vector< std::vector< REAL > > matrix, std::vector< REAL > right_hand_side, REAL pivot_tolerance=REAL(64) *std::numeric_limits< REAL >::epsilon());
\end{lstlisting}
{\footnotesize Solve a square linear system using scaled partial pivoting.\par}
{\footnotesize\textit{Returns:} Solution vector x for A*x=b.\par}
\end{minipage}\par\vspace{1pt}
\par\noindent\begin{minipage}{\linewidth}
\begin{lstlisting}[style=apireference]
template <std::floating_point REAL, typename FUNCTION>
REAL bisect_root(FUNCTION &&function, REAL lower, REAL upper, const numerical_settings< REAL > &settings={});
\end{lstlisting}
{\footnotesize Refine a sign-changing root interval by bisection.\par}
{\footnotesize\textit{Returns:} Refined root parameter.\par}
\end{minipage}\par\vspace{1pt}
\par\noindent\begin{minipage}{\linewidth}
\begin{lstlisting}[style=apireference]
template <std::floating_point REAL, typename FUNCTION>
REAL golden_section_minimize(FUNCTION &&function, REAL lower, REAL upper, const numerical_settings< REAL > &settings={});
\end{lstlisting}
{\footnotesize Minimize a scalar function on a closed interval by golden-section search.\par}
{\footnotesize\textit{Returns:} Parameter at the refined sampled minimum.\par}
\end{minipage}\par\vspace{1pt}
\par\noindent\begin{minipage}{\linewidth}
\begin{lstlisting}[style=apireference]
template <std::floating_point REAL, typename FUNCTION>
std::vector< REAL > find_roots(FUNCTION &&function, REAL lower, REAL upper, const numerical_settings< REAL > &settings={});
\end{lstlisting}
{\footnotesize Discover and refine roots on a bounded sampled interval.\par}
{\footnotesize\textit{Returns:} Sorted, tolerance-deduplicated root parameters.\par}
\end{minipage}\par\vspace{1pt}
\Needspace{9\baselineskip}
\subsection{vector2.hpp}
\par\noindent\begin{minipage}{\linewidth}
\begin{lstlisting}[style=apireference]
template <std::floating_point REAL>
std::ostream & operator<<(std::ostream &output, const vector2< REAL > &value);
\end{lstlisting}
{\footnotesize Write a 2D vector as whitespace-separated X and Y components.\par}
{\footnotesize\textit{Returns:} Reference to output.\par}
\end{minipage}\par\vspace{1pt}
\par\noindent\begin{minipage}{\linewidth}
\begin{lstlisting}[style=apireference]
template <std::floating_point REAL>
std::istream & operator>>(std::istream &input, vector2< REAL > &value);
\end{lstlisting}
{\footnotesize Read whitespace-separated X and Y components into a 2D vector.\par}
{\footnotesize\textit{Returns:} Reference to input.\par}
\end{minipage}\par\vspace{1pt}
\par\noindent\begin{minipage}{\linewidth}
\begin{lstlisting}[style=apireference]
template <std::floating_point REAL>
vector2< REAL > operator+(vector2< REAL > left, const vector2< REAL > &right) noexcept;
\end{lstlisting}
{\footnotesize Add two vectors.\par}
{\footnotesize\textit{Returns:} Vector sum.\par}
\end{minipage}\par\vspace{1pt}
\par\noindent\begin{minipage}{\linewidth}
\begin{lstlisting}[style=apireference]
template <std::floating_point REAL>
vector2< REAL > operator-(vector2< REAL > left, const vector2< REAL > &right) noexcept;
\end{lstlisting}
{\footnotesize Subtract two vectors.\par}
{\footnotesize\textit{Returns:} Vector difference.\par}
\end{minipage}\par\vspace{1pt}
\par\noindent\begin{minipage}{\linewidth}
\begin{lstlisting}[style=apireference]
template <std::floating_point REAL>
vector2< REAL > operator*(vector2< REAL > value, REAL scalar) noexcept;
\end{lstlisting}
{\footnotesize Multiply a vector by a scalar.\par}
{\footnotesize\textit{Returns:} Scaled vector.\par}
\end{minipage}\par\vspace{1pt}
\par\noindent\begin{minipage}{\linewidth}
\begin{lstlisting}[style=apireference]
template <std::floating_point REAL>
vector2< REAL > operator*(REAL scalar, vector2< REAL > value) noexcept;
\end{lstlisting}
{\footnotesize Multiply a scalar by a vector.\par}
{\footnotesize\textit{Returns:} Scaled vector.\par}
\end{minipage}\par\vspace{1pt}
\par\noindent\begin{minipage}{\linewidth}
\begin{lstlisting}[style=apireference]
template <std::floating_point REAL>
vector2< REAL > operator/(vector2< REAL > value, REAL scalar);
\end{lstlisting}
{\footnotesize Divide a vector by a scalar.\par}
{\footnotesize\textit{Returns:} Divided vector.\par}
\end{minipage}\par\vspace{1pt}
\par\noindent\begin{minipage}{\linewidth}
\begin{lstlisting}[style=apireference]
template <std::floating_point REAL>
REAL dot(const vector2< REAL > &left, const vector2< REAL > &right) noexcept;
\end{lstlisting}
{\footnotesize Compute the dot product.\par}
{\footnotesize\textit{Returns:} Dot product.\par}
\end{minipage}\par\vspace{1pt}
\par\noindent\begin{minipage}{\linewidth}
\begin{lstlisting}[style=apireference]
template <std::floating_point REAL>
REAL cross(const vector2< REAL > &left, const vector2< REAL > &right) noexcept;
\end{lstlisting}
{\footnotesize Compute the signed scalar cross product.\par}
{\footnotesize\textit{Returns:} Signed scalar cross product.\par}
\end{minipage}\par\vspace{1pt}
\par\noindent\begin{minipage}{\linewidth}
\begin{lstlisting}[style=apireference]
template <std::floating_point REAL>
vector2< REAL > lerp(const vector2< REAL > &first, const vector2< REAL > &second, REAL fraction) noexcept;
\end{lstlisting}
{\footnotesize Linearly interpolate vectors.\par}
{\footnotesize\textit{Returns:} Interpolated vector.\par}
\end{minipage}\par\vspace{1pt}
\Needspace{9\baselineskip}
\subsection{vector3.hpp}
\par\noindent\begin{minipage}{\linewidth}
\begin{lstlisting}[style=apireference]
template <std::floating_point REAL>
std::ostream & operator<<(std::ostream &output, const vector3< REAL > &value);
\end{lstlisting}
{\footnotesize Write a 3D vector as whitespace-separated X, Y, and Z components.\par}
{\footnotesize\textit{Returns:} Reference to output.\par}
\end{minipage}\par\vspace{1pt}
\par\noindent\begin{minipage}{\linewidth}
\begin{lstlisting}[style=apireference]
template <std::floating_point REAL>
std::istream & operator>>(std::istream &input, vector3< REAL > &value);
\end{lstlisting}
{\footnotesize Read whitespace-separated X, Y, and Z components into a 3D vector.\par}
{\footnotesize\textit{Returns:} Reference to input.\par}
\end{minipage}\par\vspace{1pt}
\par\noindent\begin{minipage}{\linewidth}
\begin{lstlisting}[style=apireference]
template <std::floating_point REAL>
vector3< REAL > operator+(vector3< REAL > left, const vector3< REAL > &right) noexcept;
\end{lstlisting}
{\footnotesize Add two vectors.\par}
{\footnotesize\textit{Returns:} Vector sum.\par}
\end{minipage}\par\vspace{1pt}
\par\noindent\begin{minipage}{\linewidth}
\begin{lstlisting}[style=apireference]
template <std::floating_point REAL>
vector3< REAL > operator-(vector3< REAL > left, const vector3< REAL > &right) noexcept;
\end{lstlisting}
{\footnotesize Subtract two vectors.\par}
{\footnotesize\textit{Returns:} Vector difference.\par}
\end{minipage}\par\vspace{1pt}
\par\noindent\begin{minipage}{\linewidth}
\begin{lstlisting}[style=apireference]
template <std::floating_point REAL>
vector3< REAL > operator*(vector3< REAL > value, REAL scalar) noexcept;
\end{lstlisting}
{\footnotesize Multiply a vector by a scalar.\par}
{\footnotesize\textit{Returns:} Scaled vector.\par}
\end{minipage}\par\vspace{1pt}
\par\noindent\begin{minipage}{\linewidth}
\begin{lstlisting}[style=apireference]
template <std::floating_point REAL>
vector3< REAL > operator*(REAL scalar, vector3< REAL > value) noexcept;
\end{lstlisting}
{\footnotesize Multiply a scalar by a vector.\par}
{\footnotesize\textit{Returns:} Scaled vector.\par}
\end{minipage}\par\vspace{1pt}
\par\noindent\begin{minipage}{\linewidth}
\begin{lstlisting}[style=apireference]
template <std::floating_point REAL>
vector3< REAL > operator/(vector3< REAL > value, REAL scalar);
\end{lstlisting}
{\footnotesize Divide a vector by a scalar.\par}
{\footnotesize\textit{Returns:} Divided vector.\par}
\end{minipage}\par\vspace{1pt}
\par\noindent\begin{minipage}{\linewidth}
\begin{lstlisting}[style=apireference]
template <std::floating_point REAL>
REAL dot(const vector3< REAL > &left, const vector3< REAL > &right) noexcept;
\end{lstlisting}
{\footnotesize Compute the dot product of two vectors.\par}
{\footnotesize\textit{Returns:} Dot product.\par}
\end{minipage}\par\vspace{1pt}
\par\noindent\begin{minipage}{\linewidth}
\begin{lstlisting}[style=apireference]
template <std::floating_point REAL>
vector3< REAL > cross(const vector3< REAL > &left, const vector3< REAL > &right) noexcept;
\end{lstlisting}
{\footnotesize Compute the cross product of two vectors.\par}
{\footnotesize\textit{Returns:} Right-handed cross product.\par}
\end{minipage}\par\vspace{1pt}
\par\noindent\begin{minipage}{\linewidth}
\begin{lstlisting}[style=apireference]
template <std::floating_point REAL>
vector3< REAL > lerp(const vector3< REAL > &first, const vector3< REAL > &second, REAL fraction) noexcept;
\end{lstlisting}
{\footnotesize Linearly interpolate two vectors.\par}
{\footnotesize\textit{Returns:} Interpolated vector.\par}
\end{minipage}\par\vspace{1pt}
\endgroup
